% Modernised reading — fonds Alexandre Grothendieck, shelfmark 10, pages 1-197.
%
% This is an INTERPRETATION, not a transcription. It re-reads the pages in
% today's notation and vocabulary, states the hypotheses the manuscript leaves
% implicit, and drops the critical apparatus so that the mathematics reads
% straight through. Everything the brackets and underlines carried — a doubtful
% reading, a notation he used differently, a missing passage — moves into a
% footnote.
%
% It is held to being mathematically correct as it stands, not merely faithful
% to the page. Where the manuscript is loose or mistaken, the true statement is
% what appears here, and the footnote says what the page has. In any dispute of
% substance the transcription governs: whoever cites Grothendieck must cite the
% batch-NN.fr.tex files.
%
% Scope: ONE document for the whole shelfmark, its ten batches of transcription
% read as a single argument — which is what they are. The dictionary between a
% tensor category and a bialgebra is built twice, typed at pages 117-128 and
% again by hand from nothing at pages 156-179; the first build is what makes
% the second legible, and the second is what makes the three computed rings of
% pages 181-197 mean anything. The question the typed § 12 asks at page 129 —
% is a ⊗-functor compatible with everything necessarily faithful? — is the
% question the whole folder circles. Split across ten files none of that could
% be said once.
%
% NOTATION fixed once, for a hundred and ninety-seven pages. The folder reuses
% four letters for eight things, and reading one for another makes nonsense of
% both ends of it:
%
%   Gamma — THREE objects.
%     Gamma = Z(I), the free abelian group on the indecomposable classes
%       (pages 180-182), which this reading always writes Z(I) or R(C);
%     Gamma = X(T), the character group of a torus (pages 190, 197), written
%       X(T) here;
%     Gamma = the character group of the band D(Gamma) of a tannakian category
%       (pages 59-65), written X(D) here.
%   A — FOUR objects.
%     A = the coalgebra / commutative hyperalgebra dual to U (pages 120,
%       168-179), written A throughout, and always called the coalgebra of
%       coefficients;
%     A(X) = f_*(O_X) in the SGA 3 typescript (pages 42-57), which this reading
%       writes O(X) to keep it apart;
%     A(G) = the representation ring (pages 146-154, 181-197), written R(G)
%       here, and R_split(G) where the basis is the indecomposables;
%     A = End(M) in the polarisation notes (pages 22-30), written End(M).
%   U — one notion, three scripts (U, script U, k(G)): the topological algebra
%     of endomorphisms of the fibre functor. Always U here.
%   xi — THREE: the class of the tannakian gerbe in H^2 (pages 59-77), written
%     xi; the Tate character of a graded fibre functor (page 94), written tau;
%     the generator E_2 of the Green ring (pages 187-197), written xi.
%   F — the fibre functor everywhere, EXCEPT pages 66-83 where F_M is the
%     Frobenius of an F-crystal, and page 191 where F = k is the first factor
%     of a semi-direct product. Flagged where it occurs.
%
% Substantive departures from the pages, each footnoted where it occurs:
%   - page 185's Proposition says « indécomposables » where its own proof and
%     its Corollaire give « irréductibles »: k[X]/X^2 is indecomposable of
%     dimension 2;
%   - page 186's Prop. 2 says « groupes algébriques abéliens » where the
%     argument, which passes through a ONE-dimensional abelian Lie algebra,
%     gives only the one-dimensional unipotent group;
%   - page 188's E_5 = xi^4 - 3xi^2 + 5 must be + 1, as its own left-hand
%     column requires;
%   - page 188's sum runs to n with c_{nn} = 1 where its own table, and page
%     195, require n+1 and c_{n,n+1} = 1;
%   - page 101 calls the Tate motive « de rang 2 »; it is of rank 1 and weight 2;
%   - page 108's criterion for the Galois group to be profinite is stated for a
%     quotient pro-group « formé des G_i^0 »; the quotient is BY the identity
%     components;
%   - page 4's compatibility of Weil forms is an equivalence only on the forms
%     of a FIXED object, and the margin of that page says why not in general;
%   - pages 89 and 92's dictionary for filtered fibre functors is conditional
%     on Q being a torsor, which the folder does not prove; page 104 proves
%     instead that H' is pro-unipotent, which settles it for a countably
%     generated category;
%   - page 98's Proposition is proved in characteristic 0 only, and he says so;
%   - page 125's hexagon is flagged by him as unverified and is given here in
%     its standard form;
%   - pages 26 and 146 each carry a formula written twice with the same factor
%     where the two cases differ.
%
% What the folder announces and never establishes is collected in the spine
% section and recorded as absence, not stepped over: the treatise itself (the
% nine chapters of page 114); the non-existence of a fibre functor for the
% Koszul category (page 121, his own « ??? »); the element tau built from sigma
% (page 128); the four statements « A prouver » of page 132; the name left
% blank at page 134; the comparison Z[X(T)] → the Chern ring of R(T) (page
% 151); the bijectivity of CK(G) → Z[X(T)]^W without torsion (page 154); the
% local section of Q (page 89); the characteristic-p half of the Proposition of
% page 98; the essential image of the Betti-Hodge functor (page 107); which
% ⊗-categories come from an affine group scheme (page 112, his own « ??? »);
% an intrinsic description of the irreducibles on T/W (page 182); and the two
% exercises of pages 79-80 « que je n'ai pas fait ».
%
% Pages 32-57 and 85-86 are not on the subject of the folder: they are a
% typescript of SGA 3 (n° 10 and 11) and a leaf on classifying topoi with a
% census of the simple types on its back. They are read here briefly, because
% two of their results — every affine group scheme over a field is a projective
% limit of algebraic ones, and a group scheme of finite type over a field is
% affine iff it acts faithfully on a vector space — are exactly what the
% tannakian work leans on. Nothing on the leaves says that is why they are
% filed here.
%
% Pass: Opus 5 (claude-opus-5), 2026-09-12 — first pass, unchecked against the
% pages by a human.
\documentclass[11pt,a4paper]{article}
% Le préambule commun à toutes les transcriptions du fonds Grothendieck.
%
% Il tient en une idée : **l'incertitude se compose, elle ne se résout pas.**
% Une transcription qui lisse un mot douteux a détruit la seule chose pour
% laquelle on la faisait — un lecteur doit pouvoir savoir, à chaque endroit,
% ce qui était sur le feuillet et ce qui a été ajouté. Les six macros qui
% suivent sont donc l'appareil critique, et non de la décoration.
%
% Les mêmes commandes sont reconnues par `scripts/render.mjs`, qui produit la
% vue de lecture du site. Une macro ajoutée ici doit l'être aussi là-bas,
% faute de quoi le rendu échouera bruyamment — ce qui est voulu : un rendu
% silencieusement faux est pire qu'un rendu absent.

\ProvidesPackage{grothendieck}[2026/08/08 Transcriptions du fonds Grothendieck]

\usepackage{amsmath,amssymb,amsthm}
\usepackage{mathrsfs}
\usepackage{stmaryrd}
% Les diagrammes commutatifs sont partout dans ce fonds : c'est la langue
% même de la géométrie algébrique de Grothendieck, et une transcription qui
% les rendrait en prose ne transcrirait rien.
\usepackage{tikz-cd}
% Français par défaut — c'est la langue des feuillets — mais l'anglais est
% chargé aussi : la traduction et le résumé sont des documents anglais, et la
% typographie française (espaces avant les deux-points) y serait une faute.
% Ces documents basculent par \selectlanguage{english} en tête de corps.
\usepackage[english,french]{babel}
\usepackage{fontspec}
\usepackage{geometry}
\geometry{a4paper,margin=25mm,marginparwidth=28mm}
\usepackage{marginnote}
\usepackage{xcolor}
% ulem, not soul. `soul`'s \st and \ul are built on a character-by-character
% scan that XeLaTeX's font machinery breaks: the document fails with an
% undefined control sequence rather than degrading. ulem does the same two jobs
% and is Unicode-engine safe. `normalem` stops it hijacking \emph, which the
% transcriptions use for Grothendieck's own emphasis.
\usepackage[normalem]{ulem}
\usepackage{hyperref}
\hypersetup{colorlinks=true,linkcolor=black,urlcolor=black,pdfborder={0 0 0}}

% --- Métadonnées du lot -----------------------------------------------------
% Reprises telles quelles par le script de rendu, qui les affiche en tête de
% la vue de lecture. Elles fixent ce que le fichier prétend couvrir : sans
% elles, une transcription flotte sans point d'attache dans un fonds de seize
% mille feuillets.

\newcommand{\folder}[1]{\gdef\grothFolder{#1}}
\newcommand{\batch}[1]{\gdef\grothBatch{#1}}
\newcommand{\leaves}[2]{\gdef\grothFirst{#1}\gdef\grothLast{#2}}
\newcommand{\foldertitle}[1]{\gdef\grothFoldertitle{#1}}
\gdef\grothFolder{?}\gdef\grothBatch{?}\gdef\grothFirst{?}\gdef\grothLast{?}\gdef\grothFoldertitle{}

% --- Le résumé --------------------------------------------------------------
%
% Il ouvre la lecture modernisée, sous le titre « Résumé », et il est composé
% en petit. Ce n'est pas un ornement : c'est le seul passage du document qui
% ne soit pas des mathématiques, et un lecteur doit pouvoir le reconnaître
% avant de l'avoir lu — puis le sauter s'il connaît déjà le sujet, ou s'y
% arrêter s'il ne le connaît pas. Le filet qui le suit marque la reprise du
% corps.

\newenvironment{resume}
  {\small\setlength{\parskip}{0.45em}}
  {\par\medskip\hrule\medskip}

% --- Mots-clés ---------------------------------------------------------------
%
% En anglais, en clôture du résumé de la lecture modernisée : le vocabulaire
% moderne sous lequel chercher ce que les feuillets construisent. Le manifeste
% du site (`npm run manifest`) les extrait pour étiqueter la cote entière —
% cette ligne est la source unique des tags, il n'y a pas de fichier à côté.
\newcommand{\keywords}[1]{\par\smallskip\noindent
  {\footnotesize\textsc{Keywords} --- \textit{#1}}\par}

% --- L'appareil critique ----------------------------------------------------

% Le numéro de feuillet, en marge. C'est l'ancre qui permet de régler un
% désaccord sur une formule en regardant *un* feuillet plutôt qu'une chemise
% de six cents. Le site s'en sert aussi pour tourner les pages du fac-similé
% quand on fait défiler la transcription.
\newcommand{\leaf}[1]{%
  \marginnote{\footnotesize\color{gray!70}\textsf{#1}}%
  \label{leaf:#1}\ignorespaces}

% Illisible : on ne devine pas. Un mot inventé qui se lit comme les autres est
% le pire résultat possible.
\newcommand{\ill}{\textcolor{red!60!black}{[\,\ldots\,]}}

% Lecture douteuse : le mot est donné, et signalé comme incertain.
\newcommand{\uncertain}[1]{\uline{#1}}

% Ajout de l'éditeur — une lettre manquante, un mot sous-entendu.
\newcommand{\add}[1]{\textcolor{blue!50!black}{[#1]}}

% Biffé par Grothendieck lui-même. Ce qu'il a rayé fait partie du document :
% une rature dit souvent où la pensée a changé de direction.
\newcommand{\struck}[1]{\sout{#1}}

% Note du transcripteur, sur l'état matériel du feuillet.
\newcommand{\note}[1]{\par\smallskip{\small\color{gray!60!black}\textit{[#1]}}\par\smallskip}

% Note marginale de Grothendieck, distincte du corps du texte.
\newcommand{\marginal}[1]{\par\smallskip\noindent\hspace{1em}%
  {\small\color{gray!60!black}#1}\par\smallskip}

% --- Titre ------------------------------------------------------------------

\AtBeginDocument{%
  \begin{center}
    {\large\bfseries Fonds Alexandre Grothendieck --- cote n\textsuperscript{o}~\grothFolder}\\[2pt]
    {\itshape\grothFoldertitle}\\[4pt]
    {\small feuillets \grothFirst--\grothLast\ (lot \grothBatch)}\\[2pt]
    {\footnotesize\color{gray!60!black}Transcription établie d'après le fac-similé numérisé
     par l'Université de Montpellier.\\
     Les lectures incertaines sont soulignées, les illisibles notées [\,\ldots\,],
     les ajouts entre crochets.}
  \end{center}
  \vspace{1em}}

\endinput


\folder{10}
\pages{1}{197}
\dating{[à partir de 1958]}
\watermark{Édition de démonstration\\interprétation personnelle de l'œuvre}
\foldertitle{Catégories tensorielles : notes manuscrites (s.d.), tapuscrits (s.d.) — lecture modernisée du dossier entier}

\begin{document}

\section*{Résumé}

\begin{resume}
Cent quatre-vingt-dix-sept feuillets, tapés et manuscrits, tournent autour
d'une seule question : \emph{un groupe est-il connu quand on connaît toutes ses
représentations} ?

La question a une réponse facile quand le groupe est commutatif. Prenons le
cercle, le groupe des nombres complexes de module $1$. Ses représentations de
dimension $1$ sont les applications $z \mapsto z^{n}$, une pour chaque entier
$n$, et la liste de ces applications --- les \emph{caractères} --- est
exactement le groupe des entiers. Le cercle est reconstitué à partir de cette
liste, et c'est tout le contenu de la transformation de Fourier : la liste des
caractères est un groupe, et le groupe de départ est la liste des caractères de
cette liste.

Pour un groupe non commutatif, la liste ne suffit pas. Les représentations de
dimension quelconque ne se multiplient pas entre elles comme des nombres :
quand on en prend deux et qu'on les combine --- ce qu'on appelle leur produit
tensoriel, le procédé qui, de deux façons d'agir sur deux espaces, fait une
façon d'agir sur l'espace des tableaux à double entrée --- le résultat n'est
pas une représentation de la liste, mais une somme de plusieurs d'entre elles,
et le détail de cette décomposition est une information supplémentaire. L'idée
qui commande tout le dossier est que \emph{cette information supplémentaire est
exactement ce qui manque}. Une fois qu'on se donne non seulement la collection
des représentations, mais la loi qui dit comment elles se combinent, le groupe
revient.

Une collection d'objets munie d'une telle loi s'appelle une
\emph{$\otimes$-catégorie}. Le dossier est l'atelier d'un traité que
Grothendieck n'a pas écrit --- son plan en neuf chapitres est sur un feuillet,
page 114 --- et dont le sujet est le catalogue exact de ce qu'il faut ajouter à
une $\otimes$-catégorie pour qu'elle soit la collection des représentations
d'un groupe. Il y a d'abord une \emph{manière d'oublier}, qui à chaque objet
associe l'espace vectoriel nu sur lequel le groupe devait agir ; c'est ce qu'on
appelle un \emph{foncteur fibre}, et la première découverte est que se donner
un tel oubli revient à se donner une algèbre, et se donner la loi $\otimes$
revient à munir cette algèbre d'une opération de plus. Il y a ensuite les
\emph{contraintes} : que le produit soit associatif, qu'il soit commutatif,
qu'il ait une unité, que chaque objet ait un dual. Le travail de ces pages est
de montrer qu'elles sont \emph{indépendantes} les unes des autres, et que
chacune, traduite du côté de l'algèbre, devient une propriété nommable ---
associativité, commutativité, unité, antipode.

C'est là que l'affaire devient intéressante, parce qu'on peut en enlever une.
La cohomologie d'un espace topologique se multiplie avec un signe : échanger
deux classes de degrés impairs change le signe du produit. C'est la règle de
Koszul --- « Koko », dans son abréviation --- et Grothendieck constate qu'une
$\otimes$-catégorie où la commutativité est tordue de cette façon n'est la
collection des représentations d'aucun groupe, sur aucune extension du corps de
base. Ce n'est pas un accident de calcul : c'est le signe qu'il manque une
notion. Il l'appelle \emph{foncteur fibre tordu}, et ce qu'on appelle
aujourd'hui la solution --- une catégorie de représentations d'un groupe muni
d'un élément central de carré $1$, qui décide de la parité --- est déjà là,
page 123, sous la forme du critère : il suffit que l'homomorphisme $\mu_{2} \to
G$ soit central.

Ailleurs il enlève la commutativité tout entière, et il compte ce qu'on perd.
L'exemple des pages 129 à 131 est le plus net du dossier : pour un groupe fini
$G$, il fabrique une $\otimes$-catégorie dont les manières d'oublier sont
exactement les extensions de $G$ par les scalaires inversibles, donc sont
classées par un groupe de cohomologie $H^{2}(G, k^{*})$ qui n'est pas nul même
pour $G$ commutatif et $k$ algébriquement clos, et dont un élément non nul
\emph{ne meurt dans aucune extension du corps}. Conclusion, qu'il souligne :
« on ne peut négliger la contrainte de commutativité comme on aurait aimé
pouvoir le faire ».

Le reste du dossier est fait de ce qu'on peut ajouter à une
$\otimes$-catégorie, et de trois cas calculés jusqu'au bout. On peut ajouter
une \emph{graduation} du foncteur fibre --- c'est ce que fait la cohomologie,
qui vient en degrés --- et une graduation revient à un homomorphisme d'un
groupe multiplicatif dans le groupe cherché. On peut ajouter une
\emph{filtration}, qui est moins qu'une graduation : une suite emboîtée de
sous-espaces, au lieu d'une décomposition en somme directe. La différence entre
les deux est l'objet des pages 88 à 98 et 101 à 105, et elle est mesurée par un
espace de manières de passer de l'une à l'autre --- un \emph{torseur} --- dont
l'existence d'un point est précisément le fait que la filtration se scinde. Le
résultat est que sur un corps parfait elle se scinde toujours, mais
\emph{jamais canoniquement} : c'est exactement la situation de la théorie de
Hodge, où la cohomologie de De Rham est filtrée, la cohomologie de Hodge
graduée, et le passage de l'une à l'autre n'est pas donné avec la variété. On
peut enfin ajouter une \emph{polarisation}, c'est-à-dire distinguer parmi les
formes bilinéaires celles qui sont « positives » ; c'est le sujet des trente
premières pages, reprises trois fois, et l'invariant qui en sort --- le nombre
de cônes riemanniens --- est une question qu'il pose et laisse ouverte.

Les trois cas calculés forment les quinze dernières pages, et ils sont
élémentaires au sens où l'on peut les vérifier à la main. Pour le groupe
additif d'un corps de caractéristique nulle, l'anneau des représentations est
un anneau de polynômes en une variable de degré $2$ ; pour le groupe
multiplicatif, c'est l'anneau du groupe des entiers ; pour leur produit
semi-direct, un anneau de polynômes en une variable au-dessus du précédent. Et
les pages 66 à 83, les plus difficiles du dossier, font le même travail dans un
cas beaucoup plus profond : les cristaux en caractéristique $p$, où la structure
supplémentaire est un opérateur de Frobenius, et où le groupe qui sort est un
tore dont le groupe des caractères est le groupe des nombres rationnels --- ce
qu'on appelle aujourd'hui le tore des pentes, et dont les représentations sont
la classification de Dieudonné et Manin.

Ce que le lecteur trouvera enfin, et qui ne demande aucune connaissance
préalable, c'est Grothendieck en train de travailler : le soulagement quand un
obstacle tombe (« On respire ! »), l'aveu (« sauf erreur de signe, car je n'ai
nulle part fait attention aux signes »), l'exercice laissé (« Exercice amusant,
que je n'ai pas fait »), le plaisir nommé (« Ceci est un exercice plaisant et
délectable »), et surtout sa manière de juger un énoncé : le théorème de Manin,
qu'il récrit page 83, est donné « sous une forme plus jolie et débarassé
d'hypothèses restrictives parasites ». C'est la phrase qui dit le mieux ce que
ce dossier cherche : non pas de nouveaux théorèmes, mais la forme sous laquelle
les théorèmes connus cessent d'avoir besoin d'hypothèses.

\keywords{tannakian category, fibre functor, Tannaka reconstruction, affine
group scheme, pro-algebraic group, bialgebra, coalgebra of coefficients,
comodule, grouplike element, primitive element, rigidity, commutativity
constraint, Koszul sign rule, super-tannakian category, parity element, twisted
fibre functor, universal R-matrix, triangular Hopf algebra, Drinfeld twist,
graded vector spaces, pointed fusion category, band, gerbe, Brauer group,
diagonalisable group, numerical character, character table, split Grothendieck
ring, Green ring, based ring, lambda-ring, ring of symmetric functions,
plethysm, induction product, Chern class, Chevalley restriction, Weyl character,
representation ring of a torus, F-crystal, isocrystal, unit-root crystal,
Newton polygon, slope decomposition, Dieudonne-Manin classification, slope
torus, p-divisible group, Barsotti-Tate group, crystalline cohomology, Frobenius
gerbe, local class field theory, tensor filtration, parabolic subgroup,
unipotent radical, Levi decomposition, Rosenlicht vanishing, splitting of
filtrations, Hodge structure, polarisation, Deligne torus, Mumford-Tate group,
Hodge conjecture, motivic Galois group, Betti realisation, De Rham realisation,
weight homomorphism, Weil form, Riemann algebra, self-dual cone, Rosati
involution, motives with coefficients, Larson-Sweedler theorem, Hopf module,
Steenrod algebra, Chevalley's theorem, Lazard's theorem, faithfully flat
descent, fpqc quotient, Jordan block, Clebsch-Gordan rule, additive group, wild
representation type}
\end{resume}

\section*{Le fil du dossier, et les feuillets qui s'y sont glissés}

\subsection*{Les stations}

Le dossier n'est pas un texte suivi. C'est un atelier : quatre tapuscrits
corrigés, six faisceaux de sa main, un plan de traité, une table des matières,
et deux intrus. Voici les stations, dans l'ordre des feuillets, avec ce que
chacune apporte au fil.

\begin{itemize}
\item \textbf{Pages 3 à 16}, de sa main, sous sa chemise « Polarisations sur
les $\otimes$-catégories graduées ». L'axiomatique d'une polarisation :
formes de Weil, compatibilité, prépolarisation, prépolarisation maximale,
admissibilité, polarisation. Se termine sur le cas $\mathcal{M} =
\operatorname{Rep}(G)$ et l'élément $C$ de carré $\eta$ qui la définit.

\item \textbf{Pages 18 à 30}, d'une encre et d'un papier plus anciens, sous son
titre « Notes anciennes ». La même matière prise par l'autre bout : non plus
l'axiomatique, mais l'algèbre $\operatorname{End}(M)$ munie de l'involution
adjointe, les cônes de formes positives et leur polarité. Ce sont deux
\emph{moutures} d'un même sujet et non une répétition ; leur écart est le
contenu.

\item \textbf{Pages 32 à 57}, tapuscrit d'une autre main : SGA 3,
n\textsuperscript{os} 10 et 11. Limites projectives de préschémas en groupes,
puis préschémas en groupes affines, jusqu'aux deux théorèmes de Chevalley. Hors
sujet, et lu ici brièvement.

\item \textbf{Pages 59 et 60}. Le premier morceau du dossier qui touche
vraiment au sujet : une catégorie tannakienne dont le lien est un groupe
diagonalisable est classée par un homomorphisme du groupe des caractères dans
le groupe de Brauer, et elle est reconstruite à partir des algèbres d'Azumaya
correspondantes.

\item \textbf{Pages 61 à 65}, « Etude numérique des catégories tannakiennes ».
Les quatre \emph{caractères numériques} d'une $\otimes$-catégorie, et le
résultat négatif : ils ne la déterminent pas.

\item \textbf{Pages 66 à 83}, « Catégories tannakiennes définies par des
cristaux ». Le calcul le plus poussé du dossier. Aboutit à la structure
complète de la gerbe des $F$-isocristaux sur un corps fini, et à Dieudonné et
Manin récrits.

\item \textbf{Pages 85 et 86}, un feuillet d'un tapuscrit sur les topos
classifiants, et au dos le recensement des types simples. Second intrus.

\item \textbf{Pages 88 à 98}, de sa main, sous sa chemise « Filtration des
foncteurs fibres ». Ce qu'est une filtration admissible, ce que le gradué
associé est, et la géométrie des sous-groupes paraboliques qui en sort.

\item \textbf{Pages 101 à 108}, « Notes Saavedra », adressées à Saavedra et le
nommant. Six exemples : graduations, filtrations, pré-structures de Hodge,
structures de Hodge réelles, critère de profinitude, et la polarisation
renvoyée à plus tard.

\item \textbf{Pages 110 à 112}, de sa main. Les catégories de modules d'un
schéma en groupes affine, et la question de savoir dans quelle mesure le groupe
s'en récupère.

\item \textbf{Pages 114 à 116}. Le plan du traité en neuf chapitres, et la
table des matières détaillée du chapitre 1. Ce sont les deux feuillets qui
disent ce que le dossier voulait devenir.

\item \textbf{Pages 117 à 136}, tapuscrit « $\otimes$-catégories ». Le cœur :
les axiomes, les $\otimes$-foncteurs, la dissociation des contraintes, le
dictionnaire entre catégorie et bigèbre, les foncteurs fibres tordus, les
données de symétrie, le changement de base, les exemples instructifs, et un
théorème de liberté rapporté par Husemöller.

\item \textbf{Pages 138 et 139}, de sa main : les motifs à coefficients dans un
corps de nombres, et le calcul du dual.

\item \textbf{Pages 142 à 154}, de sa main : les anneaux de représentations. Le
formalisme $A^{G}$, $\boxtimes$, $T^{n}$, l'anneau gradué des groupes
symétriques, puis $R(T)$ pour un tore et l'homomorphisme caractéristique.

\item \textbf{Pages 156 à 179}, de sa main, d'un trait rapide : le dictionnaire
refait de rien. Une classe abélienne de $k$-modules, l'algèbre topologique
qu'elle reconstruit, les structures supplémentaires une par une, le groupe des
éléments de type groupe, et le schéma à trois termes qui clôt la rédaction.

\item \textbf{Pages 180 à 197}, de sa main : le groupe ordonné des classes
indécomposables, l'hypergroupe discret, et les trois anneaux calculés ---
$R(k)$, $R(k^{*})$, $R(k \rtimes k^{*})$.
\end{itemize}

\subsection*{Le dictionnaire, construit deux fois}

Il faut le dire ici, parce que c'est la seule chose qui rende le dossier
lisible comme un tout. La correspondance
\[
\otimes\text{-catégorie} \quad\longleftrightarrow\quad
\text{bigèbre} \quad\longleftrightarrow\quad
\text{schéma en groupes affine}
\]
est établie \emph{deux fois}, indépendamment.

La première fois aux pages 117 à 128, tapée, en six pages sèches : on part
d'une catégorie abélienne $k$-linéaire $\mathcal{C}$ munie d'un foncteur
$F \colon \mathcal{C} \to \operatorname{Modf}(k)$ fidèle et exact, on pose
$U = \operatorname{End}(F)$, on observe que $\mathcal{C}$ s'identifie aux
$U$-modules continus de dimension finie, et on traduit chaque contrainte.

La seconde fois aux pages 156 à 179, à la main, en repartant de zéro : la même
algèbre $U$ est reconstruite comme le sous-anneau fermé de $\prod_{i}
\operatorname{End}_{k}(A_{i})$ formé des familles cohérentes, les
$U$-modules utilisables sont appelés \emph{permis}, et les structures
supplémentaires sont ajoutées une par une jusqu'au schéma de la page 179.

Ce ne sont pas deux rédactions d'un même texte. La première est le dictionnaire
énoncé ; la seconde est le dictionnaire \emph{démontré}, avec ses topologies,
ses conditions de finitude et son groupe. Et c'est la seconde qui produit ce
dont les quinze dernières pages ont besoin : le groupe $G$ des éléments de type
groupe, et le fait que la catégorie soit celle des $G$-modules dont les
coefficients tombent dans $A$.

\subsection*{Trois « polarisations », et trois $\Gamma$}

Le mot \emph{polarisation} sert quatre fois dans le dossier, pour quatre
choses : la polarisation d'une variété projective (le choix d'un plongement, ou
d'une classe de diviseur ample) ; la polarisation d'un motif de poids $n$, qui
est une forme $M \otimes M \to T(n)$ « définie positive » ; la polarisation
d'une pré-structure de Hodge, page 106, qui est la même chose transportée en
théorie de Hodge ; et l'objet axiomatique $P$ des pages 6 à 16, qui est la
donnée, pour chaque objet, de l'\emph{ensemble} des formes qu'on décide
d'appeler positives. C'est ce dernier sens qui est le sujet du dossier ; les
trois autres sont ses modèles.

Et le mot \emph{effectivité} sert deux fois : l'effectivité d'un
$F$-isocristal, page 67, qui est une condition de bornitude sur les itérés du
Frobenius ; et le chapitre 9 du plan, page 114, « Questions d'effectivité (?) »,
qui est la même notion prise comme structure supplémentaire abstraite.

Pour les trois $\Gamma$, les quatre $A$ et les trois $\xi$, voir le
commentaire de tête du fichier : ce sont les homonymies qui font le plus de
dégâts, et cette lecture les distingue par la notation.

\subsection*{Ce qui n'est pas établi}

Le dossier est un atelier, et un atelier se juge aussi à ce qu'il laisse en
plan. Neuf choses sont annoncées et non faites, et elles sont enregistrées ici
plutôt que franchies :

\begin{enumerate}
\item \textbf{Le traité.} Les neuf chapitres de la page 114 ne sont écrits
nulle part dans le dossier. Le chapitre 1 a sa table des matières, pages 115 et
116 ; les pages 117 à 136 en sont un brouillon partiel.
\item \textbf{Page 121.} « A prouver que lorsqu'on prend $\mathcal{C}$ comme
dessus à la Koko, il n'y a pas de foncteur fibre sur aucune extension $K$ de
$k$ », avec son propre « ??? » en marge. Non prouvé.
\item \textbf{Page 128.} Peut-on construire, à partir de la donnée de symétrie
$\sigma$, l'élément $\tau$ qui la transporte après changement de base ? La
seule voie qu'il voit est l'invariance $U^{+}\!\cdot\tau = 0$. Non tranché.
\item \textbf{Page 132.} Quatre énoncés « A prouver » sur
$\operatorname{End}(\mathbf{1})$ absolument plat, la présentation finie des
$\operatorname{Hom}$, et la décomposition de $\mathcal{C}$ en limite inductive.
\item \textbf{Page 134.} « Théorème de ——— (rapporté par Husemöller) » : le nom
est laissé en blanc par la machine, et cette lecture le laisse en blanc.
\item \textbf{Page 151.} « Il faudrait prouver que $\mathbf{Z}[X(T)] \to$
l'anneau de Chern attaché à $R(T)$ ».
\item \textbf{Page 154.} L'homomorphisme $CK(G) \to \mathbf{Z}[X(T)]^{W}$ est
« sans doute » bijectif en l'absence de torsion. Non prouvé.
\item \textbf{Pages 89 et 92.} Que le pseudo-torseur $Q$ admette localement une
section. « Cependant, il faut démontrer que loc. $Q$ admet une section ». La
page 104 établit à la place que $H'$ est pro-unipotent, ce qui suffit pour une
catégorie à engendrement dénombrable.
\item \textbf{Page 98.} La Proposition en caractéristique $p > 0$ : « il faut
sans doute travailler un peu \ldots ».
\end{enumerate}

Et quatre questions qu'il pose en toutes lettres sans y répondre : l'image
essentielle du foncteur de Betti-Hodge, page 107 (« on n'a aucune idée ») ;
quelles $\otimes$-catégories proviennent d'un schéma en groupes affine, page
112, avec son « ??? » ; une caractérisation intrinsèque des irréductibles sur
$T/W$, page 182 ; et la signification de l'invariant « nombre de cônes
riemanniens dans $\Phi$ », page 30.

\pagerange{3}{16}
\section*{Polarisations : l'axiomatique (pages 3 à 16)}

\subsection*{Formes de Weil}

Soit $\mathcal{M}$ une $\otimes$-catégorie rigide sur un corps $K \subset
\mathbf{R}$, graduée par $\mathbf{Z}$ et munie d'un objet de Tate $K(1)$
inversible de degré $-2$. Pour $M$ de degré $i$ on considère les formes
\[
\varphi \colon M \otimes M \longrightarrow K(-i) ,
\]
$(-1)^{i}$-symétriques. Une telle $\varphi$, si elle est non dégénérée,
définit sur $\operatorname{End}(M)$ une \emph{anti-involution} $u \mapsto
u^{*}$ par
\[
\varphi(ux, y) = \varphi(x, u^{*}y) ,
\qquad\text{soit}\qquad
u^{*} = \tilde{\varphi}^{-1} \circ \check{u} \circ \tilde{\varphi} ,
\]
où $\tilde{\varphi} \colon M \xrightarrow{\ \sim\ } \check{M} \otimes K(-i)$
est l'isomorphisme que la non-dégénérescence fournit. La bijection
$\operatorname{Hom}(M \otimes M, K(-i)) \simeq \operatorname{End}(M)$,
$\psi \leftrightarrow u$ pour $\psi(x,y) = \varphi(x, uy)$, transforme la
transposition $\psi \mapsto {}^{t}\psi$ en $u \mapsto \varepsilon u^{*}$ avec
$\varepsilon = (-1)^{i}$ : les formes $\varepsilon$-symétriques correspondent
donc aux $u$ tels que $u^{*} = \varepsilon u$.

$\varphi$ satisfait la \emph{condition de Weil} si
\[
\operatorname{Tr}(u^{*}u) > 0 \qquad\text{pour tout } u \in
\operatorname{End}(M) ,\ u \neq 0 .
\]
Autrement dit, la forme trace $(u,v) \mapsto \operatorname{Tr}(u^{*}v)$ est
définie positive sur $\operatorname{End}(M)$ : l'algèbre involutive
$(\operatorname{End}(M), {}^{*})$ est \emph{formellement réelle}. Lorsque
$\operatorname{End}(M)$ est de dimension finie sur $K$, cela entraîne qu'elle
est semi-simple, et que la restriction de $\varphi$ à tout sous-objet est
encore de Weil.\footnote{La page 4 ajoute en marge, cerclée et reliée à la
parenthèse, l'assertion que $M$ est alors semi-simple. Elle demande deux
hypothèses de plus, que la marge donne en abrégé et que cette lecture
explicite : que $M$ soit de longueur finie, et que tout sous-objet admette un
supplément direct.}

\subsection*{La compatibilité n'est pas transitive}

Deux formes de Weil $\varphi$ sur $M$ et $\psi$ sur $N$, de même degré, sont
dites \emph{compatibles} si $\varphi \oplus \psi$ est encore de Weil, ce qui
revient à
\[
\operatorname{Tr}\bigl( u^{*(\varphi,\psi)} u \bigr) > 0
\qquad\text{pour tout } u \colon M \to N ,\ u \neq 0 ,
\]
l'adjoint étant calculé relativement à $\varphi$ et $\psi$. La relation est
réflexive et symétrique. Elle \emph{n'est pas transitive}, et la raison est
donnée dans la marge de la page 4 : si $M$ et $N$ sont tels que
$\operatorname{Hom}(M,N) = \operatorname{Hom}(N,M) = 0$, la condition est vide,
de sorte que deux formes de Weil $\varphi$, $\varphi'$ quelconques sur $M$ sont
toutes deux compatibles à une même $\psi$ sur $N$ sans avoir aucune raison
d'être compatibles entre elles.

Ce qui est vrai, et c'est ce sur quoi tout repose, est que sur les formes de
Weil d'un objet \emph{fixé} la compatibilité est une relation d'équivalence.
Le calcul est celui de la page 5 : si $\psi = \varphi(\cdot, \alpha\,\cdot)$
avec $\alpha \in \operatorname{End}(M)$ inversible, alors
\[
u^{\mathrm{ad}(\psi)} = \alpha^{-1} u^{*} \alpha ,
\qquad
u^{\mathrm{ad}(\varphi,\psi)} = u^{*}\alpha ,
\]
de sorte que $\varphi$ et $\psi$ sont compatibles si et seulement si
$\operatorname{Tr}(u^{*}u\alpha) > 0$ pour $u \neq 0$ ; et cette même condition
dit que $\psi$ est elle-même de Weil. Les formes de Weil compatibles à
$\varphi$ correspondent donc bijectivement aux $\alpha$ vérifiant cette
inégalité, qui forment un cône convexe ouvert --- et la transitivité s'ensuit
parce que la condition est symétrique en les deux formes.\footnote{L'insertion
cerclée de la page 5, « Donc, dans l'ens.\ des formes (de Weil) sur $M$, la
relation de compatibilité est une relation d'équivalence », est exactement cet
énoncé. Il est essentiel pour la Proposition qui suit : sans lui, l'unicité de
la prépolarisation maximale tomberait.}

\subsection*{Prépolarisations, admissibilité, polarisations}

Une \emph{prépolarisation} de $\mathcal{M}$ est la donnée, pour chaque objet
$M$ de degré $i$, d'une partie non vide
\[
P(M) \subset \operatorname{Hom}(M \otimes M, K(-i))
\]
formée de formes $(-1)^{i}$-symétriques de Weil, soumise à l'axiome
\[
\text{(P1)}\qquad
\varphi \in P(M) ,\ \psi \in P(N) ,\ u \colon M \to N \ (u \neq 0)
\ \Longrightarrow\ \operatorname{Tr}\bigl(u^{\mathrm{ad}(\varphi,\psi)}u\bigr) > 0 ,
\]
pour $M$, $N$ de même degré. Cela signifie que, pour chaque degré, les formes
de $P$ sont mutuellement compatibles. L'existence d'une prépolarisation force
$\mathcal{M}$ à être semi-simple.

\textbf{Proposition (page 5).} Toute prépolarisation $P$ est contenue dans une
unique prépolarisation maximale $\widehat{P}$, donnée par
\[
\widehat{P}(M) = \{\,\varphi \text{ de Weil sur } M \ \text{compatible à
toutes les formes de } P(M)\,\} .
\]

C'est une opération de \emph{polarité} : $\widehat{P}$ est l'orthogonal de $P$
pour la compatibilité, et $\widehat{\widehat{P}} = \widehat{P}$. Trois
corollaires suivent : $\widehat{P}$ est stable par somme directe de deux
objets de même degré, par restriction à un sous-objet, et par passage au dual,
$\check{\varphi} \in \widehat{P}(\check{M}(-i))$.\footnote{Sur la page, le
troisième corollaire a son argument surchargé et la lecture $\check{M}(-i)$ y
est douteuse ; c'est le degré qui la commande, $\check{M}$ étant de degré
$-i$.}

L'admissibilité est ce qui accroche entre eux les degrés, que (P1) laisse
indépendants :
\[
\text{(P2)}\quad \chi \in P(M \otimes N) \Longrightarrow
\chi \text{ compatible à } \varphi \otimes \psi
\qquad
\text{(P3)}\quad K(i) \otimes K(i) \xrightarrow{\ \sim\ } K(2i)
\ \in\ \widehat{P}(K(i)) .
\]
En présence de (P2), (P3) équivaut à : $\varphi \in P(M) \Rightarrow
\varphi(i) \in \widehat{P}(M(i))$. Une \emph{polarisation} est une
prépolarisation maximale et admissible. Les formes d'une polarisation sont
alors stables par $\oplus$, par restriction, par $\otimes$, par dualité, par
torsion à la Tate, et contiennent les isomorphismes canoniques $K(i) \otimes
K(i) \simeq K(2i)$ ainsi que leurs multiples positifs.

\subsection*{Les trois réductions, et la chaîne $\alpha \Rightarrow \beta
\Rightarrow \gamma$}

Trois remarques réduisent le problème (pages 7 et 8). D'abord, une polarisation
est connue dès qu'on la connaît sur un système de sous-$\otimes$-catégories qui
engendre $\mathcal{M}$ avec $K(1)$ et $K(-1)$. Ensuite, les formes de Weil de
$M$ sont celles de $M \otimes_{K} \mathbf{R}$, de sorte que les polarisations
de $\mathcal{M}$ et de $\mathcal{M}_{\mathbf{R}}$ se correspondent : on est
ramené aux $\otimes$-catégories \emph{sur $\mathbf{R}$}. Enfin, si
$\mathcal{M} = \varinjlim \mathcal{M}(\alpha)$ le long de
sous-$\otimes$-catégories pleines croissantes, alors
\[
\operatorname{Pol}(\mathcal{M}) = \varprojlim \operatorname{Pol}(\mathcal{M}(\alpha)) ,
\]
et si les $\mathcal{M}(\alpha)$ sont de type fini,
$\operatorname{Pol}(\mathcal{M}) \neq \emptyset$ si et seulement si chaque
$\operatorname{Pol}(\mathcal{M}(\alpha)) \neq \emptyset$. On est ainsi ramené
aux $\otimes$-catégories semi-simples de type fini sur $\mathbf{R}$.

Les pages 8 et 9 passent aux objets complexifiés : $M \in
\operatorname{Ob}\mathcal{M}^{i}_{\mathbf{C}}$, c'est-à-dire $M_{0}$ dans
$\mathcal{M}^{i}$ muni d'un endomorphisme $J_{M}$ de carré $-1$, et
$\overline{M}$ le même avec $J_{\overline{M}} = -J_{M}$. Une forme $M
\otimes_{\mathbf{C}} \overline{M} \to \mathbf{C}(-i)$ équivaut à une forme
$\varphi_{0} \colon M_{0} \otimes M_{0} \to \mathbf{R}(-i)$ vérifiant
$\varphi_{0}(J x, y) = -\varphi_{0}(x, J y)$ et $\varphi_{0}(J x, J y) =
\varphi_{0}(x,y)$ ; elle est dite hermitienne ou antihermitienne selon que
$\varphi_{0}$ est symétrique ou alternée. Le point est un argument de
\emph{moyennisation} : si $\psi_{0}$ est une polarisation de $M_{0}$, alors
\[
\varphi_{0}(x,y) = \psi_{0}(x,y) + \psi_{0}(J_{M}x, J_{M}y)
\]
en est une, et elle est $J_{M}$-invariante. Il s'ensuit
\[
\check{M} \simeq \overline{M}(-i)
\qquad\text{pour tout } M \in
\operatorname{Ob}\mathcal{M}^{i}_{\mathbf{C}} .
\]
D'où la chaîne d'implications de la page 9 :
\begin{quote}
$\alpha$) $\mathcal{M}$ admet une polarisation sur $\mathbf{R}$ ;

$\Downarrow$

$\beta$) $\mathcal{M}$ est semi-simple et $\check{M} \simeq \overline{M}(-i)$
pour tout $M$ de $\mathcal{M}^{i}_{\mathbf{C}}$ ;

$\Downarrow$

$\gamma$) $\mathcal{M}$ est semi-simple et $\check{M} \simeq M(-i)$ pour tout
$M$ de $\mathcal{M}^{i}_{\mathbf{R}}$.
\end{quote}
La seconde implication vient de ce que $M_{\mathbf{C}} \simeq N_{\mathbf{C}}$
entraîne $M \simeq N$. Et \emph{aucune des deux ne se renverse} : pour
$\mathcal{M} = \operatorname{Rep}(G)$ avec $G$ algébrique sur $\mathbf{R}$, on
peut avoir $\beta$) sans $\alpha$) en prenant $SO(\varphi)/\text{centre}$ pour
$\varphi$ une forme en un nombre pair de variables avec un nombre impair de
carrés négatifs.\footnote{Le contre-exemple pour $\gamma$) sans $\beta$) est
annoncé entre crochets sur la page et ne s'y lit pas.}

\subsection*{Le critère, sur un système de générateurs}

Supposons $\mathcal{M}$ semi-simple sur $K \subset \mathbf{R}$, engendrée comme
$\otimes$-catégorie par $M_{1}, \ldots, M_{r}$ de degrés $d_{1}, \ldots,
d_{r}$ et par $K(\pm 1)$, et donnons-nous sur chaque $M_{i}$ une forme de Weil
$\varphi_{i}$, $(-1)^{d_{i}}$-symétrique. Pour un système d'entiers
$m = (m_{1}, \ldots, m_{r}, m_{r+1})$ avec $m_{1}, \ldots, m_{r} \geqslant 0$,
posons
\[
M_{m} = \Bigl( \bigotimes_{i=1}^{r} \otimes^{m_{i}} M_{i} \Bigr)(m_{r+1}) ,
\qquad
\varphi_{m} = \Bigl( \bigotimes_{i=1}^{r} \otimes^{m_{i}} \varphi_{i} \Bigr)(m_{r+1}) ,
\qquad
\delta(m) = -2 m_{r+1} + \sum_{i=1}^{r} m_{i} d_{i} .
\]

\textbf{Critère (pages 10 et 11).} Pour qu'il existe une polarisation de
$\mathcal{M}$ dont les $\varphi_{i}$ soient des formes de polarisation, il faut
et il suffit que
\[
\operatorname{Tr}\bigl( u^{*(\varphi_{m}, \varphi_{m'})} u \bigr) \geqslant 0
\qquad\text{pour tous } m, m' \text{ avec } \delta(m) = \delta(m')
\text{ et tout } u \in \operatorname{Hom}(M_{m}, M_{m'}) .
\]

La suffisance est la construction directe : tout $M \neq 0$ de degré $n$ se
plonge dans un $M_{m}$ de même degré, on pose $P(M)$ = l'ensemble des formes
compatibles à $\varphi_{m}|M$, et l'hypothèse dit exactement que ce $P(M)$ ne
dépend pas du plongement choisi. Que ce soit une prépolarisation, et qu'elle
soit admissible, se lit sur $\varphi_{m} \otimes \varphi_{m'} =
\varphi_{m+m'}$ et sur $\varphi_{m}(i) = \varphi_{(m_{1}, \ldots, m_{r},
m_{r+1}-i)}$.\footnote{Le signe de l'inégalité est $\geqslant$ sur la page,
non $>$, et l'annotation « $\delta(m)$ » y est portée au-dessus des deux
sommes ; le degré de $K(1)$ étant $-2$, la condition de même degré est bien
celle écrite ici.}

\subsection*{$\mathcal{M} = \operatorname{Rep}(G)$, et l'involution de Cartan}

Les six dernières pages du faisceau sont l'exemple qui commande tout : $G$ un
groupe algébrique sur $\mathbf{R}$ muni de
\[
\mathbf{G}_{m} \xrightarrow{\ i\ } G \xrightarrow{\ \varepsilon\ } \mathbf{G}_{m} ,
\qquad \eta = i(-1) ,
\]
et $C \in G(\mathbf{R})$ avec $C^{2} = \eta$. Pour $M$ une représentation de
poids $i$, on dit que $\varphi \colon M \otimes M \to \mathbf{R}(-i)$ est une
\emph{$C$-polarisation} si
\[
\varphi^{C}(x,y) := \varphi(x, Cy)
\]
est symétrique et définie positive. Le calcul essentiel, page 11, est que pour
$u \colon M \to N$ équivariant
\[
u^{*(\varphi^{C}, \psi^{C})} = C_{M}^{-1}\, u^{*(\varphi,\psi)}\, C_{N} ,
\qquad\text{d'où}\qquad
\operatorname{Tr}\bigl( u\, u^{*(\varphi,\psi)} \bigr)
= \operatorname{Tr}\bigl( u\, u^{*(\varphi^{C},\psi^{C})} \bigr) \geqslant 0 ,
\]
de sorte que la positivité pour $\varphi^{C}$ donne la condition de Weil pour
$\varphi$. Le produit tensoriel de deux $C$-polarisations en est encore une
\emph{pourvu que $\varepsilon(C) = 1$}, et non $-1$ : c'est une insertion de sa
main, et elle est nécessaire, puisque $\varepsilon$ mesure l'action de $C$ sur
l'objet de Tate. Sous cette condition, l'application qui à $M$ associe
l'ensemble de ses $C$-polarisations \emph{est} une polarisation de
$\operatorname{Rep}(G)$.

Le théorème de la page 13 est, dans le vocabulaire d'aujourd'hui, que $C$
définit une \emph{involution de Cartan} :

\textbf{Proposition.} Si $C \in G'(\mathbf{R})$, $G' = \operatorname{Ker}
\varepsilon$, avec $C^{2} = \eta$, définit une polarisation de
$\operatorname{Rep}(G)$, alors $\operatorname{Centr}_{G'}(C)(\mathbf{R})$ est
un sous-groupe compact maximal de $G'(\mathbf{R})$.

La compacité est immédiate : plongeant $G'$ dans un produit de
$\operatorname{Aut}_{\mathbf{R}}(M_{\alpha})$ pour des représentations où
$\eta$ opère par $\pm 1$, le centralisateur de $C$ laisse invariantes les
formes définies positives $\varphi^{C}_{\alpha}$, donc tombe dans un produit de
groupes orthogonaux compacts. La maximalité est l'argument long de la page 13,
qui se ramène à montrer que si $G'$ est compact alors $C$ est central : on
décompose $V = V^{+} \oplus V^{-}$ selon les valeurs propres $\pm 1$ de $C$,
la forme est positive sur $V^{+}$ et négative sur $V^{-}$, on conjugue $G'$
dans le compact maximal $O(V^{+}) \times O(V^{-})$ de $O(\varphi)$, et $C =
C^{+} \oplus C^{-}$ y devient $\mathrm{id} \oplus (-\mathrm{id})$, qui est
central.\footnote{C'est, sous une autre forme, la caractérisation devenue
standard : une polarisation d'une structure de Hodge fournit une involution de
Cartan $g \mapsto CgC^{-1}$ du groupe adjoint réel, le compact maximal étant
son groupe des points fixes, et $C$ n'est autre que l'image de $i$ par le
morphisme du tore de Deligne. Le dossier ne dispose d'aucune de ces
formulations et démontre l'énoncé directement.} Le cas $\eta \neq 1$ se traite
par le caractère $\varepsilon \colon G \to \mu_{2}$ défini par $gCg^{-1} =
\varepsilon(g) C$ : si $\varepsilon(g) = -1$ on aurait
$\varphi^{C}(gx,gx) = -\varphi^{C}(x,x)$, les deux membres respectivement $> 0$
et $< 0$, absurde.\footnote{Ce $\varepsilon$ n'est pas celui de la page 11 ; la
page réemploie la lettre.}

\subsection*{L'ensemble des polarisations est un torseur}

Deux éléments $C$, $C'$ de carré $\eta$ définissent la même polarisation dès
qu'ils sont conjugués sous $G(\mathbf{R})$. Réciproquement (page 15), en
conjuguant par $G'(\mathbf{R})^{\circ}$ on peut supposer $C' = zC$ avec $z$
dans la $2$-torsion du centre $Z'(\mathbf{R})$ ; et si $z \neq 1$, il existe
une représentation $V \neq 0$ où $z$ opère par $-\mathrm{id}$, donc
$\varphi^{C'} = -\varphi^{C} < 0$, absurde. Ainsi :
\begin{quote}
l'ensemble des polarisations de $\operatorname{Rep}(G)$ définissables par un
$C$ --- celles qu'il appelle \emph{hodgiennes} --- est un torseur sous la
$2$-torsion de $Z(\mathbf{R})$.
\end{quote}
L'action est explicite : $(zP)$ est la polarisation dont les formes sont les
$\varphi$ telles que $\varphi^{z}$ soit une $P$-polarisation.

Le cas général, page 16, est traité sans supposer $\mathcal{M} =
\operatorname{Rep}(G)$ : pour $z$ un $\otimes$-automorphisme de
$\mathrm{id}_{\mathcal{M}}$ de carré $1$, on pose $\varphi^{z}(x,y) =
\varphi(x, zy)$ et l'on trouve
\[
u^{*(\varphi^{z}, \psi^{z})} = z_{M}^{-1} u^{*(\varphi,\psi)} z_{N}
= u^{*(\varphi,\psi)} ,
\]
la seconde égalité parce que $z$, étant un automorphisme du foncteur identique,
commute à tout. Donc $\varphi$ est de Weil si et seulement si $\varphi^{z}$
l'est, la compatibilité se transporte, et $\varphi^{z} \otimes \psi^{z} =
(\varphi \otimes \psi)^{z}$. Si de plus $z$ opère par l'identité sur $K(1)$,
alors ${}^{z}P$ est encore une polarisation, et elle est distincte de $P$ dès
que $z \neq 1$ : la $2$-torsion de ce groupe d'automorphismes \emph{opère
librement} sur $\operatorname{Pol}(\mathcal{M})$. Reste la question, posée et
non tranchée : cette action est-elle transitive ?\footnote{La condition
$z^{2} = 1$ sert, comme sa marge le note, à ce que $\varphi$
$\varepsilon$-symétrique entraîne $\varphi^{z}$ $\varepsilon$-symétrique.}

\pagerange{18}{30}
\section*{Polarisations : les « Notes anciennes » (pages 18 à 30)}

\subsection*{L'adjoint, et l'involution déterminée mod intérieurs}

Ces feuillets, d'un papier jauni et d'une encre noire, portent le titre « Notes
anciennes » de sa main. Ils prennent la même matière par l'autre bout : non
plus l'axiomatique du choix des formes positives, mais l'algèbre
$\operatorname{End}(M)$ et les cônes qui vivent dedans.

Soit $M$ un motif de poids $p$ et $\varphi$ une forme non dégénérée à valeurs
dans $\mathbf{Q}(p)$. Posant $M^{*} = \check{M}(p)$ et
$\mathcal{Q}(M) = B(M, M ; \mathbf{Q}(p))$ l'espace des formes bilinéaires, on
a les identifications
\[
\operatorname{End}(M)
\xrightarrow[\ \sim\ ]{\ \Theta_{\varphi}\ }
\mathcal{Q}(M) \simeq \operatorname{Hom}(M, M^{*}) ,
\qquad
\Theta_{\varphi}(u) = \varphi \circ (u \otimes 1) ,
\]
et la symétrie de $\mathcal{Q}(M)$ devient, transportée sur
$\operatorname{End}(M)$, l'\emph{adjoint} $u \mapsto u^{\varphi}$. Le signe
$(-1)^{\delta}$ que la page 20 écrit devant $\tilde{\varphi}^{-1}$ est celui
qui rend l'identité valable dans les deux cas, $\varphi$ symétrique avec $p$
pair et $\varphi$ alternée avec $p$ impair ; l'exposant $\delta$ n'est pas
défini sur ces deux pages, et les deux dernières lignes de la page 20
aboutissent à la même égalité pour les deux cas, ce que le signe est justement
censé éviter.\footnote{La page 20 est donnée telle quelle par la
transcription : les deux arguments du membre de droite de
$\varphi(u^{\varphi}x,y) = (-1)^{\delta}\varphi(ux,x)$ y sont biffés d'une
croix, et les cas symétrique et antisymétrique s'achèvent tous deux sur
$\varphi(ux,y) = \varphi(x, u^{\varphi}y)$. Il manque là un signe, et il n'est
pas rétabli.}

Les quatre règles de la page 22 sont que $u \mapsto u^{\varphi}$ est additive,
$\mathbf{Q}$-linéaire, involutive et anti-multiplicative. Remplaçant $\varphi$
par $\psi = \varphi a$ avec $\psi(x,y) = \varphi(ax,y)$, la condition que
$\psi$ soit du même type s'écrit $a^{\varphi} = a$, et alors
\[
u^{\psi} = a^{-1} u^{\varphi} a .
\]
Donc \emph{l'involution de $A = \operatorname{End}(M)$ est déterminée modulo
automorphismes intérieurs} --- et plus précisément, $\varphi$ étant fixée,
modulo les intérieurs par un élément $\varphi$-hermitien inversible. La page 23
ajoute la remarque juste : la donnée de $\varphi$ est \emph{plus fine} que
celle de l'involution, puisque si $\operatorname{End}(M)$ est commutatif toutes
les $\varphi$ donnent la même involution, et $\lambda\varphi$ donne toujours la
même que $\varphi$.

\subsection*{Algèbre de Riemann}

La forme
\[
\langle u, v \rangle = \operatorname{Tr}(v^{\varphi} u) = \operatorname{Tr}(u^{\varphi} v)
\]
est symétrique sur $A$ et s'identifie à $\check{\varphi} \otimes \varphi$ sous
$A \simeq \check{M} \otimes M$. Quand elle est définie positive --- elle est
nécessairement non dégénérée, donc définie dès qu'elle est positive --- il
appelle $A$ une \emph{algèbre de Riemann}. C'est, dans le vocabulaire
d'aujourd'hui, une algèbre semi-simple de dimension finie munie d'une
\emph{involution positive}.\footnote{La notion et sa classification sont celles
d'Albert ; l'exemple qui la commande, et celui d'où le dossier part, est
$\operatorname{End}(A) \otimes \mathbf{Q}$ d'une variété abélienne munie de
l'involution de Rosati attachée à une polarisation. Le dossier ne nomme ni
l'une ni l'autre.}

\subsection*{Le cône est la donnée intrinsèque}

Voici le point de ces pages, et il est souligné sur la page 25. Soit $A^{+}$
l'adhérence du cône des combinaisons à coefficients $\geqslant 0$ des
$u^{\varphi}u$. Transporté dans l'espace $\Phi^{s}$ des formes symétriques, il
devient le cône fermé $\Phi^{+}$ engendré par les
\[
\varphi^{u}(x,y) = \varphi(ux, uy) ,
\]
c'est-à-dire \emph{le plus petit cône convexe fermé contenant $\varphi$ et
stable par les automorphismes de $M$}. La théorie des algèbres de Riemann donne
alors : tout élément intérieur à $A^{+}$ est une unité, donc tout élément
intérieur à $\Phi^{+}$ est une forme non dégénérée ; et si $\psi$ est intérieur
à $\Phi^{+}$, l'involution $u \mapsto u^{\psi} = a^{-1}u^{\varphi}a$ s'en
déduit par un automorphisme intérieur et a donc les mêmes propriétés de
positivité, d'où $\Phi^{+}(\psi) = \Phi^{+}(\varphi)$. Conclusion :
\begin{quote}
« C'est la donnée du cône $\Phi^{+}$ qui, intrinsèquement, constitue une
structure supplémentaire importante. »
\end{quote}
Et il note aussitôt la limite de la chose : parmi les cônes possibles, on ne
peut en distinguer de meilleurs, puisque si $\varphi$ est positive $-\varphi$
l'est aussi et définit $-\Phi^{+}$.

\subsection*{La polarité}

Deux énoncés de dualité ferment le faisceau. D'abord (page 25) $A^{+}$ est son
propre polaire dans $A^{s}$ pour la forme $\operatorname{Tr}(vw)$ : c'est un
cône \emph{autopolaire}. Ensuite (page 26), identifiant $A$ à $\Phi(M) =
\check{M} \otimes \check{M}(p)$ par $\Theta_{\varphi}$ d'une part, et à
$\Phi(M^{*}) = M \otimes M(-p)$ par $u \mapsto {}^{t}u$ et $\check{\varphi}$
d'autre part, la forme $\operatorname{Tr}(vw)$ devient l'accouplement évident
de contraction entre $\Phi(M)$ et $\Phi(M^{*})$, et
\[
\Phi^{+}(M) \quad\text{et}\quad \Phi^{+}(M^{*}) \quad\text{sont polaires l'un
de l'autre.}
\]
Chacun est donc le polaire de l'autre, et chacun détermine l'autre --- la
bipolarité valant parce que les cônes sont convexes et fermés. Le critère de la
page 27 en est la réciproque : si deux cônes convexes fermés $P \subset
\Phi(M)$, $P^{*} \subset \Phi(M^{*})$ sont stables par automorphismes,
d'accouplement $\geqslant 0$, et contiennent une paire $(\varphi,
\varphi^{*})$ de formes admissibles polaires l'une de l'autre, alors $P$ et
$P^{*}$ sont nécessairement admissibles et polaires.

\subsection*{Le cadre abstrait, et l'invariant laissé ouvert}

Les pages 27 et 28 dégagent le cadre : une algèbre $A$ semi-simple de type fini
sur $\mathbf{Q}$, munie d'une forme trace telle que $\operatorname{tr}(uv)$ soit
non dégénérée ; un module $\Phi$ sur $A \otimes A$, de rang $1$ pour chacune
des deux structures ; son dual $\Phi^{*}$ ; et une symétrie de $\Phi$ qui
échange les deux structures de $A$-module. Pour $\varphi$ fixé par la symétrie,
la bijection $A \to \Phi$ qu'il définit donne sur $A$ l'application $u \mapsto
u^{\varphi}$, involutive, déterminée modulo intérieurs par les $a$ avec
$a^{\varphi} = a$ --- et anti-multiplicative, ce que la marge de la page 28
note : $(uv)^{\psi} = v^{\psi}u^{\psi}$, indépendamment du choix de
$\psi$.\footnote{La page 28 est un brouillon rapide dont les deux tiers de la
prose de liaison sont perdus ; les formules relevées ici sont celles que la
transcription donne.}

Les pages 29 et 30 concluent. $\varphi$ est dite \emph{riemannienne} si $A$,
munie de la trace et de $u \mapsto u^{\varphi}$, est une algèbre de Riemann.
Toute $\varphi$ riemannienne définit le cône $\Phi^{+}(\varphi)$, et
$A^{\circ} \to \Phi^{*}$ transporte le cône dual en $\Phi^{*+}$ ; la théorie
des algèbres hilbertiennes donne que les deux cônes sont polaires, que leurs
intérieurs sont exactement les formes riemanniennes, et que $\Phi^{+}(\psi) =
\Phi^{+}(\varphi)$ pour $\psi$ intérieur. On trouve ainsi dans $\Phi$ un nombre
\emph{fini} de cônes convexes fermés $P_{\alpha}$, stables par les $u \otimes
u$, d'intérieurs deux à deux disjoints, et dont la réunion des intérieurs est
l'ensemble des formes riemanniennes. Si $\varphi$ est riemannienne, $-\varphi$
l'est aussi mais dans le cône opposé. D'où la question, qui clôt le faisceau et
reste ouverte :
\begin{quote}
« Quelle est la signification de l'invariant : nombre de cônes riemanniens dans
$\Phi$ ? »
\end{quote}
C'est un invariant discret du couple $(A, \Phi)$, et donc, pour l'algèbre des
endomorphismes d'une variété abélienne munie de son involution de Rosati, un
invariant de son type au sens de la classification des involutions positives.
Le dossier ne le calcule dans aucun cas.

\pagerange{32}{57}
\section*{Deux numéros de SGA 3 (pages 32 à 57)}

Vingt-six feuillets tapés par une autre main, corrigés par lui au crayon et à
l'encre bleue, ne sont pas sur le sujet du carton : ce sont les
n\textsuperscript{os} 10 et 11 de SGA 3, sur le passage à la limite projective
dans les préschémas en groupes et sur les préschémas en groupes affines. Rien
sur les feuillets ne dit pourquoi ils sont classés ici. On les lit brièvement,
pour deux résultats dont le travail tannakien du dossier ne peut pas se
passer.

\subsection*{Le premier : tout groupe affine est pro-algébrique}

\textbf{Proposition (page 51).} Soient $k$ un corps et $G$ un $k$-groupe affine
d'algèbre $A$. Alors $G$ est limite projective d'un système filtrant croissant
de $k$-groupes affines de type fini, dont les morphismes de transition sont
fidèlement plats.

La démonstration est en deux temps, et ils sont tous deux utiles ailleurs. Le
lemme 11.12 dit que tout $x \in A$ est contenu dans une sous-$k$-algèbre de
type fini $V$ stable par la diagonale $\Delta$ et par l'involution --- ce qui
se voit en écrivant $\Delta(x) = \sum \varphi_{i} \otimes a_{i}$ et
$\Delta(\varphi_{i}) = \sum \varphi_{j} \otimes b_{ji}$, d'où $\Delta(a_{i}) =
\sum b_{ij} \otimes a_{j}$, et en prenant pour $V$ l'algèbre engendrée par les
$b_{ij}$ et leurs images par l'involution. L'ensemble de ces sous-algèbres est
stable par engendrement, donc $A$ en est la limite inductive, chacune est
l'algèbre affine d'un $G_{i}$ de type fini, et $G = \varprojlim G_{i}$. Le lemme
11.14 donne enfin que les transitions sont fidèlement plates, parce qu'un
morphisme de $k$-groupes affines de type fini est fidèlement plat si et
seulement si l'homomorphisme d'algèbres correspondant est injectif.

C'est exactement ce dont les pages 108 et 112 auront besoin pour parler de
$G = \varprojlim G_{i}$ et de la composante neutre du pro-groupe, et ce que le
dictionnaire des pages 120 et 156 suppose quand il fait de $U$ une algèbre
\emph{profinie}.

\subsection*{Le second : affine, c'est agir fidèlement sur un espace vectoriel}

\textbf{Proposition 11.11 (pages 48 à 50).} Soit $G$ un groupe algébrique sur
un corps $k$. Sont équivalentes : (i) $G$ affine ; (ii) $G$ quasi-affine ;
(iii) $G$ opère fidèlement sur un $k$-schéma quasi-affine ; (iv) $G$ opère
fidèlement sur un $k$-espace vectoriel, pas nécessairement de dimension
finie ; (v) $G$ est isomorphe à un sous-groupe fermé d'un
$\mathrm{Gl}(n)_{k}$.

Le passage qui compte est (iv) $\Rightarrow$ (v), et il repose sur le
corollaire 11.10 : un module quasi-cohérent où $G$ opère est réunion filtrante
de ses sous-modules de type fini stables. Si $G$ opère fidèlement sur $V$, les
noyaux $K_{i}$ des actions induites sur les $V_{i}$ de dimension finie ont pour
intersection le groupe unité ; $G$ étant noethérien, l'un des $K_{i}$ est déjà
trivial, donc $G \to \underline{\operatorname{Aut}}_{V_{i}}$ est un
monomorphisme, donc une immersion fermée.\footnote{Le corollaire 11.10 lui-même
vient de la proposition 11.9, et la Remarque 11.10.1 crédite J.-P. Serre d'une
hypothèse affaiblie : il suffit que $\mathcal{E}$ devienne projectif sur une
base plate convenable, et sur une base régulière de dimension $\leqslant 1$ la
seule platitude de $G$ suffit. La référence, d'abord donnée à un exposé de
séminaire, est remplacée de sa main par l'article des \emph{Publ.\ Math.}}

Le cas non algébrique de (iv) $\Rightarrow$ (i) passe par la limite projective
$G' = \varprojlim G_{i}$ des plus petits sous-groupes fermés de
$\operatorname{Aut}(V_{i})$ par lesquels $G$ se factorise, et par le lemme
11.11.1 : tout sous-schéma en groupes \emph{ouvert} d'un groupe sur un corps
est \emph{fermé}, le complémentaire étant réunion de translatés de l'ouvert.

\subsection*{Ce que le reste contient}

Le n\textsuperscript{o} 10 est une machine de descente : dans la situation de
EGA IV 8.8, la catégorie des $S$-préschémas de présentation finie est déterminée
par celles des $S_{i}$, et un grand nombre de propriétés d'un morphisme se
lisent à un niveau fini. Les résultats 10.2 à 10.16 étendent cela aux groupes,
aux préschémas à groupe d'opérateurs, aux monomorphismes centraux et
invariants, à la représentabilité d'un quotient (lemmes 10.7 et 10.8, où la
représentabilité de $X/R$ équivaut à l'existence d'un $Y$ et d'un $p$
fidèlement plat rendant cartésien le carré évident), aux transporteurs et
normalisateurs (10.11), et au faisceau des commutateurs (10.12).

Y sont encore la proposition 10.6 et le corollaire 10.13, qui ramènent à un
niveau fini la représentabilité d'un quotient $G/G'$ et la condition
$(D_{s}, D_{s}) = D_{s}$ sur les fibres d'un sous-groupe invariant.

Le n\textsuperscript{o} 11 contient, outre les deux propositions ci-dessus, le
lemme 11.1 sur $\mathcal{O}(X) \otimes \mathcal{O}(Y) \to \mathcal{O}(X
\times_{S} Y)$ --- en marge duquel sa main écrit « fausse formule », et il a
raison de le signaler : l'isomorphisme demande une hypothèse de platitude, que
les trois cas a), b), c) fournissent chacun à sa façon ; la correction de la
page 44, qui rétablit le théorème de Lazard dans sa forme exacte --- tout
module plat est limite inductive filtrante de modules libres de type fini, là
où la machine avait tapé l'implication inverse et fausse ; les définitions
11.15 des vecteurs semi-invariants, de leur \emph{poids} et de leur caractère
associé, où le mot « poids » est déplacé de sa main sur « invariant » et
l'ancienne définition biffée ; et les deux théorèmes de Chevalley. Le premier
(11.16) : tout sous-groupe fermé $H$ de $G$ affine algébrique est le plus grand
sous-groupe sous lequel un nombre fini d'éléments indépendants de $A$, de même
poids, soient semi-invariants --- la démonstration prend $W = V \cap I$ de
dimension $d$ et travaille dans $\bigwedge^{d}V$. Le second (11.17) : pour $H$
fermé invariant dans $G$ affine, le faisceau quotient fpqc $G/H$ est
représentable par un groupe \emph{affine} ; le cas algébriquement clos est
supplié par un feuillet de sa main (page 55), qui est le « VOIR CI CONTRE »
cerclé de la page 56, et qui remarque en passant que la référence à
« BIBLE » est un corollaire de 11.16.

Entre les deux se placent les énoncés qui font marcher la machine et que cette
lecture ne détaille pas : le corollaire 11.2, qui fait de $X \rightsquigarrow
\operatorname{Spec}\mathcal{O}(X)$ un foncteur commutant aux produits finis,
donc transformant groupes en groupes ; la définition 11.3 de
l'\emph{enveloppe affine} $G_{\mathrm{af}} = \operatorname{Spec}
\mathcal{O}(G)$, et le n\textsuperscript{o} 11.6 qui met en bijection les
opérations de $G$ et celles de $G_{\mathrm{af}}$ sur un module, en donnant les
deux axiomes (CM1) et (CM2) sous forme géométrique et sous forme
coalgébrique ; la proposition 11.5, qui identifie les morphismes de
$S$-foncteurs de $X$ dans $\mathcal{H}om_{\mathcal{O}_{S}}
(\mathbb{V}(\mathcal{F}), \mathbb{V}(\mathcal{E}))$ aux homomorphismes
$\mathcal{E} \to \mathcal{O}(X) \otimes \mathcal{F}$, et dont la démonstration
est l'endroit où Lazard est rétabli ; le lemme 11.7, qui définit l'opération
\emph{induite} sur un sous-module stable ; et le lemme 11.8, qui est le cœur
technique de tout le paragraphe --- pour $\mathcal{O}(G)$ libre de base
$(\varphi_{i})$ et $\rho(x) = \sum \varphi_{i} \otimes x_{i}$, le sous-module
engendré par les $x_{i}$ est stable, de type fini, et contient $x$.

La remarque 11.11.2 note, et c'est ce dont les pages 108 et 112 se serviront,
qu'une fois $G$ su affine l'homomorphisme $G \to \varprojlim G_{i}$ est un
isomorphisme : $G$ est sa propre limite projective de groupes algébriques, ce
que 11.13 redémontre par une autre méthode.

Le dernier énoncé du tapuscrit, le corollaire 11.18 --- tout sous-groupe
algébrique d'un groupe diagonalisable est diagonalisable --- est annulé par des
boucles au crayon et le feuillet reste blanc dessous.

\pagerange{59}{60}
\section*{Une catégorie tannakienne de lien diagonalisable (pages 59 et 60)}

Deux feuillets d'un tapuscrit étranger au précédent, et les premiers du dossier
qui touchent son sujet.

Soit $\mathcal{M}$ une catégorie tannakienne sur un corps $k$ dont le lien est
un groupe diagonalisable $D(\Lambda)$, $\Lambda$ étant un groupe abélien à
action galoisienne triviale. Cela revient, la page le dit en insertion, à
demander que $\mathcal{M}$ soit semi-simple et que pour tout objet simple $S$
de rang $d(S)$ on ait $\operatorname{End}(S)$ de rang $d(S)^{2}$ : l'objet est
« isotypique », et son algèbre d'endomorphismes est une algèbre à division de
degré maximal.

Une telle catégorie est classée, à équivalence près respectant toutes les
structures et contraintes, par un élément
\[
\xi \in H^{2}(k, D(\Lambda)) \simeq \operatorname{Hom}(\Lambda, \operatorname{Br}(k)) ,
\]
donc par un homomorphisme $u \colon \Lambda \to
\operatorname{Br}(k)$.\footnote{L'isomorphisme utilise deux choses que la page
laisse implicites : que $\Lambda$ est un module galoisien trivial, sans quoi le
$H^{2}$ ne se calcule pas ainsi, et que $\operatorname{Br}(k) = H^{2}(k,
\mathbf{G}_{m})$, ce qui est le théorème pour un corps.}

La structure qui s'ensuit est entièrement lisible. $\mathcal{M}$ est graduée de
type $\Lambda$, les objets homogènes sont les isotypiques, il y a pour chaque
$i \in \Lambda$ exactement un simple $M_{i}$ de degré $i$, et son rang est
l'ordre $d(i)$ de $u(i)$ dans $\operatorname{Br}(k)$. L'anneau des classes est
donc
\[
K(\mathcal{M}) = \Bigl\{ \textstyle\sum n_{i} e^{i} \in \mathbf{Z}[\Lambda]
\ :\ d(i) \mid n_{i} \ \text{pour tout } i \Bigr\} \subset \mathbf{Z}[\Lambda] ,
\]
de base canonique les $d(i)e^{i}$, et $\operatorname{End}(M_{i})$ est le corps
gauche de centre $k$ d'invariant $u(i)$.\footnote{Que ce sous-groupe soit un
sous-anneau n'est pas dit sur la page et mérite un mot : $u$ étant un
homomorphisme, $u(i+j) = u(i)u(j)$, donc $d(i+j)$ divise $\operatorname{ppcm}
(d(i), d(j))$, qui divise $d(i)d(j)$, qui divise $n_{i}m_{j}$. La
multiplication reste donc dans le sous-groupe. Ce que la page écrit est
l'inclusion $K(\mathcal{M}) \subset \mathbf{Z}[\Lambda]$, qui s'obtient par
changement de base à la clôture ; la description ci-dessus est l'image.}

La page 60 reconstruit $\mathcal{M}$ à partir des algèbres. Pour $\Lambda =
\mathbf{Z}^{(I)}$ libre, la donnée de $u$ est celle d'éléments $\xi_{i} \in
\operatorname{Br}(k)$ ; choisissant un isotypique non nul $E_{i}$ de degré $i$
et posant $A_{i} = \operatorname{End}(E_{i})$, l'algèbre d'Azumaya de classe
$\xi_{i}$, on considère la catégorie sous-additive dont les objets sont les
$\underline{\gamma} \in \mathbf{N}^{(I)}$ --- lu comme $\bigotimes
E_{i}^{\gamma_{i}}$ --- avec
\[
\operatorname{Hom}(\underline{\gamma}, \underline{\gamma}')
= \begin{cases}
0 & \text{si } \underline{\gamma} \neq \underline{\gamma}' \\
\bigotimes_{i} A_{i}^{\otimes \gamma_{i}} & \text{sinon,}
\end{cases}
\]
le $\otimes$ étant l'addition sur les objets et le produit tensoriel sur les
morphismes. Son enveloppe karoubienne est une $\otimes$-catégorie abélienne,
qui s'identifie à la sous-catégorie pleine de $\mathcal{M}$ des objets à degrés
dans $\mathbf{N}^{(I)}$, et qui est munie d'un foncteur fibre sur $\bar{k}$ ;
puis, $A_{i}$ étant de rang $d_{i}^{2}$ et $E_{i}$ de rang $d_{i}$, les
$L_{i} = \det E_{i} = \bigwedge^{d_{i}} E_{i}$ sont préinversibles, et en les
inversant on retrouve $\mathcal{M}$.

Deux avertissements de sa main ferment le feuillet, et ils annoncent le
programme du dossier : « attention à contrainte de commutativité » en marge de
la construction, et « A expliciter ! » à la fin --- dans la description des
$\otimes$-foncteurs de $\mathcal{M}$, « il semble qu'il doive revenir au même
d'exiger que le rang est conservé ». C'est déjà la question du
n\textsuperscript{o} 12 des pages 129 à 131 : jusqu'où peut-on négliger la
contrainte de commutativité ?

\pagerange{61}{65}
\section*{Les caractères numériques, et leur insuffisance (pages 61 à 65)}

\subsection*{Ce qu'on peut compter}

Soit $\mathcal{M}$ une catégorie abélienne $k$-linéaire à
$\operatorname{Hom}$ de dimension finie. Si $M$ est simple, la
sous-catégorie abélienne qu'il engendre est celle des espaces vectoriels de
dimension finie sur le corps $K = \operatorname{End}(M)$. Deux conséquences
immédiates et utiles :
\begin{itemize}
\item $M$ reste semi-simple après une extension $k'/k$ si et seulement si
$K \otimes_{k} k'$ est semi-simple ;
\item cela vaut pour \emph{toute} $k'$ si et seulement si cela vaut pour une
extension parfaite, si et seulement si le centre $Z$ de $K$ est séparable sur
$k$. On dit alors que $M$ est \emph{absolument semi-simple}.
\end{itemize}
Le groupe $K(\mathcal{M})$ est libre sur l'ensemble $\Sigma(\mathcal{M})$ des
classes de simples,\footnote{L'énoncé demande que les objets soient de longueur
finie, ce que la seule finitude des $\operatorname{Hom}$ ne donne pas ; c'est
Jordan-Hölder qui fait la liberté sur les simples, et la page ne le dit pas.}
et la décomposition de $M_{k'}$ se lit sur $K \otimes_{k} k' = \prod_{i}
M_{n_{i}}(K'_{i})$ : $r$ composants isotypiques, le $i$-ème à $n_{i}$
composants simples. Pour $k'/k$ quasi-galoisienne de groupe $\pi$, les classes
obtenues sont conjuguées et
\[
\Sigma(\mathcal{M}) \simeq \Sigma(\mathcal{M}_{k'})/\pi .
\]

Si $\mathcal{M}$ est semi-simple, elle est connue à équivalence près par
$\Sigma$ et par l'application
\[
\sigma \longmapsto \bigl( Z_{\sigma},\ \xi_{\sigma} \in
\operatorname{Br}(Z_{\sigma}) \bigr) ,
\]
$Z_{\sigma}$ étant le centre de $\operatorname{End}(M_{\sigma})$ et
$\xi_{\sigma}$ la classe de celle-ci. La variance en $k'$ se calcule : la fibre
de $\Sigma' \to \Sigma$ au-dessus de $\sigma$ est l'ensemble des corps
composants de $Z_{\sigma} \otimes_{k} k'$, et l'application
$\mathbf{Z}^{(\Sigma)} \to \mathbf{Z}^{(\Sigma')}$ envoie $\sigma$ sur
$\sum_{i} n_{i}(\sigma, i)$ avec
\[
n_{i} = d/d_{i} ,
\]
$d$ étant l'ordre de $\xi_{\sigma}$ et $d_{i}$ celui de son image dans
$\operatorname{Br}(Z'^{i}_{\sigma})$.

Les quatre \emph{caractères numériques} d'une $\otimes$-catégorie sont alors
(pages 62 et 63) : l'ensemble $\Sigma$ ; la multiplication de
$K(\mathcal{M}) = \mathbf{Z}^{(\Sigma)}$, c'est-à-dire les constantes de
structure $c^{\rho}_{\sigma,\tau} \in \mathbf{N}$ de
\[
\sigma\tau = \sum_{\rho} c^{\rho}_{\sigma,\tau}\,\rho\, ;
\]
l'application $\sigma \mapsto (Z_{\sigma}, \xi_{\sigma})$ ; et la fonction rang
$\mathcal{M} \to \mathbf{Z}$.\footnote{La dernière ligne de la page 62 est de
sa main et y écrit $\sigma^{\rho}_{\sigma,\tau}$ là où la formule tapée deux
lignes plus haut porte $c^{\rho}_{\sigma,\tau}$. Les deux notations sont
données telles quelles dans la transcription ; cette lecture n'en retient
qu'une.} Les deux premières données équivalent à celle de l'anneau
$K = K(\mathcal{M})$ avec sa partie effective $K^{\mathrm{eff}} =
\mathbf{N}^{(\Sigma)}$ --- et $\Sigma$ se récupère comme l'ensemble des
éléments non nuls minimaux de $K^{\mathrm{eff}}$, ce qui est la remarque qui
rend la liste économique. La quatrième permet de calculer les polynômes
caractéristiques d'un endomorphisme en termes de normes et traces réduites dans
les algèbres d'Azumaya $\operatorname{End}(M_{i})$ : pour $M$ isotypique de
rang $n$, de centre $Z$ de rang $r$, avec $n = n'r$ et $K$ de rang $d^{2}$ sur
$Z$,
\[
\operatorname{det}_{M} f = \operatorname{N}_{Z/k}\bigl( \operatorname{Nr}_{K/Z}(f) \bigr)^{n'/d} ,
\qquad
\operatorname{Tr}_{M} f = \frac{n'}{d}\,\operatorname{Tr}_{Z/k}\bigl( \operatorname{Trr}_{K/Z}(f) \bigr) .
\]

\subsection*{Et ce que cela ne détermine pas}

Le résultat de ce paragraphe est négatif, et il est énoncé avec soin :
\begin{quote}
la connaissance des caractères numériques d'une catégorie tannakienne
$\mathcal{M}$ \emph{ne permet pas en général} de la reconstituer à équivalence
près --- même si $k$ est algébriquement clos de caractéristique $0$, de sorte
que la troisième donnée devient vide, et même si $\mathcal{M}$ est
semi-simple ; et même dans le cas particulier où $\mathcal{M}$ est la catégorie
des représentations d'un groupe fini.
\end{quote}
La parenthèse qui suit est une consigne qu'il se donne : « à vérifier dans
littérature qu'un groupe fini n'est pas déterminé par la connaissance de la
table de multiplication de son ensemble de caractères, et les rangs de ceux-ci
\ldots ». La réponse est oui, et elle est classique : le groupe diédral d'ordre
$8$ et le groupe des quaternions d'ordre $8$ ont la même table de
caractères.\footnote{Cette lecture fournit le contre-exemple ; le dossier pose
la question et ne la tranche pas.}

L'exception qu'il signale est celle de la page 59 : si le lien est
diagonalisable, $\Sigma$ devient un sous-groupe de $K^{*}$ qui détermine le
lien, et la troisième donnée fournit l'homomorphisme $\Sigma \to
\operatorname{Br}(k)$ qui suffit à tout déterminer. Reste le cas d'un lien de
type multiplicatif, et c'est ce que fait la partie c).

\subsection*{Lien de type multiplicatif}

Par changement de base à une clôture séparable $k'$, l'ensemble $\Sigma' =
\Sigma(\mathcal{M}_{k'})$ --- déduit canoniquement de la donnée numérique, avec
l'action de $\pi = \operatorname{Gal}(k'/k)$ --- reçoit de la multiplication de
$K(\mathcal{M}_{k'})$ une structure de groupe respectée par $\pi$. Cela donne
le lien $G = D(\Lambda)$, $\Lambda$ étant $\Sigma'$ avec son action. $\Sigma$
est l'ensemble des orbites, $Z_{\sigma}$ le corps des invariants du
stabilisateur d'un point de l'orbite $\sigma$, et pour $M_{\sigma}$ de type
$\sigma$ l'homomorphisme $G \to
\underline{\operatorname{Aut}}(M_{\sigma})$ se factorise par le centre en
\[
u_{\sigma} \colon G \longrightarrow \underline{Z}^{*}_{\sigma}
\simeq \prod\nolimits_{Z_{\sigma}/k} \mathbf{G}_{m} ,
\qquad u_{\sigma}(\xi) = \xi_{\sigma} ,
\]
où $\xi \in H^{2}(k_{\mathrm{pl}}, G)$ est la classe de la gerbe tannakienne.
La connaissance du caractère numérique équivaut donc à celle de $G$ et des
classes $u_{\sigma}(\xi)$ pour tous les homomorphismes de $G$ dans des tores
$\prod_{Z/k}\mathbf{G}_{m}$ : la réciproque tient parce que le rang se calcule,
\[
n(\sigma) = \operatorname{card}(\sigma) \cdot \operatorname{ord}(\xi_{\sigma}) .
\]
La question devient : $\xi$ est-il connu quand on connaît ses images
$u_{\sigma}(\xi)$ ? Oui pour $G$ diagonalisable. Oui encore pour un tore de
dimension $1$, et c'est l'\emph{Exemple 1} : un tel tore non trivial est tordu
par une extension quadratique $Z$, la suite exacte
\[
0 \longrightarrow G \xrightarrow{\ u\ } \underline{Z}^{*}
\xrightarrow{\ N_{Z/k}\ } \mathbf{G}_{m} \longrightarrow 0
\]
donne que le noyau de $H^{2}(k,G) \to H^{2}(k, \underline{Z}^{*})$ est un
quotient de $H^{1}(k, \mathbf{G}_{m})$, nul par Hilbert 90. Le caractère
numérique détermine donc $\mathcal{M}$ dans ce cas. L'\emph{Exemple 2},
sur $G' = \operatorname{Ker}(n \cdot \mathrm{id}_{G})$, est annulé par de
longues diagonales et le tapuscrit s'interrompt là.

\pagerange{66}{83}
\section*{La gerbe des $F$-isocristaux (pages 66 à 83)}

C'est le calcul le plus poussé du dossier, et le seul qui aille jusqu'à la
structure complète d'une gerbe. Dans tout ce paragraphe $k$ est un corps de
caractéristique $p > 0$, $W$ un $p$-anneau de corps résiduel $k$, $K$ son corps
des fractions, $\sigma$ le Frobenius. Attention : \emph{ici $F$ n'est pas un
foncteur fibre}, c'est le Frobenius d'un cristal.

\subsection*{Les objets}

Les cristaux en modules localement libres de type fini sur le site cristallin
absolu de $k$ s'explicitent, via $W$, comme les modules libres de type fini
$M$ sur $W$ munis pour chaque $n$ d'une connexion absolue plate sur
$M/p^{n}M$, ces connexions se recollant. Pour $k$ parfait, $W = W(k)$ est
l'anneau des vecteurs de Witt et les connexions sont une donnée vide : la
catégorie est celle des $W$-modules libres de type fini, et celle des
\emph{isocristaux} (cristaux à isogénie près) est $\operatorname{Modf}(K)$.

Un \emph{$F$-cristal} est un couple $(M, F_{M})$ avec $F_{M} \colon M^{\sigma}
\to M$. Ces objets forment une catégorie $\mathbf{Z}_{p}$-linéaire dont les
$\operatorname{Hom}$ sont des $\mathbf{Z}_{p}$-modules de type fini --- et non
des $W$-modules, même pour $k$ parfait, ce qu'il souligne. Elle n'est pas
tannakienne, même tensorisée par $\mathbf{Q}_{p}$ : il manque un
$\operatorname{Hom}$ interne. Le remède est en deux temps. $M$ est
\emph{non dégénéré} s'il existe $i \geqslant 0$ et $V \colon M \to M^{\sigma}$
avec $VF = p^{i}$ ; localisant en $p$ la catégorie des $F$-cristaux non
dégénérés, on obtient les \emph{$F$-isocristaux effectifs}. L'objet de Tate
inverse $K(-1)$ est le cristal unité avec $F = p \cdot \mathrm{id}$ ; il est
faiblement inversible, et en lui adjoignant formellement un inverse on trouve
\emph{une catégorie tannakienne sur $\mathbf{Q}_{p}$} munie d'un foncteur fibre
évident sur $K$ :
\[
\mathrm{FCriso}(k) .
\]
Pour $k$ parfait elle s'identifie à la catégorie des $(E, F_{E})$ avec $E$ de
dimension finie sur $K$ et $F_{E}$ un automorphisme $\sigma$-linéaire ; et
$(E,F_{E})$ est effectif si et seulement si $F_{E}$ laisse invariant un réseau,
c'est-à-dire si et seulement si la famille des itérés $F_{E}^{n}$, $n
\geqslant 0$, est \emph{bornée} dans $\operatorname{End}(E)$.

La gerbe associée sur $\mathbf{Q}_{p}$ est notée $\underline{G}(k)$, et le
programme est de la déterminer avec ses structures supplémentaires : la
structure d'effectivité, l'objet de Tate inverse (qui donne un homomorphisme
canonique vers la gerbe triviale de groupe $\mathbf{G}_{m}$), et pour $k$
parfait la section canonique sur $K$.

\subsection*{Pourquoi c'est intéressant : la cohomologie cristalline}

La raison de s'y intéresser est un foncteur $X \mapsto H^{*}_{\mathrm{cris}}(X)$
des schémas propres et lisses sur $k$ vers les $F$-cristaux gradués, additif,
compatible à $\otimes$ et à toutes les associativités et commutativités ---
\emph{à condition d'utiliser la règle de Koszul} pour la contrainte de
commutativité. Modulo des vérifications qu'il attribue à Berthelot (« quand
Berthelot aura bien travaillé »), il arrive même dans les $F$-cristaux non
dégénérés, avec $H^{2}_{\mathrm{cris}}(\mathbf{P}^{1}) = W(-1)$ et
$H^{2d}_{\mathrm{cris}}(X) \simeq W(-d)$ pour $X$ géométriquement connexe de
dimension $d$. Admettant le yoga des motifs --- et il faut ici la structure
graduée de la catégorie des motifs, pour pouvoir « canuler par la règle de
Koko » la commutativité géométrique --- on obtient un homomorphisme de gerbes
\[
\underline{G}(k) \longrightarrow \Pi(k)_{\mathbf{Q}_{p}} ,
\]
$\Pi(k)$ étant la gerbe des motifs, compatible à l'augmentation de Tate et aux
structures d'effectivité. Et pour $k = \mathbf{F}_{q}$ la connaissance de
l'image de $H^{*}(X)$ dans $\mathrm{FCriso}(k)$ permet de calculer la fonction
$\zeta_{X}$ « par la formule qu'on devine et que je me dispense d'écrire ».

\subsection*{Pente nulle : les cristaux unités}

\textbf{Proposition 5.1 (page 70).} Le foncteur $V \mapsto V
\otimes_{\mathbf{Z}_{p}} W$ est une équivalence de la catégorie des faisceaux
$p$-adiques constructibles sur $k$ --- c'est-à-dire des
$\mathbf{Z}_{p}$-modules de type fini munis d'une opération continue de
$\pi = \operatorname{Gal}(\bar{k}/k)$ --- sur la sous-catégorie pleine des
$F$-cristaux $M$ tels que $F_{M} \colon M^{\sigma} \to M$ soit un
\emph{isomorphisme}.

Le corollaire donne, pour $k$ parfait, quatre conditions équivalentes sur
$(E,F_{E})$ : être dans l'image essentielle ; admettre un réseau stable par
$F_{E}$ \emph{et} par $F_{E}^{-1}$ ; avoir la famille des $F^{n}$, $n \in
\mathbf{Z}$, bornée ; être effectif ainsi que son dual. Ces objets sont dits
\emph{de pente nulle}, et l'on obtient un \emph{épimorphisme de gerbes}
\[
\underline{G}(k) \longrightarrow \Pi_{p}(k) ,
\]
où $\Pi_{p}(k)$ est la gerbe des $\mathbf{Q}_{p}$-faisceaux constructibles,
équivalente --- une clôture $\bar{k}$ choisie --- à la gerbe triviale de
l'enveloppe $\mathbf{Q}_{p}$-algébrique de $\pi$. Tenant compte de
l'augmentation de Tate, c'est même un épimorphisme
\[
\underline{G}(k) \longrightarrow \Pi_{p}(k) \times
[\mathbf{G}_{m}]_{\mathbf{Q}_{p}} ,
\]
les objets de la forme $\sum_{i} V_{i} \otimes \mathbf{Q}_{p}(i)$ étant les
$F$-isocristaux dits \emph{de niveau zéro}.\footnote{L'équivalence de la
proposition 5.1 est celle qu'on appelle aujourd'hui l'équivalence des cristaux
à racine unité : la catégorie des $F$-cristaux unités sur $k$ est celle des
représentations $p$-adiques continues de $\pi_{1}(k)$. Le dossier l'établit
directement.}

\subsection*{Changement de corps, et la question 6.2}

Un homomorphisme $k \to k'$ donne $\mathrm{FCriso}(k) \to \mathrm{FCriso}(k')$
et, en sens inverse, $\underline{G}(k') \to \underline{G}(k)$, compatible aux
épimorphismes ci-dessus. La \textbf{Proposition 6.1} dit que pour $k'/k$ finie
étale c'est un \emph{monomorphisme} d'image la sous-gerbe « ouverte » définie
par l'image de $\pi_{1}(k')$ ; en particulier, pour $k'/k$ finie galoisienne de
groupe $g$,
\[
0 \longrightarrow \underline{G}(k') \longrightarrow \underline{G}(k)
\longrightarrow [g]_{\mathbf{Q}_{p}} \longrightarrow 0 .
\]
C'est, dit-il, formellement équivalent à l'effectivité de la descente de
$\operatorname{Spec}(k') \to \operatorname{Spec}(k)$ pour la catégorie fibrée
des $\mathrm{FCriso}$, dont la démonstration est « $\pm$ triviale ».

La \textbf{Question 6.2} est la même pour $k'/k$ galoisienne possiblement
infinie, et le cas intéressant est $k' = \bar{k}$ :
\[
0 \longrightarrow \underline{G}(\bar{k}) \longrightarrow \underline{G}(k)
\longrightarrow \Pi_{p}(k) \longrightarrow 0 \, ?
\]
Le n\textsuperscript{o} 7 annonce que la réponse est affirmative pour $k$ fini,
d'où, pour tout $k$,
\[
\underline{G}(k) \simeq \underline{G}(\mathbf{F}_{p})
\times_{\Pi_{p}(\mathbf{F}_{p})} \Pi_{p}(k) ,
\]
et pour $k$ parfait une réponse affirmative revient à dire que cet
homomorphisme est une équivalence. Cela ramènerait toute la structure de
$\underline{G}(k)$ à celle de $\underline{G}(\mathbf{F}_{p})$, plus la
connaissance de $\pi_{1}(k) \to \pi_{1}(\mathbf{F}_{p}) = \hat{\mathbf{Z}}$ ---
et, pour la section canonique sur $K$, plus l'opération de $\pi_{1}(k)$ sur la
clôture non ramifiée de $K$. Quant à l'effectivité, elle est invariante par
changement de base, donc se ramène au cas algébriquement clos : les
$F$-isocristaux effectifs seront ceux \emph{à pentes} $\geqslant 0$.

\subsection*{La gerbe sur le corps premier}

C'est le cœur du calcul. Le foncteur fibre canonique identifie
$\underline{G}(\mathbf{F}_{p})$ à la gerbe triviale de son lien ; et comme
$\mathrm{FCriso}(\mathbf{F}_{p})$ avec son foncteur fibre est la catégorie des
automorphismes de $\mathbf{Q}_{p}$-espaces de dimension finie, ce lien est
l'\emph{enveloppe $\mathbf{Q}_{p}$-proalgébrique du groupe discret
$\mathbf{Z}$}. Appelons-le $G$. Il est commutatif, donc produit d'un groupe
vectoriel $G_{u}$ par un groupe de type multiplicatif $G_{r}$, et $G_{r}$ est
l'enveloppe multiplicative de $\mathbf{Z}$, c'est-à-dire le groupe de type
multiplicatif de groupe des caractères
\[
X(G_{r}) = \bar{\mathbf{Q}}_{p}^{*} .
\]

De l'autre côté, $\Pi_{p}(\mathbf{F}_{p})$ est la gerbe triviale de l'enveloppe
$\mathbf{Q}_{p}$-algébrique de $\hat{\mathbf{Z}}$. Un faisceau constructible
sur $\mathbf{F}_{p}$ est un $(E, F_{E})$ dont l'opération $n \mapsto F_{E}^{n}$
de $\mathbf{Z}$ se prolonge continûment à $\hat{\mathbf{Z}}$ ; cela ne dépend
que de la composante semi-simple de $F_{E}$, et cela signifie exactement que
les valeurs propres de $F_{E}$ sont des \emph{unités} $p$-adiques. Le lien de
$\Pi_{p}(\mathbf{F}_{p})$ est donc le quotient $G_{u} \times G'_{r}$ de $G$ où
$G'_{r}$ correspond au sous-groupe des unités dans
$\bar{\mathbf{Q}}_{p}^{*}$. D'où une structure d'extension
\[
0 \longrightarrow H \longrightarrow G \longrightarrow G' \longrightarrow 0 ,
\]
où $H$ est de type multiplicatif de groupe des caractères
$\bar{\mathbf{Q}}_{p}^{*}/\text{unités}$, canoniquement $\mathbf{Q}$ par la
valuation $p$-adique normalisée par $v_{p}(p) = 1$. L'action de Galois sur ce
quotient est triviale, donc
\[
\boxed{\ H = D(\mathbf{Q}) \ }
\]
est \emph{diagonalisable de groupe des caractères $\mathbf{Q}$}.\footnote{C'est
ce qu'on appelle aujourd'hui le \emph{tore des pentes} : ses représentations
sont les $\mathbf{Q}$-graduations, et c'est par là que la décomposition
isopentique d'un $F$-isocristal est un fait de théorie des représentations et
non un théorème de structure. Le dossier le construit, et n'a pas le nom.}
L'augmentation de Tate de $G$ correspond à $\mathbf{Z} \to
\bar{\mathbf{Q}}_{p}^{*}$, $1 \mapsto p$, et celle induite sur $H$ à
l'inclusion canonique $\mathbf{Z} \hookrightarrow \mathbf{Q}$ --- grâce au
choix de la valuation normalisée.

Le dévissage se résume ainsi :

\begin{tikzcd}[column sep=small, row sep=small, nodes={font=\scriptsize}]
& \underline{G}(\bar{\mathbf{F}}_{p}) \arrow[d, "\simeq"] & \\
0 \arrow[r] & \underline{H} \arrow[r] & \underline{G}(\mathbf{F}_{p}) \arrow[r, "i"] & \Pi_{p}(\mathbf{F}_{p}) \arrow[r] & 0 \\
& D(\mathbf{Q}) \arrow[u, no head, "\text{lien}" description] & G \arrow[u, no head, "\text{lien}" description] & G' \arrow[u, no head, "\text{lien}" description] &
\end{tikzcd}

\subsection*{La classe de la gerbe noyau}

Soit $\underline{H}$ la gerbe noyau de $i$, de lien $H = D(\mathbf{Q})$. Sa
classe est
\[
\xi \in H^{2}(\mathbf{Q}_{p}, H)
= \operatorname{Hom}(\mathbf{Q}, \operatorname{Br}(\mathbf{Q}_{p}))
= \operatorname{Hom}(\mathbf{Q}, \mathbf{Q}/\mathbf{Z}) ,
\]
et vaut $\xi = \partial(\eta)$ où $\eta \in H^{1}(\mathbf{Q}_{p}, G')$ est la
classe du torseur $P$ qui définit $i$. Le calcul de la page 77 identifie $P$ :
il y a un homomorphisme canonique $\pi_{1}(\mathbf{F}_{p}) = \hat{\mathbf{Z}}
\to G'(\mathbf{Q}_{p})$, donc un foncteur des torseurs sous
$\hat{\mathbf{Z}}$ vers les torseurs sous $G'$, et $P$ est l'image du torseur
évident $P_{o}$ correspondant à \emph{l'extension non ramifiée maximale de
$\mathbf{Q}_{p}$}. Le diagramme d'extensions
\[
\begin{array}{ccccccccc}
0 &\to& H &\to& G &\to& G' &\to& 0 \\
 && \big\| && \downarrow && \downarrow && \\
0 &\to& H &\to& H' &\to& \hat{\mathbf{Z}}_{\mathbf{Q}_{p}} &\to& 0
\end{array}
\]
où $H'$ est le quotient de $G_{r}$ correspondant au sous-groupe de
$\bar{\mathbf{Q}}_{p}^{*}$ engendré par les racines de l'unité et par toutes
les racines $p^{1/n}$ de $p$, ramène le calcul à $\eta_{o}$, la classe de
$P_{o}$. Et l'on vérifie que $\xi$ correspond, par l'isomorphisme
$\operatorname{Br}(\mathbf{Q}_{p}) \simeq \mathbf{Q}/\mathbf{Z}$ du corps de
classes local, à l'homomorphisme canonique de passage au quotient
\[
\mathbf{Q} \longrightarrow \mathbf{Q}/\mathbf{Z}
\qquad
\text{« sauf erreur de signe, car je n'ai nulle part fait attention aux signes ».}
\]

\subsection*{Dieudonné et Manin}

Que $\underline{G}(\bar{\mathbf{F}}_{p}) \to \underline{G}(\mathbf{F}_{p})$ se
factorise par $\underline{H}$ est clair. Par Dieudonné et Manin, les objets de
$\mathrm{FCriso}(\bar{k})$ sont semi-simples et proviennent de
$\mathrm{FCriso}(\mathbf{F}_{p})$, donc l'homomorphisme est un monomorphisme ;
et pour qu'il soit une équivalence il suffit qu'un objet simple non constant de
$\operatorname{Rep}(\underline{H})$ donne un objet non constant sur
$\bar{\mathbf{F}}_{p}$, ce que la classification donne. Ainsi
$\underline{G}(\bar{k}) \simeq \underline{H}$ pour tout corps
algébriquement clos $\bar{k}$ de caractéristique $p$, indépendamment du choix
de la clôture.

Et $\operatorname{Rep}(\underline{H})$ se catalogue. Les classes de simples
sont indexées par $\mathbf{Q}$ --- l'élément correspondant s'appelle la
\emph{pente} --- et l'objet $\mathbf{Q}_{p}[T]/(T^{r} - p^{s})$ muni de la
multiplication par $T$, avec $r \geqslant 1$ et $(r,s) = 1$, est simple de
pente $s/r$. On obtient donc le catalogue complet
\[
E_{r,s} , \qquad \operatorname{rang} E_{r,s} = r ,
\qquad \operatorname{End}(E_{r,s}) = K_{r,s}
\ \text{de classe } s/r \ \text{dans } \operatorname{Br}(\mathbf{Q}_{p}) .
\]
$\operatorname{Rep}(\underline{H})$ est la limite inductive des
$\operatorname{Rep}(\underline{H}, r)$, où $(H,r)$ est la gerbe de lien
$\mathbf{G}_{m}$ de classe le générateur canonique $1/r$ de la $r$-torsion de
$\operatorname{Br}(\mathbf{Q}_{p})$, c'est-à-dire « la » gerbe définie par le
corps gauche $K_{r,1}$ ; ce sont les $F$-isocristaux dont les pentes sont
multiples de $1/r$. Deux exercices sont posés et non faits : expliciter la
$\otimes$-structure en termes des $E_{r,s}$, et décrire canoniquement les
modules de Dieudonné des réduits des groupes de Lubin-Tate.\footnote{La marge
de la page 79 propose la notation $E_{\lambda}$ pour $E_{r,s}$ avec $\lambda =
s/r$, et écrit $E_{\lambda} \otimes E_{\lambda'} \simeq
E_{\lambda\lambda'}^{n(\lambda,\lambda')}$. L'indice du membre de droite est
$\lambda\lambda'$ sur la page ; c'est $\lambda + \lambda'$ que la structure
demande, la pente étant additive pour $\otimes$. La transcription le donne tel
quel et le signale ; cette lecture rétablit la somme.}

\subsection*{La pente, pour $k$ quelconque}

Pour $k$ quelconque, un choix de clôture donne
$\mathrm{FCriso}(k) \to \operatorname{Rep}(\underline{H})$ et
$\underline{H} \to \underline{G}(k)$, et un $F$-isocristal est dit de pente
$\lambda$ si son image l'est. On voit aussitôt : pente nulle équivaut à
« trivial » au sens du n\textsuperscript{o} 5, et l'effectivité équivaut aux
pentes $\geqslant 0$. Et pour $k$ \emph{parfait} l'homomorphisme de liens
$H \to G(k)$ est \emph{central}, donc $H$ opère sur le foncteur identique :
tout $F$-isocristal se décompose en composantes isopentiques,
\[
E = \sum_{\lambda \in \mathbf{Q}} E_{\lambda} ,
\]
et $E$ est trivial si $E_{\lambda} = 0$ pour $\lambda \neq 0$, effectif si
$E_{\lambda} = 0$ pour $\lambda < 0$. Le cristal de Tate $K(1)$ est de pente
$-1$, d'où l'écriture unique
\[
E = \sum_{i \in \mathbf{Z}} E_{i}(-i) ,
\qquad E_{i} \ \text{à pentes dans } [0,1) .
\]

\subsection*{Ce que les pentes dans $[0,1]$ veulent dire}

Les trois premières pages du lot suivant achèvent le tapuscrit. Les
$F$-isocristaux à pentes dans $(0,1]$ sont exactement les modules de Dieudonné
des groupes formels $p$-divisibles sur $k$, la catégorie des modules étant
\emph{opposée} à celle des groupes modulo isogénie ; et l'amplitude dans
l'intervalle \emph{fermé} $[0,1]$ caractérise les modules de Dieudonné des
groupes de Barsotti-Tate, le $F$-isocristal unité répondant au groupe ind-étale
$\mathbf{Q}_{p}/\mathbf{Z}_{p}$. Pour $k$ parfait, l'existence d'un $V_{i}$
effectif avec $FV_{i} = V_{i}F = p^{i}\,\mathrm{id}$ signifie que l'amplitude
est contenue dans $[0,i]$.

Pour $\mathbf{F}_{q}$, $q = p^{a}$, le foncteur fibre canonique sur
$K_{a} = K(\mathbf{F}_{q})$ porte une structure de plus : l'automorphisme
\emph{linéaire} $F_{E}^{a}$. Le sous-groupe $G_{a} \subset G_{K_{a}}$
correspondant, ouvert d'indice $a$ --- il répond aux racines $a$-èmes de
l'unité dans le groupe des caractères --- est muni du Frobenius relatif
$f_{q} = f^{a}$. Et comme $a \cdot \mathrm{Id}_{G}$ est un monomorphisme
($x \mapsto x^{a}$ étant surjectif dans $\bar{\mathbf{Q}}_{p}^{*}$), il induit
un isomorphisme $G_{K_{a}} \xrightarrow{\sim} G_{a} \simeq
G(\mathbf{F}_{q})$ transformant le Frobenius absolu en le relatif. D'où :

\textbf{Proposition 10.1.} $\mathrm{FCriso}(\mathbf{F}_{q})^{\mathrm{tann}}
\otimes_{\mathbf{Q}_{p}} K_{a}$ est équivalente à
$\mathrm{FCriso}(\mathbf{F}_{p})^{\mathrm{tann}} \otimes_{\mathbf{Q}_{p}}
K_{a}$, avec $f_{q}$ correspondant à $f^{a}$ et correspondance des foncteurs
fibres canoniques. Sous cette équivalence, le foncteur de changement de base
$\mathbf{F}_{p} \to \mathbf{F}_{q}$ est $(E, F_{E}) \mapsto (E, F_{E}^{a})$.

Lisant l'homomorphisme de liens $\underline{H} \xrightarrow{u_{a}} G_{a} \to G$
sur les groupes de caractères, on trouve
\[
u_{a}^{*}(\lambda) = \frac{v_{p}(\lambda)}{a} ,
\]
et le théorème de Manin en sort :

\textbf{Proposition 10.3.} Soit $(E, F_{E})$ un $F$-isocristal sur
$\mathbf{F}_{q}$, $q = p^{a}$, et soient $\lambda_{1}, \ldots, \lambda_{N}$
($N = \operatorname{rang} E$) les racines du polynôme caractéristique de
$F_{E}^{a}$. Alors les pentes de $(E, F_{E})$ sont les $v_{p}(\lambda_{i})/a$,
et si les pentes distinctes sont $s_{i}/r_{i}$ avec $(s_{i}, r_{i}) = 1$,
$r_{i} \geqslant 1$, de multiplicité $d_{i}$, alors $r_{i} \mid d_{i}$, et
l'image de $(E, F_{E})$ dans $\mathrm{FCriso}(\bar{\mathbf{F}}_{p})$ est
$\sum_{i} E_{s_{i}/r_{i}}^{d_{i}/r_{i}}$.

C'est le résultat de Manin « énoncé sous une forme plus jolie et débarassé
d'hypothèses restrictives parasites », et sa marge propose la simplification
qui achève le nettoyage : prendre la valuation $v_{q}$ normalisée par
$v_{q}(q) = 1$, et ne diviser par rien. Le tapuscrit s'arrête sur un
« Cor. » de sa main que rien ne suit.

\pagerange{85}{86}
\section*{Deux feuillets étrangers (pages 85 et 86)}

Un feuillet d'un tout autre tapuscrit définit le \emph{topos classifiant d'un
pro-groupe}. Pour $\underline{G} = (G_{i})_{i \in I}$ un système projectif
strict de groupes --- transitions surjectives --- une opération de
$\underline{G}$ sur un ensemble $E$ est la donnée d'une famille $(E_{i})$ de
parties de réunion $E$, d'une opération de chaque $G_{i}$ sur $E_{i}$, et de la
condition que pour $j \geqslant i$, $E_{i}$ soit l'ensemble des points de
$E_{j}$ fixes sous le noyau de $G_{j} \to G_{i}$. Ces ensembles forment un
topos $B_{\underline{G}}$, par le critère de Giraud. Dans le cas profini, une
opération de $\underline{G}$ est une opération continue de $G = \varprojlim
G_{i}$, c'est-à-dire à stabilisateurs ouverts. Et $B_{\underline{G}}$ n'est de
la forme $\hat{C}$ que si le système est essentiellement constant --- donc,
pour un groupe profini, que s'il est fini.

Au dos, de sa main seule, le recensement des types simples : $A_{n}$ avec sa
forme adjointe et ses formes intermédiaires $A_{n,d}$ pour $d \mid n$ ;
$B_{n}$, $C_{n}$ ; $D_{n}$ avec $D_{n}'$ et $D_{n}''$ pour $n$ pair
$\geqslant 6$ ; $E_{6}$, $E_{7}$, $E_{8}$, $F_{4}$, $G_{2}$ ; et en regard les
formes portant un tilde. Le feuillet se clôt sur « 21 cas à examiner » et deux
alternatives dont seule la seconde se lit, « un \ldots des Chevalley ». Rien
n'indique de quoi ces vingt-et-un cas sont les cas.

\pagerange{88}{98}
\section*{Filtrations d'un foncteur fibre (pages 88 à 98)}

Onze feuillets de sa main, sous sa chemise « Filtration des foncteurs fibres »,
qu'il numérote lui-même. C'est la station où le dossier passe de
l'algèbre à la géométrie des groupes.

\subsection*{Ce qu'est une filtration admissible}

Soient $\mathcal{M}$ une $\otimes$-catégorie rigide sur un corps $k$, $S$ un
schéma sur $k$, et $F$ un foncteur fibre relatif de $\mathcal{M}$ vers les
faisceaux localement libres sur $S$. Une filtration de $F$ par des
sous-foncteurs $F^{(p)}$ est \emph{admissible} si
\begin{itemize}
\item[(0)] chaque $F^{(p)}(M)$ est localement facteur direct ;
\item[(i)] $F^{(p)}(M \oplus N) = F^{(p)}M \oplus F^{(p)}N$ ;
\item[(ii)] $F^{(p)}(M \otimes N) = \sum_{q+r=p} F^{(q)}(M) \otimes F^{(r)}(N)$ ;
\item[(iii)] $F^{(p)}(\check{M})$ est l'orthogonal de $F^{(-p)}(M)$ ;
\item[(iv)] pour chaque $M$ la filtration est localement exhaustive et
séparée.
\end{itemize}
Le foncteur $M \rightsquigarrow \operatorname{gr}F(M)$ est alors à valeurs dans
les modules localement libres par (0), additif par (i), multiplicatif par (ii),
et commute à la dualité par (iii). Les conditions (iii) et (iv) forcent
$F^{(0)}(\mathbf{1}) = \mathbf{1}$ et $F^{(1)}(\mathbf{1}) = 0$, donc
$\operatorname{gr}F(\mathbf{1}) = \mathbf{1}$ concentré en degré $0$. Si
$\mathcal{M}$ est semi-simple, $\operatorname{gr}F$ est exact --- sinon il faut
en faire une condition de plus, et il le dit.

\subsection*{Le torseur, et ce qui reste à démontrer}

Posant $G = \underline{\operatorname{Aut}}_{S}(F)$, un schéma en groupes affine
sur $S$, le foncteur $\operatorname{gr}F$ est le foncteur fibre tordu
$F^{P}$ par le $G$-torseur
\[
P = \underline{\operatorname{Isom}}_{S}(F, \operatorname{gr}F) ,
\]
et la graduation de $\operatorname{gr}F$ est un homomorphisme
$i \colon \mathbf{G}_{m,S} \to G^{P} = G'$. Parmi les isomorphismes de $F$ sur
$\operatorname{gr}F$ on distingue ceux qui respectent la filtration \emph{et}
induisent l'identité sur les gradués ; ils forment un sous-foncteur $Q$ de $P$,
qui est \emph{a priori} un pseudo-torseur sous un sous-groupe $H'$ de $G'$.
\begin{quote}
« Cependant, il faut démontrer que loc.\ $Q$ admet une section. »
\end{quote}
Ce n'est pas fait dans le dossier. Ce qui est fait, page 104, est autre chose
et suffit dans les cas utiles : $H'$ est toujours une limite projective de
groupes unipotents, donc si $\mathcal{M}$ est à engendrement dénombrable, tout
torseur sous $H'$ est trivial.\footnote{Il faut être précis sur ce que cela
règle. La trivialité d'un torseur sous $H'$ dit que \emph{si} $Q$ est un
torseur, il a un point ; elle ne dit pas qu'il n'est pas vide. Le dictionnaire
des pages 92 et 104 reste donc conditionné, dans le dossier, à un énoncé
d'existence qu'il annonce et ne prouve pas.}

Le groupe $H'$ se détermine (page 90) en choisissant un générateur $M$ : alors
$G' = \underline{\operatorname{Aut}}_{\mathcal{O}_{S}}(F'(M))$ et $H'$ est
formé des $g'$ tels que $g'x_{p} \equiv x_{p}$ modulo $F'^{(p+1)}(M)$ pour tout
$p$. Décomposant $\mathfrak{g}'$ suivant les caractères de $\mathbf{G}_{m}$
agissant par $i$, on trouve l'énoncé qu'il encadre :
\[
\boxed{\ \mathfrak{h}' = \sum_{n \geqslant 1} \mathfrak{g}'_{n}\ }
\qquad\text{et}\qquad
\mathfrak{b}' = \sum_{n \geqslant 0} \mathfrak{g}'_{n} ,
\]
$\mathfrak{h}'$ nilpotente donc unipotente, $\mathfrak{b}'$ \emph{parabolique},
et le centralisateur de $i$ un sous-groupe de Levi.

\subsection*{Le dictionnaire, et les homomorphismes complémentaires}

\textbf{Résumé (page 92).} Pour $k$ convenable et $\mathcal{M}$ semi-simple, la
donnée d'un foncteur fibre \emph{filtré admissible} sur $\mathcal{M}$ à valeurs
dans $S$ équivaut à celle
\begin{enumerate}
\item[(i)] d'un foncteur fibre \emph{gradué admissible} $F'$, c'est-à-dire d'un
foncteur fibre $F'$ et d'un homomorphisme $i \colon \mathbf{G}_{m,S} \to G' =
\underline{\operatorname{Aut}}_{S}F'$ ;
\item[(ii)] d'un torseur à droite $Q'$ sous le groupe uni-parabolique de $G'$
défini par $i$,
\end{enumerate}
le foncteur étant $F = F'^{Q'}$ avec la filtration descendue.

Si $\mathcal{M}$ porte une structure intrinsèque qui impose à $i$ d'être
central, et qu'on s'est donné $i_{1}$, il convient de prendre $i_{2}$ défini
par
\[
i_{1}(\lambda)\,i_{2}(\lambda) = f(\lambda) ,
\]
et $i_{1}$, $i_{2}$ sont dits \emph{complémentaires} : leurs paraboliques sont
\emph{opposés} et leur intersection est un Levi commun. S'il y a de plus la
structure de Tate, définissant $\tau \colon G' \to \mathbf{G}_{m}$, alors
$\tau\,i_{1}(\lambda) = \lambda^{n}$ et l'entier $n$ est le degré du motif de
Tate ; dans le cas qui l'intéresse, $\tau\,i_{1}(\lambda) = \lambda$.

\subsection*{L'exemple qui commande : De Rham, Hodge, Betti}

Soient $S$ un préschéma connexe sur $\mathbf{Q}$ et $\mathcal{M}$ la catégorie
des motifs sur $S$. Le \emph{foncteur de De Rham} est admissible ; le $F'$
qu'il définit n'est autre que le \emph{foncteur de Hodge} ; et comme
$G'^{\geqslant}$ a la structure qu'on vient de dire, si $S$ est réduit à un
point le torseur $Q$ est trivial et \emph{le foncteur de De Rham filtré
splitte}. Mais --- et c'est le point --- le splitting n'est \emph{pas
canonique}, en particulier il n'est pas stable par les automorphismes du
foncteur qui respectent la filtration et induisent l'identité sur les gradués.

Pour $S = \operatorname{Spec}K$ et un plongement $K \to \mathbf{C}$, le
\emph{foncteur de Betti} $F''$ à valeurs dans $\mathbf{Q}$ s'ajoute, et l'on a
deux isomorphismes canoniques :

\begin{tikzcd}[column sep=huge]
F''_{\mathbf{C}} \arrow[r, "\text{Gaga}", "\sim"'] & F_{\mathbf{C}} \arrow[r, "\text{int.\ harmoniques}", "\sim"'] & F'_{\mathbf{C}}
\end{tikzcd}

le premier par Gaga, le second par le théorème de Hodge des intégrales
harmoniques. Le composé $F_{\mathbf{C}} \simeq \operatorname{gr}
F_{\mathbf{C}}$ se déduit de $F_{\mathbf{C}} \simeq F''_{\mathbf{C}}$, qui
traduit une structure \emph{réelle} sur $F_{\mathbf{C}}$ donnée par une
involution $\sigma$, et la bigraduation est
\[
F^{p,q}_{\mathbf{C}} = F^{(p)}_{\mathbf{C}} \cdot
\sigma\bigl( F^{(q)}_{\mathbf{C}} \bigr) .
\]
C'est la théorie de Hodge dite dans le vocabulaire des foncteurs fibres : le
splitting existe, il vient de la conjugaison complexe, et il n'est pas
algébrique.

\subsection*{Parabolique et radical unipotent}

Les trois dernières pages traduisent tout cela en théorie des groupes, sans
catégories. Soient $G$ un schéma en groupes affine sur $k$ et $i \colon
\mathbf{G}_{m,k} \to G$ --- c'est-à-dire une $\otimes$-graduation du foncteur
oubli sur $\operatorname{Rep}(G)$. Elle définit une $\otimes$-filtration, et
l'on regarde
\[
U \subset P \subset G ,
\qquad
P = \underline{\operatorname{Aut}}_{\otimes,\mathrm{filt}} ,
\qquad
U = \underline{\operatorname{Aut}}_{\otimes,\mathrm{filt},\,\mathrm{id\ sur\ Gr}} .
\]
Si $G$ admet un module fidèle $V$, alors $V$ est gradué par les
$\operatorname{Gr}^{i}(V)$, filtré par $\operatorname{Fil}^{p}(V) = \sum_{i
\geqslant p}\operatorname{Gr}^{i}(V)$, et $P$ est le stabilisateur de la
filtration dans $\underline{\operatorname{Aut}}_{k}(V)$, $U$ le noyau de
$P \to \prod_{i}\underline{\operatorname{Aut}}_{k}\operatorname{Gr}^{i}(V)$ ;
dans le cas général ce sont des intersections de tels groupes.

Écrivant $\mathfrak{g} = \sum_{\alpha \in \mathbf{Z}}\mathfrak{g}^{\alpha}$ et
$V = \sum_{\beta}V^{\beta}$ suivant les caractères de $\mathbf{G}_{m}$, on a
$[\mathfrak{g}^{\alpha}, \mathfrak{g}^{\beta}] \subset
\mathfrak{g}^{\alpha+\beta}$ et $\mathfrak{g}^{\alpha}V^{\beta} \subset
V^{\alpha+\beta}$. L'algèbre de Lie de $P_{V}$ est formée des $x$ tels que
$xV^{\beta} \subset \sum_{\beta' \geqslant \beta}V^{\beta'}$ ; il suffit que
$x \in \sum_{\alpha \geqslant 0}\mathfrak{g}^{\alpha}$, et c'est nécessaire si
$V$ est fidèle, car en écrivant $x = \sum x^{\alpha}$ et en prenant $\mu$
minimal avec $x^{\mu} \neq 0$, la condition $x^{\mu}V^{\alpha} \subset
V^{\alpha+\mu}$ impose $\mu \geqslant 0$.\footnote{Le bas de la page 97 écrit
« $\nu + \mu \leqslant \alpha$ i.e.\ $\mu \leqslant 0$ », avec les deux
inégalités dans ce sens. La condition d'invariance donnée quatre lignes plus
haut demande $\alpha + \mu \geqslant \alpha$, donc $\mu \geqslant 0$, en accord
avec le $x \in \sum_{\alpha \geqslant 0}\mathfrak{g}^{\alpha}$ de la même page.
Cette lecture rétablit le sens ; la transcription donne la leçon telle
quelle.} Par passage à la limite,
\[
\mathfrak{p} = \sum_{\alpha \geqslant 0}\mathfrak{g}^{\alpha} ,
\qquad
\mathfrak{u} = \sum_{\alpha > 0}\mathfrak{g}^{\alpha} .
\]

\textbf{Proposition (page 98).} Si $G$ est connexe, de type fini et réductif,
alors $P$ est parabolique et $U$ est son radical unipotent.

En caractéristique $0$ cela résulte aisément de la théorie ; en
caractéristique $p > 0$, « il faut sans doute travailler un peu \ldots ». Le
corollaire, en revanche, est net :
\begin{quote}
\textbf{Corollaire.} Toute $\otimes$-filtration d'un foncteur fibre sur un
corps parfait splitte, car $H^{1}(k, U) = 0$ par Rosenlicht.
\end{quote}
Le crochet de sa main ajoute « ou même sur un quelconque » ; c'est un espoir et
non un énoncé, et la restriction à un corps parfait est celle que l'annulation
invoquée demande.\footnote{L'annulation du $H^{1}$ galoisien d'un groupe
unipotent est celle d'un corps parfait ; sur un corps imparfait elle peut
tomber en défaut, de sorte que le crochet entre parenthèses n'est pas
justifié par l'argument donné.}

Une \emph{Remarque} ferme le faisceau : posant $T = \operatorname{Im}i$, les
$uTu^{-1}$ pour $u$ dans un ouvert dense de $U$ engendrent $T \cdot U$ ---
parce que $[\mathfrak{g}^{\alpha}, x] = \alpha\,\mathfrak{g}^{\alpha} =
\mathfrak{g}^{\alpha}$ pour $\alpha > 0$ --- d'où, pour toute représentation
$V$ de $G$,
\[
V^{T \cdot U} = \bigcap_{u}V^{uTu^{-1}}
\qquad\text{et}\qquad
V^{G} = V^{P} .
\]
Le dernier bloc de la page est encadré et barré d'une dizaine de diagonales ;
il ne rend que ces deux formules.

\pagerange{101}{108}
\section*{Les exemples des « Notes Saavedra » (pages 101 à 108)}

Huit feuillets tapés, corrigés de bout en bout à l'encre, annoncés par les deux
mots « Notes Saavedra » de la page 99. Ils sont adressés à Saavedra, qui y est
nommé deux fois, et ils sont le seul endroit du dossier où il le soit. Six
exemples numérotés.

\subsection*{1) Graduations : le critère de centralité}

Soit $\underline{C}_{M}$ la catégorie des espaces vectoriels de dimension finie
sur $k$ munis d'une graduation de type un groupe $M$. C'est une
$\otimes$-catégorie munie du foncteur fibre « oubli de la graduation », et le
groupe associé est le groupe de type multiplicatif $D_{k}(M)$ ; pour $M =
\mathbf{Z}^{r}$ on trouve $\mathbf{G}_{m}^{r}$.

D'où l'application. Soit $\underline{C}$ une $\otimes$-catégorie sur $k$ munie
d'un foncteur fibre $F$, de groupe $G$. Les manières de graduer $F(V)$ de type
$M$, fonctoriellement en $V$ et compatiblement à $\otimes$, correspondent aux
$\otimes$-foncteurs $\underline{C} \to \underline{C}_{M}$ compatibles aux
foncteurs fibres, donc aux homomorphismes
\[
D_{k}(M) \longrightarrow G .
\]
Et voici le critère qui sera réutilisé partout : \emph{une telle graduation de
$F$ provient d'une graduation des objets eux-mêmes si et seulement si
l'homomorphisme est central}. Pour $\underline{C}$ les motifs sur $k$ munis
d'un foncteur fibre, on trouve donc un homomorphisme central canonique
$i \colon \mathbf{G}_{m} \to G$ --- la graduation par le poids ; et le motif de
Tate donne $j \colon G \to \mathbf{G}_{m}$, avec
\[
j\,i(\lambda) = \lambda^{2}
\]
puisque $T$ est de poids $2$.\footnote{Le tapuscrit écrit « le motif de Tate
(qui est de rang $2$, et de poids $2$) » ; il est de rang $1$ et de poids $2$.
Rien sur la page ne corrige. La relation $ji(\lambda) = \lambda^{2}$, elle, ne
dépend que du poids et est correcte.}

En caractéristique nulle il y a toujours un foncteur fibre naturel, le
\emph{foncteur de Hodge}, qui associe à un motif le vectoriel bigradué
$\coprod H^{q}(X, \Omega^{p})$. Il donne $\mathbf{G}_{m}^{2} \to G$,
c'est-à-dire deux homomorphismes commutant l'un à l'autre $i_{1}$, $i_{2}$,
avec $i(\lambda) = i_{1}(\lambda)i_{2}(\lambda)$ --- et il faut faire attention
que \emph{$i_{1}$ et $i_{2}$ ne sont pas centraux}, la bigraduation de Hodge ne
provenant pas d'une bigraduation du motif.

\subsection*{2) La filtration de De Rham}

Le \emph{foncteur de De Rham}, qui donne l'hypercohomologie $H^{*}(X,
\Omega^{*})$, n'est plus bigradué : gradué et filtré seulement, la filtration
étant celle de la suite spectrale d'hypercohomologie, qui commence par la
cohomologie de Hodge et qui dégénère --- d'où
$\operatorname{Gr}(H_{DR}(V)) \simeq H_{Hdg}(V)$, fonctoriellement.

Le paragraphe pose alors le problème général : déterminer les filtrations
décroissantes, discrètes, indexées par $\mathbf{Z}$, d'un foncteur fibre $F$,
compatibles aux produits tensoriels. Et il énonce immédiatement pourquoi la
réponse ne peut pas être un homomorphisme dans $G$ :
\begin{quote}
la catégorie des espaces vectoriels de dimension finie filtrés n'est pas
abélienne --- les bimorphismes n'y sont pas des isomorphismes.
\end{quote}
C'est l'obstruction structurelle, et toute la construction des pages 103 à 105
est faite pour la contourner : on remplace le « groupe des filtrations » par un
second foncteur fibre plus un torseur. On pose $F'(V) = \operatorname{Gr}F(V)$,
à valeurs dans les gradués ; $F'$ correspond à un torseur $P$ sous $G$, de
groupe tordu $G'$ ; la graduation de $F'$ est $i_{1} \colon \mathbf{G}_{m} \to
G'$ ; et l'on regarde
\[
H'_{i_{1}} \subset G'
\qquad\text{et}\qquad
Q = \underline{\operatorname{Isom}}^{\mathrm{filt}}(F, F') \subset
P = \underline{\operatorname{Isom}}(F, F') ,
\]
$H'_{i_{1}}$ étant les $\otimes$-automorphismes de $F'$ qui respectent la
filtration et induisent l'identité sur le gradué. $Q$ est \emph{a priori} un
pseudo-torseur à droite sous $H'$ --- vide ou torseur --- et « il faudrait
prouver que c'est bien un torseur ».

\subsection*{Le dictionnaire, et le splitting non canonique}

Modulo cette vérification, la donnée d'un foncteur fibre \emph{filtré} $F$
revient à celle de :
\begin{enumerate}
\item[(i)] un foncteur fibre $F'$, de groupe $G' =
\underline{\operatorname{Aut}}(F')$ ;
\item[(ii)] une graduation $i_{1} \colon \mathbf{G}_{m} \to G'$, qui définit
$H'_{i_{1}}$ ;
\item[(iii)] un torseur à droite $Q'$ sous $H'_{i_{1}}$,
\end{enumerate}
avec $F = F' \wedge^{H'_{i_{1}}} Q'$.

Et voici le résultat, qui remplace sur la page un N.B.\ barré de dix
diagonales :
\begin{quote}
le groupe $H'$ est nécessairement une limite projective de groupes algébriques
\emph{unipotents}. Donc si $\underline{C}$ est à engendrement dénombrable ---
de façon que $G'$, donc $H'$, soit limite projective d'une \emph{suite} de
groupes algébriques --- alors tout torseur sous $H'$ est trivial, et tout
foncteur fibre filtré est en fait associé à un foncteur fibre gradué.
\end{quote}
Autrement dit la filtration admet un splitting compatible aux produits
tensoriels. Mais le choix d'un tel splitting est le choix d'un point d'un
torseur sous le groupe pro-unipotent $H$, et il n'est
\begin{quote}
« pas du tout canonique ! »
\end{quote}
C'est la même conclusion que la page 94 tire du foncteur de De Rham, et le
dossier y revient trois fois : le scindage existe, il n'est pas donné avec
l'objet.\footnote{L'argument de dénombrabilité est exactement ce qui permet de
faire l'induction sur les $H'_{n}$ : un torseur sous une limite projective
\emph{dénombrable} de groupes unipotents se trivialise étage par étage, la
surjectivité à chaque cran venant de l'annulation du $H^{1}$ unipotent. Pour
une limite non dénombrable l'argument ne se conclut pas, et la page le dit en
la supposant.}

\subsection*{3) Pré-structures de Hodge}

Une \emph{pré-structure de Hodge} sur un $\mathbf{Q}$-espace $V$ de dimension
finie est une bigraduation $V_{\mathbf{C}} = \coprod_{p,q}V^{p,q}$ telle que
\begin{itemize}
\item[a)] la graduation totale soit définie sur $\mathbf{Q}$, i.e.
$\coprod_{p+q=n}V^{p,q}$ provienne d'un sous-espace $V^{n}_{\mathbf{Q}}$ ;
\item[b)] $\overline{V^{p,q}} = V^{q,p}$.
\end{itemize}
C'est, dans le vocabulaire d'aujourd'hui, une structure de Hodge rationnelle
tout court ; ce qu'il appelle \emph{structure de Hodge} est la structure de
Hodge \emph{polarisable}.\footnote{Se donner une pré-structure de Hodge sur $V$
équivaut à se donner une représentation du tore de Deligne $S =
R_{\mathbf{C}/\mathbf{R}}\mathbf{G}_{m}$ sur $V_{\mathbf{R}}$ : les deux
homomorphismes $i_{1}$, $i_{2}$ de la page 105 sont les deux composantes de
$S_{\mathbf{C}} \simeq \mathbf{G}_{m} \times \mathbf{G}_{m}$, la condition a)
est que le poids $i = i_{1}i_{2}$ soit défini sur $\mathbf{Q}$, et la condition
b) est $i_{2} = \overline{i_{1}}$, c'est-à-dire la réalité de l'homomorphisme.
Le dossier écrit exactement ces deux conditions et ne dispose pas du tore.}

Ces objets forment une $\otimes$-catégorie sur $\mathbf{Q}$ à foncteur fibre
canonique, donc un groupe $G$, et la bigraduation donne deux homomorphismes
$i_{1}$, $i_{2}$ de $\mathbf{G}_{m,\mathbf{C}}$ dans $G_{\mathbf{C}}$ qui
commutent. La condition a) dit que $i_{\mathbf{C}} = i_{1}i_{2}$ provient d'un
$i \colon \mathbf{G}_{m} \to G$, nécessairement \emph{central} --- parce que les
composantes homogènes d'une pré-structure de Hodge en sont elles-mêmes, de
sorte que $V$ est somme directe des $V^{n}_{\mathbf{Q}}$ \emph{en tant que
pré-structure}. La condition b) dit $i_{2}(\lambda) =
\overline{i_{1}(\bar{\lambda})}$.

Vient alors la traduction de la conjecture de Hodge. Le foncteur de
Betti-Hodge, qui à une variété projective lisse sur $\mathbf{C}$ associe
$H^{*}(X, \mathbf{Q})$ avec sa pré-structure, va des motifs sur $\mathbf{C}$
vers les pré-structures ; et
\begin{quote}
la conjecture de Hodge standard équivaut à ce que ce foncteur soit
\emph{pleinement fidèle}, c'est-à-dire à ce que l'homomorphisme du groupe de
Galois de Hodge vers le groupe de Galois motivique soit un \emph{épimorphisme}.
\end{quote}
La \emph{polarisation} de poids $n$ est un accouplement $\varphi \colon V
\times V \to \mathbf{Q}(n)$ dont la forme hermitienne associée
\[
\psi(x,y) = \varphi(x, \bar{y})\,(-i)^{p-q}
\qquad\text{pour } x \text{ de bidegré } (p, n-p)
\]
est définie positive.\footnote{Le bidegré est tapé $(p,q)$ et le $n$ écrit à
l'encre par-dessus le $q$ : c'est bien $q = n - p$, et le facteur est donc
$(-i)^{2p-n}$.} La théorie de Hodge assure que le foncteur arrive en fait dans
les structures de Hodge. Sur l'image essentielle, en revanche :
\begin{quote}
« On n'a aucune idée sur ce que pourrait être l'image essentielle du foncteur
précédent \ldots\ cela semble peu probable, mais on n'a aucune indication
sérieuse dans un sens ou l'autre. »
\end{quote}
Et l'exercice qu'il laisse --- expliciter en termes de $(G, i_{1})$ la
condition pour que les objets soient polarisables --- est celui qu'il appelle
« plaisant et délectable », en notant que l'on trouve des restrictions très
sérieuses sur $G$, en plus de sa réductivité.

\subsection*{4) et 5) Le Frobenius à l'infini, et le critère de profinitude}

Le n\textsuperscript{o} 4 est ce qu'il laisse à Saavedra : la structure
supplémentaire portée par la structure de Hodge d'une variété complexe donnée
comme déduite d'une variété \emph{réelle}. On trouve un élément $f_{\infty}$ de
$G(\mathbf{Q})$ d'ordre $2$ --- « un élément de Frobenius à l'infini » --- qui
est l'automorphisme du foncteur de Betti induit par la conjugaison sur
$X_{0}(\mathbf{C})$, et il reste à expliciter ses relations avec $i_{1}$ et
$i_{2}$.

Le n\textsuperscript{o} 5, pour $k$ de caractéristique nulle, est le critère :
\begin{quote}
$G$ est profini si et seulement si tout objet $M$ de $\underline{C}$ engendre
une $\otimes$-catégorie semi-simple n'ayant qu'un nombre fini d'objets simples
non isomorphes.
\end{quote}
Et $\underline{C}_{0}$, la sous-catégorie pleine des objets ayant cette
propriété, est celle des objets sur lesquels $G$ opère à travers un quotient
\emph{fini}, c'est-à-dire celle qui correspond au pro-groupe quotient de $G$
par les composantes neutres, $\varprojlim G_{i}/G_{i}^{\circ}$.\footnote{La
page écrit « le pro-groupe quotient de $G$ formé des $G_{i}^{\circ}$ ». Prise
au mot, la phrase est fausse : les $G_{i}^{\circ}$ forment un pro-\emph{sous}
groupe, $G^{\circ} = \varprojlim G_{i}^{\circ}$, et non un quotient. C'est le
quotient \emph{par} les composantes neutres qui est en jeu, et c'est lui qui
répond aux objets décrits : un objet engendre une sous-catégorie semi-simple à
nombre fini de simples exactement quand l'image de $G$ dans son groupe linéaire
est finie.} Il ajoute qu'il serait intéressant de trouver les énoncés
correspondants en caractéristique quelconque.

\subsection*{6) La polarisation, renvoyée}

Le dernier numéro est une page de programme, et c'est elle qui explique
pourquoi les trente premiers feuillets du dossier existent : la polarisation
d'un motif de poids $n$ --- le choix, parmi les formes symétriques pour $n$
pair et alternées pour $n$ impair $M \otimes M \to T(n)$, de celles qui sont
définies positives --- est une structure supplémentaire remarquable dans la
catégorie des motifs, qui se reflète en structures supplémentaires sur les
groupes de Galois motiviques.
\begin{quote}
« Il y a lieu de faire une étude axiomatique abstraite d'une telle notion de
polarisation sur une $\otimes$-catégorie générale au dessus d'un sous-corps du
corps des réels. On pourra en rediscuter à l'occasion. »
\end{quote}
Cette étude axiomatique est faite : ce sont les pages 3 à 16.

\pagerange{110}{112}
\section*{Récupérer le groupe de sa catégorie de modules (pages 110 à 112)}

Trois feuillets de sa main, écrits vite. Soit $G$ un schéma en groupes sur une
base $S$, soumis à l'une des deux conditions : $a)$ $G$ affine, ou $b)$ $G$
quasi-compact, quasi-séparé et plat sur $S$. Soient
$\underline{C}_{qc}(G)$ la catégorie des $S$-modules quasi-cohérents où $G$
opère et $\underline{C}_{f}(G)$ sa sous-catégorie pleine des modules de type
fini. Alors
\[
\underline{C}(G) \xrightarrow{\ \approx\ }
\operatorname{Ind}\underline{C}_{f}(G) ,
\]
ce qui est le point technique : \emph{tout} module à opérateurs est réunion
filtrante de ses sous-modules de type fini stables, exactement le corollaire
11.10 du tapuscrit SGA 3 des pages 47 à 51.

Le premier calcul est une restriction de Weil. Pour $f \colon T \to S$ affine et
$\mathcal{B} = f_{*}(\mathcal{O}_{T})$, la donnée d'un
$\mathcal{O}_{T}$-module quasi-cohérent $\underline{F}$ revient à celle de
$\underline{E} = f_{*}(\underline{F})$ muni d'une opération de $\mathcal{B}$ ;
et la donnée d'une opération de $G_{T}$ sur $\underline{F}$ revient à celle
d'une opération de $G$ sur $\underline{E}$ \emph{commutant à $\mathcal{B}$},
parce que
\[
\prod_{T/S}\underline{\operatorname{Aut}}(\underline{F})
\simeq \underline{\operatorname{Aut}}_{\mathcal{B}}
\Bigl( \prod_{T/S}\underline{F} \Bigr)
= \underline{\operatorname{Aut}}_{\mathcal{B}}(\underline{E}) .
\]
Donc $\underline{C}(G_{T})$ est la catégorie des objets de
$\underline{C}(G)$ munis d'une opération de $\mathcal{B}$ --- si $S$ est affine
d'anneau $A$ et $\mathcal{B}$ défini par une $A$-algèbre $B$, la catégorie des
objets de $\underline{C}(G)$ à anneau $B$ d'endomorphismes. Tout cela est donc
déterminé par $(\underline{C}_{f}(G), \varphi_{G}^{f})$, le couple formé de la
catégorie des modules de type fini et du foncteur oubli.

Vient alors la question qui est celle du dossier tout entier, posée ici dans sa
forme la plus nue. Soit $G$ affine sur un anneau $A$ :
\begin{quote}
dans quelle mesure $G$ se récupère-t-il par la connaissance de
$(\underline{C}_{f}(G), \varphi_{G}^{f})$, la première comme
$\otimes$-catégorie et le second comme $\otimes$-foncteur ?
\end{quote}
Il suffirait de récupérer $G(A)$, et l'on a l'homomorphisme canonique
\[
G(A) \longrightarrow \operatorname{Aut}_{\otimes}(\varphi_{G}^{f}) .
\]
L'injectivité se prouve en regardant l'effet de $g \in G(A)$ sur l'algèbre
affine de $G$ elle-même. La surjectivité, dit-il, « se prouve sans doute
aussi », en regardant l'effet d'un $\otimes$-automorphisme sur cette même
algèbre --- et en marge : « à détailler ».\footnote{C'est le théorème de
reconstruction : pour $G$ affine et plat sur $A$, l'homomorphisme ci-dessus est
un isomorphisme, de sorte que $G$ est bien déterminé par sa catégorie de
représentations munie du foncteur oubli. L'argument esquissé ici --- faire agir
sur l'algèbre affine, qui est le module régulier --- est celui qui le donne.}

Et la question finale, qu'il marque de trois points d'interrogation :
\begin{quote}
« À quelle condition une $\otimes$-catégorie $A$-linéaire munie d'un
$\otimes$-foncteur $\varphi^{f} \colon \underline{C}_{f} \to
\operatorname{Mod}_{f}(A)$ provient-elle d'un schéma en groupes affine sur $A$
??? »
\end{quote}
avec l'indication de ce qu'il faudrait supposer : que $\underline{C}$ soit un
$\mathcal{A}b$-topos et que $\varphi$ commute aux limites inductives. C'est la
question dont la théorie des catégories tannakiennes est la réponse ; le dossier
la pose et n'y répond pas.

\pagerange{114}{116}
\section*{Le plan du traité (pages 114 à 116)}

Deux feuillets disent ce que le dossier voulait devenir. Le premier porte le
plan en neuf chapitres, tapé, avec l'appendice et les deux derniers alinéas
ajoutés à l'encre :

\begin{enumerate}
\item Généralités sur les opérations $\otimes$ sur une catégorie. Appendice :
rappels sur les coalgèbres, comodules, bigèbres.
\item Relations entre $\otimes$-catégories avec foncteur fibre, et coalgèbres
resp.\ hyperalgèbres : dualité de Tannaka. (Structures mixtes, connexions.)
\item Exemples nombreux et divers. Dictionnaire détaillé.
\item Filtrations sur un foncteur fibre.
\item Théorème de structure des $\otimes$-catégories en termes de gerbes.
\item Polarisation d'une $\otimes$-catégorie. Groupes de Hodge.
\item $\otimes$-Catégories à lien réductif. Classification sur $\mathbf{R}$.
\item Structures de Hodge, groupes de Mumford-Tate et espaces modulaires
correspondants.
\item Questions d'effectivité (?)
\end{enumerate}
et en marge, de sa main, trois ajouts : un chapitre éventuel sur le cas non
commutatif, un autre pour les catégories triangulées, et un sur les catégories
de Picard, plus généralement les groupoïdes à loi $\otimes$, et le cas des
catégories ayant les deux lois $\oplus$ et $\otimes$.

Ce plan est la clé du dossier : chaque station en est un chapitre. Les pages
117 à 136 sont le chapitre 1 et une partie du 2 ; les pages 59 à 83 sont le
chapitre 3 et une partie du 5 ; les pages 88 à 105 sont le chapitre 4 ; les
pages 3 à 30 sont le chapitre 6 ; la page 86 est peut-être un morceau du
chapitre 7 ; les pages 101 à 107 sont le chapitre 8 ; et les pages 66 à 83, le
chapitre 9. \emph{Aucun des neuf n'est écrit.}

La table des matières du chapitre 1 (pages 115 et 116) est annotée partout, les
numéros de paragraphes renumérotés à l'encre. Elle est instructive sur deux
points. D'abord elle nomme un objet que le dossier n'emploie presque pas : le
$\operatorname{Hom}$ interne pris de l'autre côté, que la machine écrit
$\underline{\operatorname{Moh}}(Y,X)$, avec la remarque au paragraphe 2.5 que
la commutativité les identifie, $\underline{\operatorname{Hom}}(X,Y) \simeq
\underline{\operatorname{Moh}}(Y,X)$ --- contre laquelle sa main écrit
« non ! ». Ensuite elle place les « lois $\otimes$ rigides » et le
« formalisme des traces » au paragraphe 2, donc \emph{avant} les
$\otimes$-foncteurs, et elle ajoute en marge du paragraphe 0 cinq exemples de
lois $\otimes$ : le produit tensoriel de modules, une catégorie à produits
cartésiens, les torseurs sous un faisceau abélien, les familles de modules
indexées, et l'espace des lacets d'un espace topologique.

\pagerange{117}{128}
\section*{Les axiomes, et le dictionnaire (pages 117 à 128)}

\subsection*{Les axiomes, et le fait que les contraintes sont indépendantes}

Une $\otimes$-catégorie est une catégorie additive munie d'un bifoncteur
$(X,Y) \mapsto X \otimes Y$, d'un isomorphisme d'associativité
$\varphi(X,Y,Z)$ et d'un isomorphisme de commutativité $\psi(X,Y)$, soumis à :
a) l'axiome du pentagone pour $\varphi$ ; b) $\psi(X,Y)\psi(Y,X) =
\mathrm{id}$ ; c) l'axiome de l'hexagone, qui est la compatibilité de $\varphi$
et $\psi$ ; d) l'existence d'une unité $\mathbf{1}$ avec les isomorphismes
$u(X) \colon \mathbf{1} \otimes X \simeq X$ et $v(X) \colon X \otimes \mathbf{1}
\simeq X$ vérifiant $u(\mathbf{1}) = v(\mathbf{1})$, ce qui équivaut à
$\psi(\mathbf{1},\mathbf{1}) = \mathrm{id}$ ; e) la biadditivité de $\otimes$.

Le premier résultat est une observation de méthode, et c'est elle qui organise
le chapitre :
\begin{quote}
a) est indépendant de $\psi$ ; b) est indépendant de $\varphi$ ; a), b), c),
d) sont indépendants du fait que $\underline{C}$ soit additive.
\end{quote}
D'où la \emph{dissociation} du n\textsuperscript{o} 4 : la donnée d'un
$\otimes$ se décompose en aspects indépendants --- l'associativité soumise au
pentagone ; l'unité bilatère, unique à isomorphisme unique près ; la
commutativité soumise à la réflexivité --- chacun avec sa notion de foncteur
compatible, et les compatibilités croisées à postuler séparément (il pense
qu'elles ne sont pas conséquences des autres axiomes, et le dit).
L'additivité et la biadditivité ne viennent qu'en dernier lieu. L'unité elle-même
n'est pas une donnée mais un axiome d'existence, puisque deux unités sont
reliées par un unique isomorphisme compatible, qu'il construit comme le composé
$\mathbf{1} \to \mathbf{1} \otimes \mathbf{1}' \to \mathbf{1}'$.

Deux conséquences se lisent aussitôt. Pour $u$, $v$ endomorphismes de
$\mathbf{1}$ on a $u \otimes v = uv = vu$ : sous la contrainte de
commutativité, $\operatorname{End}(\mathbf{1})$ est un monoïde commutatif, et
si $\underline{C}$ est additive c'est un \emph{anneau commutatif}
\[
k = \operatorname{End}(\mathbf{1})
\]
qui opère sur chaque objet, de sorte que $\underline{C}$ est
$k$-linéaire et $\otimes$ est $k$-bilinéaire. Et sans la commutativité, ce
monoïde opère des deux côtés et les deux opérations ne coïncident pas.

\subsection*{Le dictionnaire}

Voici le n\textsuperscript{o} 6, le cœur du chapitre. Supposons
$\underline{C}$ abélienne $k$-linéaire, $k$ un corps, munie d'un foncteur
\emph{fidèle et exact} $F \colon \underline{C} \to \operatorname{Modf}(k)$.
Alors $\underline{C}$ s'identifie à la catégorie des représentations continues
de dimension finie d'une algèbre profinie associative unitaire $U$, ou encore
des comodules sur son dual topologique $A$, qui est une coalgèbre associative
unitaire. Et :

\begin{quote}
la donnée d'un $\otimes$ bilinéaire compatible avec $F$ --- au sens fort
$F(X \otimes Y) = F(X) \otimes F(Y)$ --- \emph{équivaut} à celle d'un morphisme
diagonal $U \to U \mathbin{\widehat{\otimes}_{k}} U$, ou encore d'une loi de
multiplication sur $A$.
\end{quote}
Puis, contrainte par contrainte : une donnée supplémentaire d'associativité,
d'unité ou de commutativité compatible avec $F$ est \emph{uniquement
déterminée} si elle existe, et elle existe si et seulement si $A$ est
respectivement associative, unitaire, commutative. Quand les trois conditions
sont satisfaites, on a un \emph{schéma en monoïdes affine} sur $k$ qui permet
de reconstituer la situation --- un monoïde et non un groupe, la rigidité
n'étant pas encore supposée.

La fonctorialité suit : les foncteurs $\underline{C} \to \underline{C}'$
compatibles avec $F$, $F'$ correspondent aux homomorphismes d'algèbres
topologiques $U' \to U$, ou de coalgèbres $A \to A'$ ; et un tel foncteur est
un $\otimes$-foncteur si et seulement si $U' \to U$ est compatible aux
diagonales, c'est-à-dire $A \to A'$ aux multiplications. Enfin, en marge, ce
qui sera repris au n\textsuperscript{o} 9 : les endomorphismes de $F$
compatibles avec $c = \mathrm{id}$ sont les $g \in U$ tels que $\Delta(g) = g
\otimes g$, et si l'on veut respecter l'unité il faut ajouter
$\varepsilon(g) = 1$.

Le dictionnaire se résume à ce carré, où chaque colonne est déterminée par les
deux autres :

\begin{tikzcd}[column sep=small, row sep=small, nodes={font=\scriptsize}]
(\underline{C}, F) \arrow[r, "\text{6)}"] \arrow[d, "\otimes"'] & U = \operatorname{End}(F) \arrow[r, "\text{dual}"] \arrow[d, "\Delta"'] & A \arrow[d, "\text{mult.}"] \arrow[r] & \operatorname{Spec}A \\
(\underline{C}, \otimes, F) \arrow[r] & U \rightarrow U \mathbin{\widehat{\otimes}} U \arrow[r] & A \otimes A \rightarrow A \arrow[r] & \text{monoïde affine}
\end{tikzcd}

\subsection*{Le signe de Koszul, et les foncteurs fibres tordus}

Le n\textsuperscript{o} 7 prend l'exemple des modules gradués sur un anneau
commutatif $k$, avec le produit tensoriel gradué ordinaire : tous les axiomes
sont vérifiés, le foncteur oubli est un $\otimes$-foncteur respectant tout, et
pour $k$ un corps on trouve le groupe $\mathbf{G}_{m}$, à partir duquel
$\underline{C}$ se récupère.

Puis il change $\psi$ pour la \emph{règle de Koszul}, et il donne la raison de
le faire. Soit $\underline{V}$ la catégorie des espaces compacts avec le produit
cartésien --- où, dit-il, « pour le coup on n'a vraiment pas le choix » pour
l'associativité et la commutativité --- et considérons $X \mapsto
H^{\bullet}(X,k)$. Künneth donne l'isomorphisme de compatibilité $c$, et
\emph{$c$ n'est compatible aux contraintes qu'à condition de choisir $\psi$ par
la règle de Koszul} : la cohomologie force le signe. Suit alors l'énoncé qu'il
marque « ??? » en marge et ne démontre pas :
\begin{quote}
« A prouver que lorsqu'on prend $\underline{C}$ comme dessus à la Koko, il n'y
a pas de foncteur fibre sur $\underline{C}$ sur aucune extension $K$ de
$k$. »
\end{quote}
C'est l'obstruction qui commande le reste. Le remède du n\textsuperscript{o} 8
est de changer de cible : on étudie les $\otimes$-foncteurs à valeurs dans les
modules \emph{gradués}, qu'il appelle \emph{foncteurs fibres tordus}. Se donner
une factorisation de $F$ par $\operatorname{Gradf}(k)$ revient à un
homomorphisme de coalgèbres $A \to k[T, T^{-1}]$, ou à la donnée de projecteurs
$\pi_{i} \in U$ avec $\pi_{i}^{2} = \pi_{i}$, $\pi_{i}\pi_{j} = 0$ pour
$i \neq j$, $\sum \pi_{i} = 1$. La $\otimes$-structure impose
\[
\Delta(\pi_{i}) = \sum_{j+k=i} \pi_{j} \otimes \pi_{k} ,
\]
et la condition de symétrie se lit :
\begin{quote}
en graduant $U$ par le « degré des opérateurs dans les gradués »,
$U^{i}_{\alpha} = \widehat{\sum}\,\pi_{i+j}U_{\alpha}\pi_{i}$, l'application
diagonale devient \emph{anticommutative} --- c'est-à-dire que la
multiplication de $A$ est anticommutative.
\end{quote}
Et le n\textsuperscript{o} 9 note que le groupe des $\otimes$-automorphismes du
foncteur oubli est \emph{indépendant} des contraintes a) à c) en jeu : ce sont
les $g \in U$ avec $\Delta(g) = g \otimes g$ et $\varepsilon(g) = 1$.

Le critère qui règle la question est page 123 :
\begin{quote}
il suffit que l'homomorphisme induit $\mu_{2} \to G$ soit \emph{central}, a
fortiori que $\mathbf{G}_{m} \to G$ le soit.
\end{quote}
Cela signifie que la décomposition d'un objet en composantes paires et impaires
est invariante sous $G$, donc se fait dans $\underline{C}$ elle-même. En
caractéristique $2$ « Koko ne change rien, tous nos ennuis et distinguos
disparaissent ».\footnote{C'est, à la lettre, la structure qu'on appelle
aujourd'hui super-tannakienne : la catégorie est celle des représentations d'un
schéma en groupes muni d'un élément central $\epsilon$ de carré $1$, qui décide
de la parité, et le foncteur fibre est à valeurs dans les super-espaces
vectoriels. Le critère « $\mu_{2} \to G$ central » est exactement la donnée de
cet élément. Le dossier a le critère, non le nom, et ne dit pas si la condition
est suffisante --- l'espoir qu'elle le soit est barré de plusieurs diagonales
sur la page.}

L'application, au n\textsuperscript{o} 10, est aux motifs semi-simples : les
graduations naturelles sur les $H^{\bullet}$ proviennent d'une graduation de la
catégorie des motifs, et
\begin{quote}
« ce qui n'est qu'une traduction du fait que les correspondances de Künneth
cohomologiques sur une variété algébrique projective lisse sont algébriques.
On respire ! »
\end{quote}
avec l'avertissement immédiat : ces $H^{\bullet}$ ne sont pas des foncteurs
fibres à proprement parler, mais des foncteurs fibres \emph{tordus}. « On ne
pourra s'empêcher de les appeler quand-même foncteurs fibres ! » Pour des
motifs pas nécessairement semi-simples la catégorie n'est plus graduée mais
filtrée, et ce qu'on veut est que les foncteurs cohomologiques se factorisent
par elle en foncteurs fibres \emph{véritables}, ce qui « n'est pas idiot » :
la symétrie de Koko ayant été mise une fois pour toutes sur
$\operatorname{Grad}(\underline{M})$, on s'est débarrassé de la torsion.

\subsection*{Les données de symétrie}

Le n\textsuperscript{o} 11 demande des théorèmes de structure intrinsèques, en
termes de gerbes, et le n\textsuperscript{o} 12 ouvre « quelques cogitations
dans ce sens » : $\underline{C}$ munie de $F$ additif exact fidèle
$k$-linéaire, donc de $U$ et $A$, et d'un $\otimes$ compatible, on se propose
de déterminer \emph{toutes} les données de symétrie $\psi$ pour ce $\otimes$
--- pas seulement celles compatibles à $F$ et à la symétrie ordinaire. On
trouve qu'elles correspondent aux éléments inversibles
\[
\sigma \in \bigl( U \mathbin{\widehat{\otimes}} U \bigr)^{*}
\]
vérifiant
\[
(1)\quad \sigma'\sigma = 1 ,
\qquad
\sigma\,(\Delta u)\,\sigma^{-1} = (\Delta u)' \quad \text{pour tout } u \in U ,
\]
où $x \mapsto x'$ est l'automorphisme de symétrie de $U
\mathbin{\widehat{\otimes}} U$. Il souligne ce que dit la seconde relation :
$\sigma$ transforme $\Delta$ en la diagonale opposée $\Delta'$, et $\sigma'$
fait l'inverse --- « condition de symétrie tordue de l'application
diagonale ». En présence d'une unité bilatère,
\[
(2)\quad (\mathrm{id} \otimes \varepsilon)(\sigma) =
(\varepsilon \otimes \mathrm{id})(\sigma) = 1 ,
\quad\text{i.e.}\quad
\sigma \in 1 + U^{+} \mathbin{\widehat{\otimes}} U^{+} ,
\]
et en présence d'associativité, la compatibilité $\varphi$--$\psi$ s'exprime
par une relation d'hexagone
\[
(3)\quad (\Delta \mathbin{\widehat{\otimes}} \mathrm{id})(\sigma)
= \sigma^{13}\sigma^{23} ,
\qquad
(\mathrm{id} \mathbin{\widehat{\otimes}} \Delta)(\sigma)
= \sigma^{13}\sigma^{12} .
\]
\footnote{La page écrit $(\mathrm{id} \mathbin{\widehat{\otimes}}
\Delta)(\sigma) = \sigma^{12}\sigma^{13}$ et ajoute : « NB je n'ai pas écrit le
calcul et ne garantis pas la forme exacte de la relation de l'hexagone ». La
forme standard est celle donnée ici ; la différence est une convention
d'indices, et cette lecture donne la forme correcte parce que c'est elle qui
rend l'énoncé vrai.}

Le triplet $(U, \Delta, \sigma)$ ainsi axiomatisé est ce qu'on appelle
aujourd'hui une bigèbre \emph{triangulaire}, et $\sigma$ une \emph{matrice
$R$ universelle} : les conditions (1), (2), (3) sont mot pour mot celles d'une
$R$-matrice avec $R_{21}R = 1$.\footnote{La vérification qu'il donne du cas
gradué en est la meilleure preuve. Sa boîte énonce que dans le cas du
n\textsuperscript{o} 6, gradué,
\[
\sigma = \pi^{+} \otimes \pi^{+} + \pi^{+} \otimes \pi^{-}
+ \pi^{-} \otimes \pi^{+} - \pi^{-} \otimes \pi^{-} ,
\]
et c'est exactement, réécrite dans les idempotents $\pi^{\pm} = (1 \pm g)/2$,
la $R$-matrice $\frac{1}{2}(1 \otimes 1 + 1 \otimes g + g \otimes 1 - g \otimes
g)$ de l'algèbre du groupe à deux éléments --- la $R$-matrice « super », celle
qui produit le signe de Koszul. Le dossier écrit l'élément et n'a aucun des
noms.}

\subsection*{Changement de base, et le torseur}

Reste à faire varier le corps. Pour $k'$ une $k$-algèbre commutative, on
détermine les foncteurs $G \colon \underline{C} \to \operatorname{Projf}(k')$
exacts et $k$-linéaires munis d'un $c$. La donnée de $G$ exact à droite revient
à celle d'un pro-module $L = (L_{\alpha})$ à droite sur $U$ muni d'une
structure de $k'$-module commutant aux opérations, par $G(M) = L \otimes_{U} M$
avec $L_{\alpha} = G(U_{\alpha})$ ; la $k$-linéarité dit que $L$ provient d'un
pro-module sur $U' = U \otimes k'$ ; et l'exactitude plus l'arrivée dans
$\operatorname{Projf}(k')$ disent que les $L_{\alpha}$ sont plats sur les
$U_{\alpha}$ et projectifs de type fini sur $k'$. La donnée de $c$ est alors
celle d'un isomorphisme de bi-$U'$-modules
\[
L \mathbin{\widehat{\otimes}}_{k'} L \xleftarrow{\ \sim\ }
L \mathbin{\widehat{\otimes}}_{k'} U' ,
\]
et posant $B = \operatorname{Homcont}_{k'}(L, k')$ --- une coalgèbre sur $A'$,
chaque $B_{\alpha}$ coplat sur $A_{\alpha}$ --- cela devient un isomorphisme de
comodules
\[
(*)\quad B \otimes_{k'} B \xleftarrow{\ \sim\ } B \otimes_{k'} A' ,
\]
d'où, par l'augmentation du facteur $A'$, une loi de composition
$B \otimes_{k'} B \to B$ compatible à la structure de $A'$-comodule. Sa marge
note, et c'est juste, que la donnée de $(*)$ équivaut à celle de cette loi.

Et alors tout se traduit :
\begin{quote}
$\Delta$ unitaire $\Leftrightarrow$ la loi de $B$ est unitaire ; $\Delta$
associative $\Leftrightarrow$ $B$ est une algèbre associative ; $\Delta$
commutative et $c$ compatible à la symétrie $\Leftrightarrow$ la loi de $B$ est
commutative.
\end{quote}
En termes du monoïde affine $H = \operatorname{Spec}A$, cela dit que $H_{k'}$
opère à droite sur $\operatorname{Spec}B$, et qu'après le changement de base
fidèlement plat $k' \to B$ l'action devient triviale : on a un
\emph{torseur} sous $H_{k'}$, et le nouveau foncteur fibre $G$ est déduit de
$F \otimes_{k} k'$ par torsion.\footnote{La marge ajoute « Ceci implique aussi
que $B$ est \emph{fidèlement} plat sur $k'$ », ce qui est la condition qui fait
du $\operatorname{Spec}B$ un torseur et non un pseudo-torseur. Et l'insertion
tapée en tête de la page 128 demande : « Sous réserve $B$ associative unitaire,
si $A$ est commutative, $B$ l'est-elle aussi ? Donc, la condition de comp.\ aux
sym.\ pour un $c$ est-elle surabondante, en prés.\ des autres compat.\ ? » La
réponse est oui pour une catégorie tannakienne sur un corps : les foncteurs
fibres de $\operatorname{Rep}(G)$ sur $k'$ sont les $G$-torseurs sur $k'$, dont
les algèbres sont commutatives. La page pose la question et ne la tranche pas.}

La dernière question du paragraphe est celle du $\tau$. Supposons $A$
seulement associative unitaire avec antipodisme, et donnons-nous $\sigma$.
Posant $V = G(U) = L \mathbin{\widehat{\otimes}}_{U} (U, \mathrm{int})$, le
pro-anneau qui définit $\underline{C}' = \underline{C} \otimes_{k} k'$ :
peut-on construire à partir de $\sigma$ un
$\tau \in V \mathbin{\widehat{\otimes}}_{k} V$ vérifiant les mêmes relations,
de façon que $(G,c)$ soit compatible aux symétries ? La seule chance qu'il voit
est que $\tau$ soit invariant, $U^{+}\!\cdot\tau = 0$. Et il dit ce que cela
donnerait :
\begin{quote}
« cela montrerait que dans l'étude des foncteurs fibres, il y a lieu de laisser
tomber les questions de compatibilité aux symétries, celles-ci se récupérant
toujours après coup. »
\end{quote}
Suit, en marge, le renvoi qui annule l'espoir : « Mais, cf.\ exemples
n\textsuperscript{o} 11 ».

\pagerange{129}{136}
\section*{Les exemples instructifs (pages 129 à 136)}

\subsection*{Le montage}

Soit $G$ un monoïde fini. Contrairement au yoga galoisien ordinaire, on prend
non pas $A = k(G)$ mais
\[
U = k(G) ,
\]
avec l'application diagonale déduite de la composition $G \times G \to G$ ---
c'est la bigèbre \emph{duale au sens de Cartier} de celle qu'on considère
d'habitude. La catégorie $\underline{C}$ est alors celle des familles
$(X_{g})_{g \in G}$ de vectoriels, avec le produit de convolution
\[
Z_{g} = \sum_{g'g''=g} X_{g'} \otimes Y_{g''}
\]
et pour unité la famille de Dirac. Elle a $\otimes$, $\varphi$, $\mathbf{1}$,
elle a $\psi$ si $G$ est commutatif, et une dualité si $G$ est un groupe. Un
foncteur fibre est la donnée d'une famille $(M_{g})$ avec des isomorphismes
\[
(*)\quad M_{gg'} \simeq M_{g} \otimes M_{g'} ,
\qquad
(**)\quad M_{e} \simeq k ,
\]
et cela impose que le rang des $M_{g}$ soit $0$ ou $1$.\footnote{Pour $g$
inversible, $(*)$ donne $\operatorname{rg}M_{e} = \operatorname{rg}M_{g}
\cdot \operatorname{rg}M_{g^{-1}}$, donc tous les rangs valent $1$ si $G$ est
un groupe ; pour un monoïde, un rang peut être nul, et c'est exactement ce
qu'exploite l'exemple a).}

\subsection*{a) Exact, compatible à tout, et non fidèle}

Prenons $G = \{e, p\}$ avec $p^{2} = p$, et $M_{e} = k$, $M_{p} = 0$. On n'a
plus le choix pour les isomorphismes $(*)$, qui satisfont toutes les
compatibilités. On obtient un $\otimes$-foncteur exact, compatible à $\varphi$,
$\psi$, $\mathbf{1}$, \emph{et non fidèle} --- donc pas localement isomorphe au
foncteur fibre donné. Et ce n'est possible que parce que $G$ n'est pas un
groupe.

D'où sa question, au bas de la page :
\begin{quote}
« Un $\otimes$-foncteur compatible avec tout entre $\otimes$-catégories
satisfaisant à tous les axiomes (y compris rigidité) est-il néc.\ fidèle et
exact ?? »
\end{quote}
La réponse est affirmative, et c'est l'un des premiers lemmes de la théorie
telle qu'elle s'est constituée : un $\otimes$-foncteur entre
$\otimes$-catégories abéliennes rigides sur un corps est fidèle et exact, la
dualité forçant l'un et l'autre. Le dossier pose la question et ne la tranche
pas ; son exemple montre qu'il faut bien la rigidité.

\subsection*{b) Exact, fidèle, et pas compatible à $\psi$}

Prenons maintenant $G$ un \emph{groupe}. La donnée d'un système de $M_{g}$ avec
les isomorphismes $(*)$, $(**)$ et leurs compatibilités est équivalente à la
donnée d'une \emph{extension de $G$ par $k^{*}$}, et le foncteur ainsi obtenu
est isomorphe à $F$ si et seulement si l'extension est scindée. Donc
\begin{quote}
l'ensemble des classes d'isomorphie de foncteurs fibres définis sur $k$ est en
bijection avec $H^{2}(G, k^{*})$.
\end{quote}
Or ce $H^{2}$ peut être non nul même pour $G$ commutatif --- le produit de deux
groupes cycliques isomorphes d'ordre premier à la caractéristique --- et $k$
algébriquement clos ; et surtout, pour $k$ algébriquement clos, \emph{un
élément non nul ne se tue dans aucune extension de $k$}, par le principe de
l'extension finie. Il y a donc des foncteurs fibres compatibles à un $\psi$
donné et d'autres qui ne le sont pas, et qui par suite ne sont pas localement
isomorphes aux premiers. Conclusion, soulignée :
\begin{quote}
« on ne peut négliger la contrainte de commutativité comme on aurait aimé
pouvoir le faire. De plus, cela suggère qu'en l'absence de contrainte de
commutativité, la classification des foncteurs fibres se fait par des groupes
du type de cohomologie de Brauer. »
\end{quote}
Si l'on se restreint aux $F'$ compatibles à $\psi$, on ne trouve plus que les
extensions \emph{commutatives}, qui pour $k$ algébriquement clos --- donc
$k^{*}$ divisible --- sont toutes scindées : $F'$ est alors isomorphe à $F$,
« en conformité avec la théorie galoisienne relative au cas où l'application
diagonale de $U$ est commutative ».\footnote{Cette catégorie est celle qu'on
note aujourd'hui $\mathrm{Vec}_{G}$ munie de la convolution, et le calcul est
celui des twists : les $\otimes$-foncteurs vers les espaces vectoriels
correspondent aux $2$-cocycles, donc à $H^{2}(G, k^{*})$. Le dossier fait le
calcul complet et n'a ni le mot « twist » ni le mot « catégorie de fusion ».}

La page 131 pousse plus loin. Sur cette $\underline{C}$, il y a une contrainte
de commutativité compatible aux autres données si et seulement si $G$ est
commutatif, et alors les contraintes correspondent aux fonctions $G \times G
\to k^{*}$ bilinéaires antisymétriques. Donc si $H^{2}(G, k^{*})$ est nul ---
$G$ cyclique, par exemple --- \emph{tout} choix non nul de $\sigma$ donne une
contrainte de commutativité pour laquelle il n'existe aucun foncteur fibre
compatible, sur aucune extension de $k$. Les foncteurs fibres gradués ne
sauvent pas la situation en caractéristique $2$, et sans doute en aucune. Reste
à donner l'exemple, qu'il demande et ne construit pas : une
$\otimes$-catégorie définissable par une bialgèbre, admettant une contrainte de
commutativité, mais \emph{aucune} contrainte pour laquelle un foncteur fibre
existe. Sa note du bas de page dit ce qu'il faut pour cela : $A$ commutative et
$U$ non commutative, sinon l'existence de $\sigma$ entraînerait la
commutativité de $\Delta$.

\subsection*{c) La rigidité tombe}

Soit $X$ propre sur $k$ avec $H^{0}(X, \mathcal{O}_{X}) = k$. La catégorie des
faisceaux cohérents a $\otimes$, $\varphi$, $\psi$, $\mathbf{1}$ et des
$\operatorname{Hom}$ internes, \emph{mais}
\[
\check{X} \otimes Y \longrightarrow
\underline{\operatorname{Hom}}(X,Y)
\]
n'est pas un isomorphisme : l'axiome de rigidité tombe, et « c'est pour cela
qu'il n'y a pas lieu de s'attendre à une vraie théorie de Galois ». Les
foncteurs fibres aux points de $X$ ont toutes les compatibilités mais ne sont
\emph{ni exacts ni fidèles}. Pour rétablir l'exactitude sans quitter
l'abélianité, on fixe un ouvert $U \subset X$ et on prend les faisceaux
cohérents munis sur $U$ d'une stratification : tout point de $U$ donne alors un
foncteur fibre exact --- mais non fidèle si $U \neq X$. Pour $U = X$ la
rigidité est satisfaite et l'on a une théorie de Galois, dont il doit résulter
que $\underline{C}$ n'a aucun foncteur fibre sur aucune extension de $k$.

Suivent quatre énoncés « A prouver », pour $\underline{C}$ abélienne à
$\otimes$ biadditif, $\varphi$, $\psi$, $\mathbf{1}$, dualité et rigidité, avec
$k = \operatorname{End}(\mathbf{1})$ : a) $k$ est absolument plat, donc un
composé fini de corps s'il est noethérien ; b) les
$\operatorname{Hom}(X,Y)$ sont de présentation finie sur $k$, qui est cohérent ;
c) $\underline{C}$ est limite inductive de sous-catégories stables par tout et
attachées à des composés de corps ; d) le cas d'un composé de corps se ramène à
celui d'un corps. Aucun n'est prouvé ici.

\subsection*{13) La variante sans fidélité, et l'algèbre de Steenrod}

Le n\textsuperscript{o} 13 est la « variante utile du tapis de
n\textsuperscript{o} 6 », et son intérêt est de lâcher deux hypothèses :
$\underline{C}$ n'est pas supposée additive, ni $F$ fidèle. On pose
$U = \operatorname{End}_{k}(F)$, $A$ la coalgèbre duale, et l'on trouve la
factorisation canonique de $F$ par
$\operatorname{Modf}(U/k) = \operatorname{Comodf}(A/k)$. Si l'on se donne un
bifoncteur $\otimes$ et un isomorphisme $F(X \otimes Y) \simeq F(X) \otimes
F(Y)$, on obtient une application diagonale sur $U$ : car l'anneau des
endomorphismes du bifoncteur $(X,Y) \mapsto F(X \otimes Y)$ est
$U \mathbin{\widehat{\otimes}}_{k} U$, et $U$ s'y envoie par l'isomorphisme
donné. Les compatibilités s'expriment comme au n\textsuperscript{o} 6, et ce
qui reste spécifique au cas fidèle sont les \emph{réciproques} : comment
déduire d'une diagonale sur $U$ un $\otimes$ sur $\underline{C}$.

L'exemple qu'il nomme est les opérations cohomologiques sur $\mathbf{F}_{p}$ :
sauf le morceau engendré par les $\pi_{i}$, les degrés dans $U$ sont $> 0$, et
$\widetilde{U} = k + U^{+}$ est le complété canonique d'un anneau gradué à
application diagonale anticommutative, donc $\widetilde{A}$ est un anneau
analogue.\footnote{$U$ est l'algèbre de Steenrod et $A$ son dual, qui est
précisément une algèbre de Hopf graduée commutative. C'est l'exemple qui
justifie d'avoir lâché la fidélité, et il n'est pas nommé sur la page.}

\subsection*{Le théorème de liberté}

La page 134 énonce, sous un titre dont le nom est resté en blanc :
\begin{quote}
\textbf{Théorème de ——— (rapporté par Husemöller).} Soient $A$ une bialgèbre
sur un corps $k$, associative et unitaire pour ses deux structures, cocomplète
pour sa filtration canonique de comodule ; $M$ un $A$-module muni d'un
homomorphisme diagonal unitaire $M \to M \otimes_{k} M$ qui soit $A$-linéaire,
d'une unité $k \to M$ qui soit un homomorphisme de coalgèbres unitaires --- de
sorte que l'image $e_{M}$ de $1$ vérifie $\varepsilon_{M}(e_{M}) = 1$ et
$\Delta_{M}(e_{M}) = e_{M} \otimes e_{M}$ --- $M$ cocomplet, et $u \mapsto
u.e$ injectif. Alors $M$ est un $A$-module \emph{libre} : posant
$\overline{M} = M \otimes_{A} k$, il existe un isomorphisme de $A$-modules
$M \simeq A \otimes_{k} \overline{M}$, non canonique, compatible aux projections
sur $\overline{M}$.
\end{quote}
Le nom est laissé en blanc par la machine, et cette lecture le laisse en
blanc.\footnote{L'énoncé est, dans le vocabulaire d'aujourd'hui, le théorème
fondamental des modules de Hopf : un module de Hopf sur une algèbre de Hopf est
libre, isomorphe à $A \otimes_{k} M^{\mathrm{co}A}$. Cette lecture ne
propose pas de remplir le blanc de la page.}

Deux exemples l'éclairent. Si la diagonale de $A$ est commutative, $A$ est
l'hyperalgèbre d'un groupe formel connexe $G$ ; si $\Delta_{M}$ est associative
et commutative, $M$ est l'algèbre des distributions à l'origine d'une variété
formelle connexe $V$, et l'hypothèse d'injectivité dit que $g \mapsto g.a$ est
un monomorphisme. Alors $G \times V \to V \times V$ est un monomorphisme par
Nakayama, le quotient $W = V/G$ existe, $V \to W$ est la projection d'un
$G_{W}$-torseur, et $\overline{M}$ est la coalgèbre de $W$. Le torseur n'est
pas trivial en général : on peut seulement dire que tout relèvement
$\overline{M} \to M$ comme $k$-espaces vectoriels donne l'isomorphisme
$A \otimes_{k} \overline{M} \simeq M$ de $A$-modules. Mais si $G$ est
formellement lisse --- en particulier en caractéristique nulle --- on a
splitting parfait.

Le second exemple suppose $A$ commutative pour sa loi d'algèbre, donc algèbre
affine d'un schéma en monoïdes affine $G$, et la condition de cocomplétude y
devient étrange : pour toute fonction $f$ sur $G$ il doit exister $n \geqslant
1$ tel que $f(g_{1}\cdots g_{n})$ soit somme de fonctions de $n-1$ des
variables. Sa conjecture est que pour $G$ de type fini cela équivaut à
$G$ \emph{fini unipotent}, et sans hypothèse de type fini à $G$
\emph{prounipotent}.

Le corollaire de la page 136 est net, et c'est lui qui ramène tout au
dossier :
\begin{quote}
\textbf{Corollaire.} Si $M \to M \otimes_{k} M$ induit un \emph{isomorphisme}
$M \otimes_{(A, \Delta_{A})}(A \otimes_{k} A) \to M \otimes_{k} M$, alors
$u \mapsto u.e$ est un isomorphisme pour toutes les structures.
\end{quote}
La démonstration est l'argument de rang : le rang de $M \otimes_{k}M$ sur
$A \otimes_{k}A$ est $n^{2}$, l'hypothèse donne $n = n^{2}$, donc $n = 1$ ---
et il note que l'argument ne marche que si l'on sait le rang fini, d'où la
variante : la diagonale de $\overline{M}$ est surjective, donc bijective
puisqu'unitaire, donc tout élément de $\overline{M} \otimes \overline{M}$ est
de la forme $\bar{e} \otimes x$, ce qui n'est possible que si $\bar{e}$
engendre.

Et la conclusion, qui renvoie explicitement à ses propres notes :
\begin{quote}
« Ceci montre que la conjecture de locale trivialité soulevée dans mes notes
sur les $\otimes$-catégories marche si la bialgèbre $A$ envisagée est finie sur
$k$, et sa duale est cocomplète, \emph{pourvu que} l'algèbre $B$ admette un
homomorphisme d'algèbres dans $k$ (ou une extension de $k$), i.e.\ pourvu que
$B$ ne soit pas identique à son idéal des commutateurs. »
\end{quote}
La conjecture en question est celle des pages 126 à 128 : que
$\operatorname{Spec}B$ soit un torseur et non un pseudo-torseur, c'est-à-dire
que deux foncteurs fibres soient localement isomorphes. Le corollaire en donne
un cas, sous une hypothèse dont il dit lui-même la fragilité.

\pagerange{138}{139}
\section*{Motifs à coefficients dans un corps de nombres (pages 138 et 139)}

Deux feuillets de sa main, d'une autre encre et d'un autre papier. Soient
$\underline{C}$ une $\otimes$-catégorie sur un corps $K$ et $K'$ une extension
\emph{finie}. On note $\underline{C}^{K'}$ la catégorie des objets de
$\underline{C}$ sur lesquels $K'$ opère, c'est-à-dire des $X$ munis d'un
homomorphisme de $K$-algèbres $K' \to \operatorname{End}(X)$. C'est encore une
$\otimes$-catégorie, en posant
\[
X \otimes_{K'} Y = \operatorname{coker}\Bigl(
X \otimes_{K} Y \xrightarrow{\ \lambda \otimes \mathrm{id} - \mathrm{id} \otimes \lambda\ }
X \otimes_{K} Y \Bigr)_{\lambda \in K'} ,
\]
et $\underline{\operatorname{Hom}}_{K'}(X,Y)$ comme le noyau correspondant dans
$\underline{\operatorname{Hom}}_{K}(X,Y)$. Quand $\underline{C}$ est définie
par un schéma en groupes affine $G$ sur $K$ muni d'un foncteur fibre, alors
$\underline{C}^{K'}$ est définie par $G_{K'} = G \otimes_{K} K'$, avec le
foncteur fibre déduit ; et le $\otimes$-foncteur naturel $\underline{C} \to
\underline{C}^{K'}$, $X \mapsto X \otimes_{K}K'$, a pour composé avec l'oubli
le foncteur $X \mapsto X \otimes_{K} \widetilde{K'}$, $\widetilde{K'}$ étant
l'espace vectoriel sous-jacent à $K'$.

Prenons $\underline{C}$ la catégorie des motifs semi-simples sur un corps $k$,
donc $K = \mathbf{Q}$, et $K'$ un corps de nombres. Un motif à coefficients
dans $K'$ est un $\mathbf{Q}$-motif $E_{0}$ muni de
\[
\varphi \colon K' \longrightarrow \operatorname{End}_{\mathbf{Q}}(E_{0}) .
\]
Une polarisation sur $E_{0}$, de poids $\rho$, donne $\check{E}_{0} \simeq
E_{0}(\rho)$, donc $\operatorname{End}(\check{E}_{0}) \simeq
\operatorname{End}(E_{0})$, et par là une \emph{involution} sur
$\operatorname{End}(E_{0})$. Le motif dual de $E_{\varphi}$ est alors défini
par $\lambda \mapsto \varphi(\lambda)'$, où le prime est cette involution, et
l'énoncé qu'il encadre est
\[
\boxed{\ (E_{\varphi})^{\vee} = (E_{\varphi'})(\rho)\ }
\]
avec l'avertissement immédiat :
\begin{quote}
« Ce n'est pas en général un isom.\ $(E_{\varphi})^{\vee} \simeq
E_{\varphi}(\rho)$ ! »
\end{quote}
Autrement dit, le dual d'un motif à coefficients n'est pas le tordu du motif
lui-même : c'est le tordu du motif à coefficients \emph{conjugués}.\footnote{
L'involution en jeu est celle que la polarisation induit sur l'algèbre des
endomorphismes, c'est-à-dire l'involution de Rosati dans le cas d'une variété
abélienne ; $\varphi'$ est le conjugué de $\varphi$ pour elle. Le dossier ne la
nomme pas et l'obtient en composant l'isomorphisme $\check{E}_{0} \simeq
E_{0}(\rho)$ avec la transposition.}

La vérification qu'il donne est celle du cas le plus simple, et elle est
concluante. Les $K'$-motifs de niveau zéro et de poids $0$ forment une catégorie
équivalente à celle des $K'$-espaces de dimension finie munis d'une opération
continue de $\pi = \operatorname{Gal}(\bar{k}/k)$, le dual répondant à la
contragrédiente. Les $K'$-motifs de rang $1$ à isomorphisme près correspondent
aux caractères continus $\chi \colon \pi \to K'^{*}$, avec
$(L_{\chi})^{\vee} = L_{\check{\chi}}$, $\check{\chi}(g) = \chi(g^{-1})$. Donc
\[
L_{\chi}^{\vee} \simeq L_{\chi}
\iff \chi = \check{\chi}
\iff \chi \text{ à valeurs dans } \pm 1 ,
\]
ce qui est bien une condition, et non une identité.

\pagerange{142}{154}
\section*{Les anneaux de représentations (pages 142 à 154)}

Treize feuillets de sa main, d'un trait rapide, qu'il numérote lui-même de
$(1)$ à $(34)$. Le sujet change : ce n'est plus la reconstruction d'un groupe,
c'est le calcul de l'anneau de ses représentations, et des opérations
$\lambda$-adiques qui le structurent.

\subsection*{Le formalisme}

Pour $G$ un groupe fini, $A^{G}$ est la catégorie des $k$-espaces de dimension
finie où $G$ opère --- une $\otimes$-catégorie avec involution. Elle dépend
\emph{contravariantement} de $G$ : $u \colon G \to G'$ donne
\[
A^{u} \colon A^{G'} \longrightarrow A^{G} ,
\qquad A^{\mathrm{id}} = \mathrm{id} ,
\qquad A^{uv} = A^{v}A^{u} ,
\]
et par là le produit tensoriel \emph{extérieur}
\[
M \boxtimes M' = A^{p_{1}}(M) \otimes A^{p_{2}}(M')
\ \in\ A^{G \times G'} ,
\]
d'où le $\otimes$ interne par la diagonale, $M \otimes M' = A^{\Delta_{G}}
(M \boxtimes M')$. Dans l'autre sens vient l'induction $I^{u}$,
\emph{covariante} et additive, avec $I^{uv} = I^{u}I^{v}$,
\[
A^{u}\bigl( I^{u}(M) \bigr) = M^{[G':G]}
\qquad (G \subset G') ,
\]
deux homomorphismes fonctoriels $M \to A^{u}I^{u}(M)$ et $M' \to
I^{u}(A^{u}(M'))$, et $I^{u} = A^{u^{-1}}$ quand $u$ est un isomorphisme.

Avec les injections $u_{p_{1},\ldots,p_{h}}$ d'un produit de groupes
symétriques dans $\mathfrak{S}_{p}$, on définit le \emph{produit d'induction}
\[
M * M' = I^{u_{p,q}}(M \boxtimes M')
\qquad
(M \in A^{\mathfrak{S}_{p}} , \ M' \in A^{\mathfrak{S}_{q}}) .
\]
Sa commutativité se lit sur l'élément $\alpha$ de $\mathfrak{S}_{p+q}$
vérifiant $\alpha\,u_{p,q}\,\alpha^{-1} = u_{q,p}\,s_{p,q}$, et son
associativité passe par $u_{p,q,r}$. Enfin $T^{n}(M) = \otimes^{n}M$, vu dans
$(A^{G})^{\mathfrak{S}_{n}}$, est un foncteur \emph{non additif}, et sa
propriété essentielle est celle qu'il encadre :
\[
\boxed{\ T^{n}(M \oplus M') \simeq \sum_{p+q=n} T^{p}(M) * T^{q}(M')\ }
\]
avec, pour la composition, $T^{m}(T^{n}(M)) \simeq
A^{u_{m,n}}(T^{mn}(M))$ à travers l'homomorphisme
$u_{p,q} \colon \mathfrak{S}_{p} \times \mathfrak{S}_{q} \to
\mathfrak{S}_{pq}$ construit sur les permutations $\alpha_{p,q}(\sigma)$
croissantes sur les intervalles $\,]iq, (i+1)q]\,$ et les envoyant sur les
$\,]\sigma(i)q, (\sigma(i)+1)q]\,$ ; et
$T^{n}(M \otimes M') \simeq T^{n}(M) \otimes T^{n}(M')$.

\subsection*{Les groupes de classes, et l'anneau des fonctions symétriques}

Passant aux classes, $R(G) = K(A^{G})$ est un \emph{anneau de Chern}
commutatif augmenté, contravariant en $G$, avec
\[
\alpha \boxdot \beta = R(p_{1})(\alpha) \cdot R(p_{2})(\beta) ,
\]
et l'induction $I(u)$ est additive \emph{mais pas multiplicative}. Les
composés valent
\[
R(u)\,I(u) = I(u)\,R(u) = \text{multiplication par } [G':G] ,
\]
et un automorphisme intérieur agit par l'identité : $R(\sigma_{g}) =
I(\sigma_{g}) = \mathrm{id}$.\footnote{Les deux lignes de la formule (26)
portent le même facteur $[G':G]$ sur le feuillet. Le second composé,
$I(u)R(u)$, est bien la multiplication par l'indice ; le premier, $R(u)I(u)$,
est la multiplication par l'indice seulement lorsque $u$ est l'inclusion d'un
sous-groupe \emph{normal}, et vaut en général la somme sur les classes
doubles. La transcription donne les deux lignes telles quelles ; la leçon est
celle de la page.}

L'objet qui sort est l'anneau gradué
\[
R^{\bullet} = \coprod_{n \geqslant 0} R(\mathfrak{S}_{n}) ,
\qquad R^{0} = \mathbf{Z} ,
\]
muni du produit d'induction $*$, et c'est ce qu'on appelle aujourd'hui
\emph{l'anneau des fonctions symétriques}.\footnote{La caractéristique de
Frobenius identifie $\coprod_{n}R(\mathfrak{S}_{n})$, muni du produit
d'induction, à l'anneau $\Lambda$ des fonctions symétriques ; le foncteur
$T^{n}$ devient la $n$-ième puissance tensorielle totale et la formule encadrée
la formule d'addition des puissances, et les relations (29) et (30) sont la
restriction d'un induit. Le dossier construit l'anneau et n'a pas le nom.} La
relation qui relie la multiplication des $R^{n}$ à la multiplication globale
est
\[
R(u_{p,q})(\alpha * \beta) = \frac{(p+q)!}{p!\,q!}\,\alpha \boxdot \beta ,
\qquad\text{et en particulier}\qquad
\frac{(2p)!}{(p!)^{2}}\,\alpha\beta = R(u_{p})(\alpha * \beta)
\]
où $u_{p}$ est le composé $\mathfrak{S}_{p} \xrightarrow{\Delta}
\mathfrak{S}_{p} \times \mathfrak{S}_{p} \xrightarrow{u_{p,p}}
\mathfrak{S}_{2p}$. Et la formule d'addition devient un homomorphisme de
groupes
\[
T \colon R(G) \longrightarrow 1 + \widehat{R(G)^{\bullet +}} ,
\qquad
T(\alpha) = \sum_{n} T^{n}(\alpha) , \quad T^{0}(\alpha) = 1 ,
\]
soumis aux deux règles
\[
T^{n}(\alpha + \beta) = \sum_{p+q=n} T^{p}(\alpha) * T^{q}(\beta) ,
\qquad
T^{n}(\alpha \cdot \beta) = T^{n}(\alpha)\,T^{n}(\beta) ,
\]
et fonctoriel en $G$. La page 149 récapitule le tout en quatre articles : les
anneaux de Chern $R^{n}$, les $R(u)$, les applications bilinéaires $R^{p}
\times R^{q} \to R^{p+q}$ avec la relation ci-dessus, et $T$.\footnote{La
récapitulation de la page 149 écrit $T^{n}(\alpha + \beta) = \sum_{p+q=n}
T^{p}(\alpha)T^{q}(\alpha)$, avec $\alpha$ deux fois là où l'énoncé demande
$T^{q}(\beta)$. La transcription le signale et donne la leçon ; cette lecture
rétablit $\beta$.}

\subsection*{Un tore}

Pour $T$ un tore, $R(T) \simeq \mathbf{Z}[X(T)]$, l'algèbre du groupe des
caractères. Une représentation $\rho \colon T \to Gl(n)$ se diagonalise en $n$
caractères $\alpha_{1}, \ldots, \alpha_{n}$, et l'on pose $C^{i}(\rho)$ pour
leurs fonctions symétriques élémentaires, $C(\rho) = \sum_{i}C^{i}(\rho)$ ---
de sorte que $C(\rho) = C(\rho')C(\rho'')$ pour une extension. On prolonge
$\rho \mapsto (\deg \rho, C(\rho))$ en un homomorphisme d'anneaux de Chern
\[
R(T) \simeq \mathbf{Z}[X(T)]
\longrightarrow \mathbf{Z} \times \bigl( 1 + \mathbf{Z}[[X(T)]]^{+} \bigr) ,
\qquad
\chi \longmapsto (1, 1 + [\chi]) ,
\]
et l'assertion est qu'il est \emph{injectif à image dense}. L'injectivité se
voit en composant avec $\chi \mapsto e^{\chi}$ dans
$\mathbf{Q}[[X(T)]]$ et en notant que les $e^{\chi}$ sont linéairement
indépendants. La densité se ramène, par récurrence, à un lemme sur les
fonctions symétriques : étant donné un polynôme $C^{n}$ de degré $n$ en
$X_{1}, \ldots, X_{r}$, trouver des formes linéaires $\alpha^{1}, \ldots,
\alpha^{N}$ telles que $\sigma^{1} = \cdots = \sigma^{n-1} = 0$ et
$\sigma^{n}((\alpha^{i})) = C^{n}$. Le lemme est énoncé au haut de la page 151
et \emph{sa démonstration n'y est pas}.

Le reste de la page est un schéma au crayon qui dit à quoi cela devait servir :
le fibré universel $E_{T} \to B_{T}$, la surjection $\mathbf{Z}[X(T)] \to
H^{*}(B_{T}, \mathbf{Z})$, et $R(T)$ atteignant ce dernier de deux façons ---
« alg » d'un côté, « top.\ via Chern » de l'autre --- avec la même figure pour
$H^{*}(X, \mathbf{Z})$ et l'homomorphisme caractéristique entouré. Et la tâche
laissée :
\begin{quote}
« Il faudrait prouver que $\mathbf{Z}[X(T)] \to$ l'anneau de Chern attaché à
$R(T) \simeq \mathbf{Z}[X(T)]$. »
\end{quote}

\subsection*{L'homomorphisme caractéristique}

Les trois derniers feuillets passent à un groupe algébrique. $K(G)$ est
l'anneau engendré par les classes des représentations rationnelles
irréductibles --- libre comme groupe additif, et anneau de Chern --- et
$CK(G)$ « l'anneau des classes de Chern » de $K(G)$, gradué, avec
l'homomorphisme canonique fonctoriel
\[
c_{G} \colon K(G) \longrightarrow
\mathbf{Z} \times \bigl( 1 + CK(G)^{+} \bigr) ,
\]
sur lequel les automorphismes intérieurs agissent trivialement. Pour $P$ un
fibré principal de groupe $G$ et de base $X$, chaque représentation donne un
fibré vectoriel sur $X$, d'où $K(G) \to K(X)$ et l'\emph{homomorphisme
caractéristique}
\[
CK(G) \longrightarrow CK(X) = A(X) ,
\]
compatible aux deux carrés évidents. La proposition de la page 153 calcule le
cas d'un tore de dimension $r$ :
\[
K(T) \simeq \mathbf{Z}(X(T))
\simeq \mathbf{Z}[x^{1}, \ldots, x^{r}]_{x^{1}\cdots x^{r}} ,
\qquad
CK(T) \simeq \mathbf{Z}[X(T)] ,
\]
avec $\varepsilon(\chi) = 1$ et $c_{T}(\chi) = 1 + \chi$, c'est-à-dire
$c_{T}^{1}(\chi) = \chi$ et $c^{i}_{T}(\chi) = 0$ pour $i >
1$.\footnote{L'opération $\lambda$ change de signe d'une ligne à l'autre sur le
feuillet : $\lambda x^{i} = 1 - x^{i}t$ dans l'accolade, $\lambda(\chi) = 1 +
\chi t$ à la ligne suivante. Les deux leçons sont celles de la page ; c'est la
seconde qui s'accorde avec $c_{T}(\chi) = 1 + \chi$.}

Et la page 154 pose la question qui reste ouverte. Soient $T$ un tore maximal,
$N$ son normalisateur, $W = N/T$ ; l'homomorphisme canonique $CK(G) \to
CK(T) \simeq \mathbf{Z}[X(T)]$ est compatible aux opérations de $N$, qui opère
trivialement à gauche et par $W$ à droite. On a donc
\[
(5)\qquad CK(G) \longrightarrow \mathbf{Z}[X(T)]^{W} ,
\]
et :
\begin{quote}
en général ce n'est pas une bijection. Mais c'est « sans doute » une bijection
s'il n'y a pas de torsion ; et l'homomorphisme analogue
\[
(5\,\mathrm{bis})\qquad C\bigl( K(G) \otimes_{\mathbf{Z}} k \bigr)
\longrightarrow k[X(T)]^{W}
\]
est bijectif si la caractéristique de $k$ est première à la torsion de $G$.
\end{quote}
Ni l'un ni l'autre n'est démontré ; le feuillet s'arrête là et la page 155 est
blanche.\footnote{Les deux énoncés sont ceux qu'on connaît aujourd'hui sous
d'autres noms : la restriction $R(G) \to R(T)^{W}$ est un isomorphisme pour $G$
connexe réductif à groupe dérivé simplement connexe, et $H^{*}(BG;\mathbf{Z})
\simeq (\operatorname{Sym}X(T))^{W}$ en l'absence de torsion. Le dossier
conjecture le second sous la forme (5) et en donne la variante $k$-linéaire.}

\pagerange{156}{179}
\section*{Le dictionnaire refait à la main (pages 156 à 179)}

Vingt-quatre feuillets de sa main, écrits plus vite encore, et les plus raturés
du dossier. C'est la seconde construction du dictionnaire, faite de rien, et
c'est elle qui produit les objets dont les quinze dernières pages ont besoin.

\subsection*{L'algèbre de la catégorie}

Soit $\mathcal{C}$ une classe abélienne dont les objets sont des modules sur un
anneau commutatif $k$, les morphismes des homomorphismes de modules, telle que
le foncteur d'oubli soit \emph{exact} et qu'il existe un \emph{ensemble}
$(A_{i})$ de représentants des classes d'isomorphie. Noyaux, images et conoyaux
y sont les ordinaires.

\textbf{Lemme I.} Les endomorphismes du foncteur d'oubli sont exactement les
familles $(u_{i})$ avec $u_{i} \in \operatorname{End}_{k}(A_{i})$, telles que
pour tout $\varphi \in \operatorname{Hom}_{k}(A_{i}, A_{j})$ on ait
$\varphi u_{i} = u_{j}\varphi$.

Ils forment donc un sous-anneau de $\prod_{i}\operatorname{End}(A_{i})$,
\emph{fermé} pour la topologie produit quand chaque
$\operatorname{End}(A_{i})$ porte la topologie de la convergence simple, les
$A_{i}$ étant discrets. D'où une algèbre linéairement topologisée, séparée et
\emph{complète},
\[
U ,
\]
pour laquelle tout objet $V$ de $\mathcal{C}$ est un $U$-module, l'application
$U \to \operatorname{End}(V)$ étant continue, et pour laquelle les
$k$-homomorphismes sont des $U$-homomorphismes. Un $U$-module est dit
\emph{permis} quand $U \to \operatorname{End}(V)$ est continue ; si les $A_{i}$
sont de type fini sur $k$, les $\operatorname{End}(A_{i})$ sont discrets et $U$
est linéairement topologisée avec des idéaux ouverts, et on peut se borner aux
permis de type fini.

Le foncteur canonique $I$ de $\mathcal{C}$ vers les $U$-modules permis est
alors l'objet du dictionnaire, et l'on demande deux choses de lui :
\begin{enumerate}
\item[$1^{\circ}$)] que pour tout $k$-module $V$, l'application $K(V) \to
K'(V)$ des structures « de type $\mathcal{C}$ » vers les structures de
$U$-module permis soit \emph{bijective} ;
\item[$2^{\circ}$)] que pour $V, W$ dans $\mathcal{C}$, les $k$-homomorphismes
soient \emph{exactement} les $U$-homomorphismes.
\end{enumerate}
La première structure supplémentaire est une \emph{augmentation continue}
$\varepsilon$ sur $U$, qui singularise le module trivial.

Et la première proposition de fonctorialité, soulignée sur la page 159 :
\begin{quote}
les homomorphismes continus $V \to U$ de deux telles algèbres s'identifient aux
foncteurs permis $\mathcal{C}^{U} \to \mathcal{C}^{V}$ au-dessus de
$\mathcal{C}^{k}$ ;
\end{quote}
et ceux qui sont compatibles aux augmentations correspondent aux foncteurs qui
transforment le module trivial en le module trivial.

\subsection*{Dualité, puis produit tensoriel}

La dualité (pages 160 et 161) est un foncteur $F^{\vee}$ mettant sur le dual
$V'$ une structure de $U$-module permise, au-dessus du passage au dual des
$k$-modules, avec $F^{\vee}F^{\vee} = \mathrm{id}$. Un tel foncteur est un
homomorphisme de l'anneau opposé $U^{\circ}$ dans $U$, c'est-à-dire une
application $x \mapsto \check{x}$ \emph{anti-multiplicative}, et la relation
$F^{\vee}F^{\vee} = \mathrm{id}$ en fait un \emph{anti-automorphisme
involutif}. La compatibilité avec l'augmentation s'écrit
\[
\varepsilon(\check{x}) = \varepsilon(x)
\]
et signifie que les modules triviaux vont sur les modules triviaux ; celle d'un
homomorphisme continu $\varphi \colon V \to U$ s'écrit $\varphi(\check{x}) =
\check{\overline{\varphi(x)}}$ et signifie que le foncteur permis commute au
passage au dual.

Le produit tensoriel (pages 162 à 164) est plus délicat, et ces trois feuillets
sont les plus raturés du lot. Ce qu'on veut est un bifoncteur $T(V,W)$
au-dessus de $\otimes_{k}$, c'est-à-dire une structure de $U$-module sur
$V \otimes_{k}W$, naturelle en les deux arguments et transformant les
$U$-homomorphismes en $U$-homomorphismes. Le chemin est celui-ci. Un module
permis sur $U \mathbin{\widehat{\otimes}} U$ définit fonctoriellement un
$U$-module permis, et cela revient à un homomorphisme
\[
\boxed{\ U \longrightarrow U \mathbin{\widehat{\otimes}} U\ }
\]
--- l'énoncé qu'il souligne page 163. La vérification passe par les quotients
$U_{V} = U/J_{V}$, qui sont \emph{artiniens}, par Krull-Schmidt sur les
indécomposables $W_{ijk}$ de $V_{i} \otimes V_{j}$, et par l'inclusion
\[
U/J \otimes_{k} U/J \subset
\operatorname{End}_{k}(V) \otimes_{k} \operatorname{End}_{k}(W)
= \operatorname{End}_{k}(V \otimes_{k} W)
\]
qui montre que les $U$-structures sur $V \otimes V$ sont les $\sum u_{i}
\otimes v_{i}$. « On vérifie que c'est bon. »

\subsection*{Les axiomes, rassemblés}

La page 165 énumère les axiomes supplémentaires : l'associativité de $\Delta$,
accompagnée de son axiome pour les produits tensoriels ; la commutativité,
c'est-à-dire que $V \otimes_{k}W \to W \otimes_{k}V$ soit un
$U$-homomorphisme ; l'augmentation, avec $\Delta(u) = 1 \otimes u + u \otimes
1 + \cdots$ pour les primitifs ; la compatibilité avec la dualité,
$(V \otimes_{k}W)' \simeq V' \otimes_{k} W'$ ; et le carré
\[
\Delta \circ \vee = (\vee \otimes \vee) \circ \Delta ,
\qquad\text{c'est-à-dire}\qquad
\vee \ \text{est un isomorphisme de co-algèbres.}
\]
La compatibilité \emph{combinée} de $\varepsilon$, $\vee$ et $\Delta$ (page
166) est la condition que l'accouplement naturel $V' \otimes V \to k$ soit un
$U$-homomorphisme pour $k$ trivial, ce qui s'écrit
\[
P(v \otimes 1)\Delta = P(1 \otimes v)\Delta = 1 ,
\]
$P$ étant la multiplication $U \mathbin{\widehat{\otimes}} U \to U$. Et deux
algèbres $U$, $V$ munies de $\Delta$ donnent un foncteur permis $\varphi
\colon \mathcal{C}^{V} \to \mathcal{C}^{U}$ compatible aux produits tensoriels
\emph{si et seulement si} $\varphi$ commute à $\Delta$.

La page 167 rassemble le tout en quatre articles, et c'est la définition
complète :
\begin{itemize}
\item $U$ une algèbre associative topologique avec assez de représentations de
dimension finie ;
\item $U$ munie d'une structure de co-algèbre associative $U \to U
\mathbin{\widehat{\otimes}} U$, les deux structures compatibles, c'est-à-dire
un homomorphisme continu d'algèbres ;
\item $U$ munie d'une augmentation $\varepsilon$, qui est un homomorphisme
d'algèbres \emph{et} de co-algèbres ;
\item $U$ munie d'une involution $\vee$, anti-automorphisme involutif
d'algèbres et automorphisme de co-algèbres, avec $\varepsilon(\check{x}) =
\varepsilon(x)$ ;
\end{itemize}
plus la relation $P(v \otimes 1)\Delta = P(1 \otimes v)\Delta = 1$. Et « par
dualité », les mêmes données vivent sur une algèbre commutative $A$.

\subsection*{Le groupe}

Vient alors la définition, page 168 :
\[
G = \bigl\{ s \in U \ \text{inversible} \ :\ \Delta s = s \otimes s ,\
\varepsilon(s) = 1 \bigr\} .
\]
Et deux faits, démontrés. D'abord $\Delta s = s \otimes s$ donne
$\varepsilon(s) = \varepsilon(s)^{2}$, donc $\varepsilon(s)$ vaut $0$ ou $1$,
et pas $0$ puisque $s$ est inversible : \emph{la condition $\varepsilon(s) = 1$
est automatique}. Ensuite $P(v \otimes 1)\Delta = 1$ donne $\check{s}s = 1$,
donc
\[
\check{s} = s^{-1} :
\]
l'involution est l'inverse sur $G$.

Posant $A = U'$, le dual topologique, la page 169 identifie ses éléments :
\begin{quote}
les formes linéaires continues sur $U$ sont exactement les $u \mapsto
\langle u \cdot a, b \rangle$ pour $a$, $b'$ dans un objet $V$ --- c'est-à-dire
les \emph{coefficients} des représentations linéaires continues de dimension
finie.
\end{quote}
Et trois classes d'éléments sont distinguées : le groupe $U^{*}$ des
inversibles, qui répond aux automorphismes du foncteur d'oubli ; les éléments
\emph{primitifs} ; et les endomorphismes du foncteur additif qui, à des
$A$, $B$ sur $V$, $W$, associent
\[
A \otimes 1 + 1 \otimes B \quad\text{sur } V \otimes W ,
\]
lesquels forment une \emph{algèbre de Lie}.\footnote{Les deux dernières classes
sont la même notion, énoncée deux fois : la seconde est la description
fonctorielle de la première. Ce sont les éléments primitifs de $U$,
c'est-à-dire les $\otimes$-dérivations du foncteur d'oubli, et ils forment
l'algèbre de Lie de $G$. Sa parenthèse « car $p \neq 0$ » signale que cette
algèbre de Lie ne détermine pas $G$ en caractéristique positive.}

\subsection*{« $G$ est total dans $U$ »}

La page 170 fait de $A$ une algèbre de \emph{fonctions sur $G$} : $A \to
\{\text{fonctions } G \to k\}$, homomorphisme d'algèbres transportant $1$ en
$\mathbf{1}$, compatible avec la multiplication diagonale, l'augmentation et
$\vee$. Et l'énoncé qu'il souligne :
\begin{quote}
dire que cette application est injective signifie précisément que \emph{$G$ est
total dans $U$},
\end{quote}
c'est-à-dire que l'adhérence de $G$ pour la topologie faible est tout
$U$.\footnote{C'est la condition qui, dans le vocabulaire d'aujourd'hui, dit
que les points rationnels du schéma en groupes $\operatorname{Spec}A$ sont
Zariski-denses : les coefficients séparent les points, et $A$ est vraiment une
algèbre de fonctions et non seulement une algèbre affine. L'exemple $A =
k[x]/x^{p}$ qu'il donne aussitôt est l'exemple où cela tombe en défaut : le
groupe des points est trivial et ne voit rien.}

D'où le théorème, page 171 :
\begin{quote}
si $G$ est total dans $U$, la classe des $U$-modules permis s'identifie, avec
toutes ses structures, à la classe des $G$-modules dont les coefficients
tombent dans $A$.
\end{quote}
Ainsi une bigèbre $U$ donne un groupe $G$ et une algèbre $A_{0}$ de fonctions
sur $G$ séparant ses points, stable sous $\vee$, avec $f(xy) \in A_{0} \otimes
A_{0}$. Et la page \emph{renverse le problème} : partant de $G$ et d'une telle
$A_{0}$, on reconstruit $U$ avec $A_{0} \subset A$ et une bijection de $G$ sur
le groupe des éléments de type groupe de $U$.

Les pages 172 à 176 le démontrent, en passant par une algèbre associative
$\Omega$ avec $1$ : on prend pour $A$ l'ensemble des fonctions sur $\Omega$ qui
sont coefficients de représentations linéaires de dimension finie ; la
multiplication de $\Omega$ est continue pour $\sigma(\Omega \otimes \Omega, A
\otimes A)$ et $\sigma(\Omega, A)$, ce qui est un homomorphisme $A \to A
\otimes A$ ; les $\Omega$-modules de dimension finie sont exactement les
$\hat{\Omega}$-modules permis, $\hat{\Omega}$ étant le complété de $\Omega$
pour $\sigma(\Omega, A)$ ; on a $\langle u \otimes v, f \rangle = \langle uv, f
\rangle$ et $f(uv) \in A_{0} \otimes A_{0}$ ; et la classe des
$\Omega$-modules dont les $f_{x,x'}$ tombent dans $A_{0}$ est isomorphe à la
classe des $A_{0}$-comodules et à celle des $\hat{\Omega}$-modules pour
$\sigma(\Omega, A_{0})$, la réciproque tenant « grâce à la dimension finie ».
CQFD, page 176.

Reprenant le groupe, la stabilité de $A_{0}$ se lit en deux mots : $A_{0}$ est
stable sous la multiplication, \emph{autrement dit} $A_{0}$ est une
sous-algèbre, et $A_{0}$ est stable sous $\vee$. La page 177 ajoute $1 \in
A_{0}$ pour la représentation triviale, et demande que
\[
G \longrightarrow G'(A_{0})
\]
soit bijective, c'est-à-dire que $A_{0}$ sépare les points de $G$ et que toute
fonction de $A_{0}$ provienne de $G$ ; et la structure est
\emph{algébrique} quand $A_{0}$ est de type fini sur $k$.

\subsection*{Le schéma qui clôt la rédaction}

La page 179 porte le résultat de tout le faisceau sous forme de schéma :
\[
K\text{-catégorie}
\iff
K\text{-bigèbre } U
\iff
\text{hyperalgèbre commutative } A ,
\]
avec deux termes qui montent vers le membre médian : \emph{groupe
pseudo-algébrique} --- car si $G(U)$ est total dans $U$, on retrouve $A$ --- et
au-dessous \emph{groupe algébrique}, si de plus $A$ est de type
fini.\footnote{Sur la page, « algèbre » est biffé et « bigèbre » écrit
au-dessus dans le second membre. « Pseudo-algébrique » est ce qu'on appelle
aujourd'hui \emph{pro-algébrique} : un schéma en groupes affine quelconque, par
opposition à un groupe algébrique, qui est celui d'algèbre affine de type fini.
La flèche du bas dit exactement la proposition 11.13 du tapuscrit des pages 51
à 52 : tout groupe affine est limite projective de groupes algébriques.}

C'est le théorème du dossier, et c'est le seul endroit où il soit écrit en
entier.

\pagerange{180}{182}
\section*{L'hypergroupe des classes indécomposables (pages 180 à 182)}

\subsection*{Le groupe ordonné}

Soit $\mathcal{R}$ une classe abélienne dont les objets sont de longueur finie,
$I$ l'ensemble des classes d'\emph{indécomposables}, et
\[
\Gamma = \mathbf{Z}(I)
\]
le groupe abélien libre sur $I$. Krull-Schmidt en fait un \emph{groupe
ordonné} : le cône positif $\Gamma^{+} = \mathbf{N}(I)$ s'identifie
canoniquement au monoïde des classes d'isomorphie, la somme de $\Gamma^{+}$
répondant à la somme directe.\footnote{Ce $\Gamma$ n'est pas le groupe de
Grothendieck de la catégorie abélienne : celui-ci est libre sur les
\emph{simples} et ignore les extensions non scindées. C'est ce qu'on appelle
aujourd'hui le groupe de Grothendieck \emph{scindé}, ou l'anneau de Green quand
on y met la multiplication. La distinction est essentielle pour la suite : pour
le groupe additif d'un corps de caractéristique nulle, le groupe de Grothendieck
ordinaire est $\mathbf{Z}$, et c'est le scindé qui est intéressant.}

Deux structures supplémentaires s'y transportent. Si $\mathcal{R}$ porte un
produit tensoriel et si $k$ est muni d'une structure privilégiée, alors $I$ a
un \emph{élément distingué} $e$ ; et si $\mathcal{R}$ porte une involution
$\vee$, alors $I$ porte une involution $i \mapsto \check{\imath}$, dont la
compatibilité avec $e$ s'écrit
\[
\check{e} = e .
\]
Enfin $n(i) = \dim M_{i}$ est une fonction privilégiée sur $I$, qui se prolonge
par linéarité à $\Gamma$, et sur $\Gamma^{+}$ c'est la fonction « dimension
d'un module ». On a $n(e) = 1$ et $n(\check{\imath}) = n(i)$.

Ce qui reconstruit $\mathcal{R}$, page 181, est la donnée de $I$, de $n$, de
représentants $M_{i}$ avec $M_{i} = k^{n(i)}$ comme $k$-module, et pour chaque
couple $(i,j)$ d'un sous-espace $\operatorname{Hom}_{j,i}(k^{n(i)},
k^{n(j)})$, ces données étant stables par composition : $\operatorname{Hom}_{kj}
\circ \operatorname{Hom}_{ji} \subset \operatorname{Hom}_{ki}$.

\subsection*{L'hypergroupe discret}

Supposons maintenant $\mathcal{R}$ munie d'un produit tensoriel. Alors
$\Gamma$ devient un anneau commutatif dont la multiplication, restreinte à
$\Gamma^{+}$, est celle du produit tensoriel ; $n$ y est une fonction
\emph{multiplicative} ; la compatibilité avec $e$ indique que $e$ est l'unité ;
la compatibilité avec $\vee$ indique que $\vee$ est un \emph{automorphisme}
de l'anneau ; et enfin la compatibilité simultanée des trois indique que
\begin{quote}
$i \otimes \check{\imath}$ a toujours l'unité $e$ pour composante.
\end{quote}
Et il conclut : $I$ est alors un \emph{« hypergroupe discret »}.\footnote{Le
mot est le sien, avec l'accord de la page. C'est, terme pour terme, ce qu'on
appelle aujourd'hui un \emph{anneau à base} : un anneau muni d'une base à
coefficients de structure entiers positifs, dont l'unité est un élément de la
base, muni d'une involution de la base, et soumis à la condition que l'unité
apparaisse dans $i \otimes \check{\imath}$. La fonction $n$ en est une
\emph{fonction dimension}. Le dossier écrit les quatre axiomes et n'a que le mot
« hypergroupe ».}

\subsection*{Le cas d'un groupe}

La page 182 calcule ce que cela donne quand la catégorie vient d'un groupe. Si
tout objet de $\mathcal{C}$ est complètement réductible et si $\mathcal{C}$ est
engendrée par un groupe $G$ muni d'une algèbre $A$ de fonctions sur lui, alors
$\Gamma$ s'identifie à l'algèbre des combinaisons linéaires des représentations
--- l'anneau des caractères --- et $\Gamma^{+}$ aux fonctions des
représentations linéaires de $G$, \emph{toutes les structures étant
conservées} :
\[
n(f) \rightsquigarrow f(e) \ (e \in G) ,
\qquad
\vee \rightsquigarrow \vee ,
\qquad
e \rightsquigarrow \text{la fonction } \mathbf{1} \text{ sur } G ,
\]
et il note en passant qu'il s'agit de fonctions \emph{centrales}. Si $G$ est un
groupe algébrique semi-simple, on obtient ainsi une description concrète de
$\Gamma$ : c'est l'algèbre des fonctions sur $T/W$, $T$ un tore maximal et $W$
le groupe de Weyl.\footnote{C'est le théorème de restriction de Chevalley sous
sa forme la plus simple : $R(G) \simeq \mathbf{Z}[X(T)]^{W}$, et le spectre de
cet anneau est $T/W$.} Et la tâche qui reste :
\begin{quote}
« il faudrait essayer de trouver une caractérisation intrinsèque des
irréductibles, i.e.\ des éléments de $I$ ; pour cela il faut mettre sur $T/W$
une structure d'hypergroupe, dont on doit construire les fonctions
irréductibles »
\end{quote}
--- avec un point d'interrogation porté dans la marge en regard de cette
parenthèse.\footnote{La base canonique de $R(G) = \mathbf{Z}[X(T)]^{W}$ est
celle des caractères de Weyl, donnés par la formule des caractères ; c'est la
réponse moderne à « quels sont les irréductibles ». Qu'elle se lise sur $T/W$
seul, sans la donnée du système de racines, n'est pas ce que cette formule
dit.}

Le feuillet se clôt sur le premier cas calculé : pour $G = k^{*}$ avec $k$
algébriquement clos,
\[
\Gamma(G) \simeq \text{l'algèbre du groupe } \mathbf{Z} ,
\qquad I(G) \simeq \mathbf{Z} .
\]

\pagerange{183}{197}
\section*{Trois anneaux calculés (pages 183 à 197)}

Les quinze derniers feuillets font, à la main et jusqu'au bout, ce que les
pages 142 à 154 avaient annoncé : le calcul de l'anneau des représentations
dans trois cas. Ils sont numérotés par lui II et III pour les deux derniers ; le
premier ne porte pas de numéro.

\subsection*{Le groupe additif : le problème}

Soit $G = k$, le groupe additif, $k$ algébriquement clos. Toutes les
représentations rationnelles de $G$ sont \emph{unipotentes}, donc leur
classification est celle des opérateurs unipotents, donc celle des opérateurs
nilpotents, donc celle des $k[X]$-modules annulés par une puissance de $X$ ---
les $k[X]$-modules $X$-primaires de type fini. Les indécomposables sont les
$k[X]/X^{i}$, d'où
\[
I \simeq \mathbf{N}^{+} , \qquad i \leftrightarrow k[X]/X^{i} ,
\]
l'élément neutre est $1$, l'involution est l'identité --- et la question :
\begin{quote}
« Quelle est la multiplication ? »
\end{quote}
avec, en marge, la question qui commande tout : « Caractéristique $0$ ? »

\subsection*{Pourquoi un seul irréductible}

Les pages 184 et 185 règlent le préalable. Soit $P(x)$ une application
polynomiale de $G = k$ à valeurs dans un groupe linéaire, vérifiant
\[
P(x+y) = P(x)\,P(y) .
\]
Pour une représentation \emph{irréductible}, le commutant de l'image est un
corps $D$ par le lemme de Schur, les $P(x)$ prennent leurs valeurs dans son
centre $Z$, et l'on est ramené à une application polynomiale de $k$ dans les
inversibles d'une extension commutative. L'argument est alors immédiat : si $P$
était un polynôme de degré $n > 0$, alors $P(x+y)$ serait de degré $n$ et
$P(x)P(y)$ de degré $2n$ --- absurde. Donc $P(x) = 1$, $Z = k$, et
\begin{quote}
\textbf{Proposition.} La seule représentation rationnelle \emph{irréductible}
du groupe additif est la représentation triviale de dimension $1$.
\end{quote}
\footnote{La page 185 écrit « indécomposables » et non « irréductibles ». Prise
au mot la proposition est fausse : $k[X]/X^{2}$ est indécomposable de dimension
$2$, et la page 183 vient précisément de dire que les indécomposables sont tous
les $k[X]/X^{i}$. C'est le lemme de Schur --- le commutant est un \emph{corps}
--- qui est invoqué dans la démonstration, et il ne vaut que pour les
irréductibles ; l'anneau des endomorphismes d'un indécomposable de longueur
finie est local et non un corps. Le Corollaire de la même page, qui triangule
toute représentation avec des $1$ sur la diagonale, est exactement l'énoncé
corrigé.}

\textbf{Corollaire.} Toute représentation rationnelle de $k$ se met sous forme
triangulaire avec des $1$ sur la diagonale.

Et il met aussitôt en regard la caractéristique $p$, où la situation est « tout
à fait distincte ». Les applications
\[
f(\lambda,\mu)(x) = \lambda x^{p^{h}} + \mu
\qquad (h \text{ fixé})
\]
donnent, par translation,
\[
\tau_{y}f(\lambda,\mu)(x)
= \lambda(x+y)^{p^{h}} + \mu
= \lambda x^{p^{h}} + \mu + \lambda y^{p^{h}}
= f\bigl(\lambda,\ \mu + \lambda y^{p^{h}}\bigr) ,
\]
c'est-à-dire la représentation de dimension $2$
\[
E_{h}(y) = \begin{pmatrix} 1 & y^{p^{h}} \\ 0 & 1 \end{pmatrix} ,
\]
et ces représentations sont deux à deux inéquivalentes pour $h$ variable : il y
a donc une \emph{infinité} de classes de représentations indécomposables de
dimension $2$, là où la caractéristique nulle n'en a qu'une.\footnote{C'est,
dans le vocabulaire d'aujourd'hui, le fait que le groupe additif est de type de
représentation \emph{sauvage} en caractéristique positive : les homomorphismes
$G_{a} \to G_{a}$ sont les polynômes additifs, qui forment un anneau de rang
infini, de sorte que $\operatorname{Ext}^{1}$ du trivial par le trivial est de
dimension infinie. Les $E_{h}$ en sont les classes les plus visibles.}

\subsection*{Caractéristique nulle : la Proposition 2}

La page 186 énonce, pour $k$ de caractéristique nulle, ce qui suit : les
représentations rationnelles du groupe additif correspondent aux
représentations linéaires de son algèbre de Lie, qui est \emph{de dimension
$1$} et abélienne ; et par la Proposition qui précède, ce sont exactement les
représentations \emph{nilpotentes} de cette algèbre. D'où :
\begin{quote}
\textbf{Proposition 2.} En caractéristique nulle, la classification des
représentations linéaires du groupe additif est équivalente à celle des
$k[X]$-modules de type fini $(X)$-primaires.
\end{quote}
\footnote{La page écrit « des groupes algébriques \emph{abéliens} », et ainsi
énoncé le résultat est faux : les représentations d'un tore sont semi-simples
et classées par les caractères, ce qui n'a rien à voir avec des modules
$(X)$-primaires. L'argument de la page passe par une algèbre de Lie abélienne
\emph{de dimension $1$}, donc par un groupe unipotent de dimension $1$,
c'est-à-dire par le groupe additif. Pour un groupe unipotent abélien de
dimension $r \geqslant 2$ le problème est celui de $r$ opérateurs nilpotents
qui commutent, et il n'est pas de cette forme.}

Les indécomposables $(X)$-primaires sont les $k[X]/X^{n}$ pour $1 \leqslant n
< \infty$, et le Corollaire nomme l'objet cherché : $R(k)$, l'anneau des
classes de représentations du groupe additif.

\subsection*{Le calcul : les puissances de $\xi$}

Les pages 187 et 188 sont les deux seules du faisceau qui se lisent presque de
bout en bout, et elles font le calcul. Notons $E_{n}$ la classe de
$k[X]/X^{n}$ ; l'involution est l'identité, $d(E_{n}) = n$, et $E_{1}$ est
l'unité. Posons
\[
\xi = E_{2} , \qquad d(\xi) = 2 .
\]
Alors $\xi^{n}$ est la classe de
\[
F_{n} = k[X_{1}, \ldots, X_{n}]\,/\,(X_{1}^{2}, \ldots, X_{n}^{2})
\]
muni de l'opérateur nilpotent $T$ de multiplication par la classe $\dot{S}$ de
$S = X_{1} + \cdots + X_{n}$. Le point du calcul est l'expression des
itérés : comme les carrés sont nuls, le développement multinomial de $S^{k}$ ne
laisse que les monômes à indices distincts, chacun avec le coefficient $k!$,
\[
\dot{S}^{k} = k!\,\sum_{1 \leqslant i_{1} < \cdots < i_{k} \leqslant n}
\dot{X}_{i_{1}} \cdots \dot{X}_{i_{k}} .
\]
En particulier $T^{n}$ est la multiplication par
$n!\,\dot{X}_{1}\cdots\dot{X}_{n}$ et $T^{n+1} = 0$ ; donc $T^{n}F_{n}$ est de
dimension $1$, de base $\dot{X}_{1}\cdots\dot{X}_{n}$, et
\emph{$F_{n}$ a un unique bloc de Jordan de taille $n+1$}. D'où :
$\xi^{n}$ est une combinaison linéaire à coefficients dans $\mathbf{N}$ des
$E_{1}, \ldots, E_{n+1}$, celui de $E_{n+1}$ valant $1$ :
\[
\xi^{n} = \sum_{i=1}^{n+1} c_{ni}\,E_{i} ,
\qquad c_{n,n+1} = 1 .
\]
\footnote{La page 188 écrit la somme jusqu'à $n$ avec $c_{nn} = 1$. Sa propre
table, quatre lignes plus bas, donne $\xi^{4} = 2E_{1} + 3E_{3} + E_{5}$ : le
terme de tête est $E_{n+1}$, et la page 195 écrit bien $\sum_{i=1}^{n+1}$ avec
$c_{n,n+1} = 1$ pour le calcul analogue. L'indice supérieur de la page 188 est
donné dans la transcription comme un $n$ surmonté d'un trait, sans que le
feuillet le décide ; cette lecture prend $n+1$, que le calcul exige.}

\textbf{Proposition 3.} $R(k)$ est l'anneau de polynômes $\mathbf{Z}[\xi]$ sur
$\mathbf{Z}$ engendré par $\xi = E_{2}$, de degré $2$.

Le calcul explicite, qu'il donne des deux côtés :
\begin{align*}
\xi^{0} &= E_{1}                      &  E_{1} &= 1 \\
\xi^{1} &= E_{2}                      &  E_{2} &= \xi \\
\xi^{2} &= E_{1} + E_{3}              &  E_{3} &= \xi^{2} - 1 \\
\xi^{3} &= 2E_{2} + E_{4}             &  E_{4} &= \xi^{3} - 2\xi \\
\xi^{4} &= 2E_{1} + 3E_{3} + E_{5}    &  E_{5} &= \xi^{4} - 3\xi^{2} + 1
\end{align*}
\footnote{La dernière ligne de la colonne de droite est « $+\,5$ » sur la
page ; la colonne de gauche impose $+\,1$, puisque $E_{5} = \xi^{4} - 2E_{1} -
3E_{3} = \xi^{4} - 2 - 3(\xi^{2}-1)$. La transcription donne la leçon du
feuillet et la signale ; cette lecture rétablit le $+1$.}

La règle qui les engendre est la récurrence à trois termes
\[
\xi\,E_{n} = E_{n-1} + E_{n+1} ,
\qquad\text{soit}\qquad
E_{n+1} = \xi\,E_{n} - E_{n-1} ,
\]
qui est la décomposition du produit tensoriel de deux blocs de Jordan
nilpotents : $J_{m} \otimes J_{n} \simeq \bigoplus_{i=1}^{\min(m,n)}
J_{m+n+1-2i}$.\footnote{C'est la règle de Clebsch-Gordan, et ce n'est pas une
coïncidence : un nilpotent régulier est un nilpotent principal de
$\mathfrak{sl}_{2}$, et en caractéristique nulle la décomposition des produits
tensoriels de blocs de Jordan est celle des représentations de
$\mathfrak{sl}_{2}$. L'anneau $R(k) = \mathbf{Z}[\xi]$ est donc isomorphe à
l'anneau des représentations de $\mathrm{SL}_{2}$, les $E_{n}$ répondant aux
irréductibles de dimension $n$ et $\xi$ à la représentation standard. Le
dossier calcule l'anneau et ne fait pas ce rapprochement.}

Une parenthèse ferme le feuillet, dont le premier mot --- une capitale $A$ ---
ne se lit pas : toutes les structures sont bien données, mais la structure
d'anneau « est un peu cachée ! ».

\subsection*{II. Le groupe multiplicatif}

La page 189 ouvre le second cas, et cette fois tout est semi-simple. Les
représentations irréductibles de $k^{*}$ passent par une représentation
rationnelle $P \colon k^{*} \to L^{*}$, $L$ une extension finie de $k$, avec
\[
P(\lambda) = \sum_{i=-n}^{n} \lambda^{i} c_{i} , \qquad c_{i} \in L .
\]
La relation $P(\lambda\lambda') = P(\lambda)P(\lambda')$ s'écrit
\[
\sum_{i} \lambda^{i}\lambda'^{i}c_{i}
= \sum_{i,j} \lambda^{i}\lambda'^{j}c_{i}c_{j} ,
\]
et comme c'est une identité en $\lambda$, $\lambda'$, on doit avoir
\[
c_{i} = c_{i}^{2} ,
\qquad c_{i}c_{j} = 0 \ \text{ pour } i \neq j .
\]
Il ne peut donc y avoir qu'un seul $c_{n}$ non nul, et $c_{n}^{2} = c_{n}$ avec
$c_{n}$ inversible donne $c_{n} = 1$. Donc $P(\lambda) = \lambda^{n}$, et
l'irréductibilité impose $L = k$. D'où :
\begin{quote}
\textbf{Proposition.} Les seules représentations irréductibles de $k^{*}$, $k$
un corps quelconque, sont celles de dimension $1$,
\[
\sigma_{n}(\lambda)\,t = \lambda^{n}\,t ,
\]
avec
\[
\sigma_{n}\sigma_{m} = \sigma_{n+m} ,
\qquad \check{\sigma}_{n} = \sigma_{-n} .
\]
\end{quote}
Pour une représentation quelconque, écrivant $P(\lambda) = \sum_{i=0}^{n}
c_{i}\lambda^{i}$ avec $c_{i} \in \operatorname{End}(V)$, le même calcul donne
$c_{i}c_{j} = 0$ et $c_{i}^{2} = c_{i}$, et de plus
\[
\sum_{i} c_{i} = P(1) = 1 :
\]
les $c_{i}$ sont des \emph{projecteurs deux à deux orthogonaux de somme $1$},
c'est-à-dire une décomposition de $V$. D'où la complète réductibilité :
\begin{quote}
\textbf{Proposition.} Toute représentation linéaire rationnelle de $k^{*}$ est
complètement réductible.
\end{quote}
\footnote{La page écrit une capitale $C$ là où le contexte, et le corollaire
qui suit immédiatement sur $R(k^{*})$, demandent $k^{*}$. La transcription
donne la leçon et la signale.}

\textbf{Corollaire.} $R(k^{*}) \simeq \mathbf{Z}[\mathbf{Z}]$, l'algèbre du
groupe des entiers, $\sigma_{n}$ répondant à $n$ ; l'involution et
l'augmentation y sont les usuelles. Et la généralisation à un tore : si
$\Gamma = X(G)$ est le groupe des caractères rationnels de $G$ dans $k^{*}$,
alors $\Gamma \simeq \mathbf{Z}^{n}$ et
\[
R(G) \simeq \mathbf{Z}[X(G)] .
\]

\subsection*{III. Le produit semi-direct}

La page 191 ouvre le troisième et dernier cas, qui occupe les sept derniers
feuillets :
\[
G = k \rtimes k^{*} ,
\qquad k^{*} \ \text{opérant dans } k \ \text{par } \ \sigma_{p}(\lambda)t =
\lambda^{p}t ,
\]
les deux facteurs étant notés $F = k$ et $H = k^{*}$ sur la page. En
caractéristique nulle, l'algèbre de Lie est l'extension de $F$ par $H$ avec
\[
(1)\qquad [\,Y, X\,] = pX ,
\]
$X$ et $Y$ étant les générateurs canoniques de $F$ et de $H$.

Soit $\varphi$ une représentation linéaire indécomposable rationnelle de $G$.
Restreinte à $H = k^{*}$ elle est semi-simple par le cas II, d'où la
décomposition en sous-espaces propres
\[
V = \sum_{n} V_{n} ,
\qquad Y \cdot a = n\,a \ \text{ pour } a \in V_{n} .
\]
La condition $(1)$ appliquée aux composantes $X_{ji} \in L(V_{i}, V_{j})$ de
$X$ donne
\[
(j - i)\,X_{ji} = p\,X_{ji} ,
\qquad\text{donc}\qquad
X_{ji} = 0 \ \text{ sauf si } \ j - i = p :
\]
\emph{$X$ décale le poids de exactement $p$}.

Deux cas. Si tous les $X_{ji}$ sont nuls, l'indécomposabilité laisse un seul
$V_{n}$ non nul, de dimension $1$, et la représentation est le $\sigma_{n}$ du
cas II remonté par la projection $G \to H$. Sinon, posant
\[
V' = \sum_{n \equiv \alpha \ (p)} V_{n} ,
\qquad
V'' = \sum_{n \not\equiv \alpha \ (p)} V_{n} ,
\]
les deux sont stables sous $G$ et $V = V' \oplus V''$ ; comme $V' \neq 0$,
l'indécomposabilité donne $V' = V$ : \emph{tous les poids non nuls sont
congrus modulo $p$}. Prenant $n$ minimal avec $V_{n} \neq 0$, les $V$ non nuls
sont les
\[
V_{n+pk} , \qquad k = 0, 1, \ldots
\]
Soit $k$ le plus grand indice avec $V_{n+pk} \neq 0$ --- et $k \geqslant 1$
dans ce cas. Posant
\[
u_{i} = X_{n+p(i+1),\,n+pi} , \qquad E_{i} = V_{n+pi} ,
\]
on a une suite d'applications
\[
E_{0} \xrightarrow{\ u_{0}\ } E_{1} \xrightarrow{\ u_{1}\ } E_{2}
\ \cdots\ \xrightarrow{\ u_{k-1}\ } E_{k} .
\]

\textbf{Lemme 1.} Les $u_{i}$, $0 \leqslant i \leqslant k-1$, sont injectifs.

\textbf{Lemme 2.} Ils sont bijectifs.

La démonstration du premier, page 193, est une récurrence sur $i$, triviale
pour $u_{-1} \colon 0 = E_{-1} \to E_{0}$, et l'argument de l'étape est celui-ci :
un $x$ non nul du noyau de $u_{\ell}$ engendre un $D \subset E_{\ell+1}$, on
choisit un supplémentaire $F$ de $D$, et l'on pose
\[
E'_{j} = \begin{cases}
E_{j} & j \leqslant \ell' \\
E_{j} \cap F & \ell' \leqslant j \leqslant \ell
\end{cases}
\qquad
E''_{j} = \begin{cases}
0 & j \leqslant \ell' \ \text{ou} \ j > \ell \\
D & \ell' \leqslant j \leqslant \ell
\end{cases}
\]
de sorte que $E_{j} = E'_{j} \oplus E''_{j}$ pour tout $j$, et que $\sum E'_{j}$
et $\sum E''_{j}$ soient tous deux stables et non nuls --- ce qui contredit
l'indécomposabilité.\footnote{Le haut de la page 193 est raturé de part en part
--- une quinzaine de traits et deux barres horizontales --- et ne rend que
quelques symboles. La construction des supplémentaires $E'_{j}$, $E''_{j}$ est
donnée dans la transcription à travers trois leçons successives dont deux sont
barrées ; cette lecture retient celle qui rend l'argument.}

\subsection*{Le catalogue des indécomposables}

\textbf{Conclusion (page 194).} Les $u_{i}$ sont bijectifs et les $E_{i}$ sont
de dimension $1$. La représentation ainsi obtenue est notée $e_{n,k+1}$, avec
$k+1 \geqslant 2$ ; et comme pour $k = 0$ on retrouve le cas précédent, il est
naturel de poser $e_{n,1} = \sigma_{n}$.

\begin{quote}
\textbf{Proposition.} Les représentations indécomposables de $G = k \rtimes
k^{*}$ sont, à équivalence près, les
\[
e_{n,k} , \qquad n \in \mathbf{Z} , \quad k \geqslant 1 ,
\]
deux à deux inéquivalentes, avec $e_{n,k} = \sigma_{n} \otimes e_{0,k}$ ;
$e_{0,k}$ est la représentation sur l'espace
\[
k[Z]/Z^{k}
\]
où $X$ est l'opérateur de multiplication par $Z$ modulo $Z^{k}$ et
\[
Y \cdot Z^{\alpha} = \alpha p\,Z^{\alpha} .
\]
$e_{0,k}$ est de dimension $k$, les poids de $e_{n,k}$ sont $n$, $n+p$,
\ldots, $n + p(k-1)$, et
\[
\check{e}_{n,k} = e_{-n-p(k-1),\,k} .
\]
\end{quote}
La formule du dual se lit sur les poids : le dual de $e_{n,k}$ a pour poids les
opposés, donc pour poids minimal $-n-p(k-1)$, et il est encore indécomposable
de dimension $k$.

\subsection*{Le calcul de l'anneau}

Les trois derniers feuillets calculent $R(G)$. Les homomorphismes
$F \xrightarrow{i} G \xrightarrow{j} H$ donnent
\[
R(H) \xrightarrow{\ j^{*}\ } R(G) \xrightarrow{\ i^{*}\ } R(F) .
\]
$R(H)$ est l'algèbre engendrée par les $\sigma_{n}$, $n \in \mathbf{Z}$, avec
$\sigma_{n}\sigma_{m} = \sigma_{n+m}$ et $\sigma_{0} = 1$ --- c'est
$\mathbf{Z}[\mathbf{Z}]$ du cas II --- et elle se plonge dans $R(G)$ par
\[
j^{*}(\sigma_{n}) = e_{n,1} .
\]
$R(F)$ est $\mathbf{Z}[e_{2}]$ du premier cas, de base les $e_{n}$ avec
$e_{1} = 1$, et $R(G)$ s'y envoie par
\[
i^{*}(e_{n,k}) = e_{k} .
\]
Identifiant $\Lambda = R(H)$ à une sous-algèbre de $R(G)$ et posant
$\xi = e_{0,2}$, l'énoncé à prouver est
\[
R(G) = \Lambda[\xi] ,
\]
c'est-à-dire que les $\xi^{n}$, $n = 0, 1, \ldots$, forment une \emph{base} de
$R(G)$ sur $\Lambda$. Comme les $e_{0,1}, e_{0,2}, \ldots$ en forment une, il
suffit d'établir les formules
\[
(1)\qquad \xi^{n} = \sum_{i=1}^{n+1} c_{ni}\,e_{0i} ,
\qquad c_{ni} \in \Lambda , \quad c_{n,n+1} = 1 ,
\]
qui sont triangulaires à diagonale $1$ et donc inversibles.

La démonstration, page 196, est un raisonnement en deux temps qu'il vaut la
peine de suivre. \emph{A priori} $\xi^{n} = \sum_{i} c_{ni}e_{0i}$ avec
$c_{ni} \in \Lambda^{+}$. Appliquant $i^{*}$ --- qui, sur $\Lambda$, est
l'augmentation $d$ --- on trouve
\[
e_{2}^{\,n} = \sum_{i} d(c_{ni})\,e_{i} ,
\]
et l'on sait par ailleurs, c'est le calcul des pages 187 et 188, que
$e_{2}^{\,n} = \sum_{i=1}^{n+1}a_{ni}e_{i}$ avec les $a_{ni}$ entiers
$\geqslant 0$ et $a_{n,n+1} = 1$. La comparaison donne
\[
\varepsilon(c_{ni}) = 0 \ \text{ pour } i > n+1 ,
\qquad \varepsilon(c_{n,n+1}) = 1 ,
\]
et comme les $c_{ni}$ sont dans $\Lambda^{+}$ --- combinaisons à coefficients
$\geqslant 0$ des $\sigma_{\alpha}$ --- il s'ensuit que
\[
c_{ni} = 0 \ \text{ pour } i > n+1 ,
\qquad c_{n,n+1} = \sigma_{\alpha} \ \text{ pour un certain } \alpha .
\]
Reste à voir que $\alpha = 0$, et c'est l'argument de poids de la page 197,
pris des deux côtés. D'un côté, le plus haut poids figurant dans $\xi^{n}$ est
$n$ fois le plus haut poids figurant dans $\xi$, ce qui donne
$p\alpha \leqslant pn$, donc $\alpha \leqslant 0$. De l'autre, les poids
intervenant dans $\xi = e_{0,2}$ sont $0$ et $p$, donc positifs, donc ceux de
$\xi^{n}$ aussi, et le poids apparaissant dans $\sigma_{\alpha}e_{0,p+1}$
impose $\alpha \geqslant 0$. Donc $\alpha = 0$, $\sigma_{\alpha} = 1$, c.q.f.d.

\begin{quote}
\textbf{Proposition (page 197).} $R(G)$ est l'algèbre de polynômes sur
$\Lambda = \mathbf{Z}[\Gamma]$, $\Gamma = \mathbf{Z}$, engendrée par
$\xi = e_{0,2}$ ; l'augmentation de $\Lambda$ est l'augmentation usuelle,
\[
d(\xi) = 2 ,
\]
l'involution $\vee$ est l'involution usuelle sur $\Lambda$, et
\[
\check{\xi} = e_{-p,2} = \sigma_{-p} \cdot \xi .
\]
\end{quote}
La dernière formule est le cas $n = 0$, $k = 2$ de celle de la page 194 :
$\check{e}_{0,2} = e_{-p,2}$, et $e_{-p,2} = \sigma_{-p} \otimes e_{0,2}$. Le
dossier se referme là-dessus.\footnote{La page écrit $d(\xi) = 2$ avec un $p$
barré devant le $2$ ; la dimension de $e_{0,2}$ est bien $2$, et le $p$ est une
fausse route corrigée sur le feuillet.}

\end{document}
